% !TEX program = lualatex
% ---------------------------------------------------------------------------------------------
%  Readback note 18: The grid of lifts and descents, and the crude rows
%
%  Build:  python3 tools/rb.py note <slug>
%
%  One unit of a Lean artifact, rewritten as ordinary mathematics. The prose is hand-written. Lean names,
%  line ranges and counts come from data.tex; quoted Lean statements are copied from the artifact by the build.
%
%  Rules of the text: standard terminology; every symbol is explained where it first appears and listed again
%  in the notation at the end; no letter has two meanings; a proof is complete or says what it takes for granted;
%  Lean names appear only in the LEAN lines and in the concordance.
%
%  The signature rule (the build enforces it): a function, relation or predicate that the artifact defines is
%  written only through its notation macro, which the build makes from the spec (`notation`) with one slot for
%  each explicit argument of the Lean definition. No argument is fixed and dropped, also not "for brevity" and
%  not for one section. Every argument and every hypothesis of a stated declaration stands inside its
%  statement, or in a Setting that the statement cites by number; only implicit and instance arguments (the
%  ambient objects) may stand in the conventions.
% ---------------------------------------------------------------------------------------------
\documentclass[11pt,a4paper]{article}
\ifdefined\pdfvariable\pdfvariable suppressoptionalinfo \numexpr 1+2+4+8\relax\fi
\usepackage[a4paper,hmargin=35mm,top=22mm,bottom=27mm,footskip=12mm,nohead]{geometry}

\usepackage{fontspec}
\directlua{
  luaotfload.add_fallback("gfallback", {"STIXTwoMath-Regular.otf:mode=node;", "DejaVuSans.ttf:mode=node;"})
  luaotfload.add_fallback("gmonofallback", {"DejaVuSansMono.ttf:mode=node;", "STIXTwoMath-Regular.otf:mode=node;"})
}
\setmainfont{NewCM10}[Extension=.otf,UprightFont=*-Book,ItalicFont=*-BookItalic,
  BoldFont=*-Bold,BoldItalicFont=*-BoldItalic,SmallCapsFont=*-Book,SmallCapsFeatures={Letters=SmallCaps},
  RawFeature={fallback=gfallback}]
\setsansfont{NewCMSans10}[Extension=.otf,UprightFont=*-Book,ItalicFont=*-BookOblique,
  BoldFont=*-Bold,BoldItalicFont=*-BoldOblique,RawFeature={fallback=gfallback}]
\setmonofont{JetBrainsMono}[Extension=.otf,UprightFont=*-Regular,ItalicFont=*-Italic,
  BoldFont=*-SemiBold,BoldItalicFont=*-SemiBoldItalic,Scale=MatchLowercase,
  RawFeature={-calt;-liga;fallback=gmonofallback}]
\usepackage{amsmath}
\usepackage{amsthm}
\usepackage[math-style=TeX]{unicode-math}
\setmathfont{NewCMMath-Book.otf}
\emergencystretch=2.5em\relax

\usepackage{xcolor}
\usepackage{tikz}
\usetikzlibrary{calc,arrows.meta,positioning,backgrounds}
\usepackage[protrusion=true,expansion=true]{microtype}
\usepackage{enumitem}
\usepackage{needspace}
\usepackage{booktabs}
\usepackage{array}
\usepackage{longtable}
\usepackage{fancyhdr}
\usepackage{lastpage}
\usepackage{multicol}
\usepackage[hidelinks,pdfusetitle]{hyperref}

\definecolor{ink}{HTML}{0F172A}
\definecolor{inksoft}{HTML}{475569}
\definecolor{inkfaint}{HTML}{94A3B8}
\definecolor{rulegrey}{HTML}{CBD5E1}
\definecolor{arith}{HTML}{C2410C}
\definecolor{est}{HTML}{7C3AED}
\definecolor{ana}{HTML}{0891B2}
\definecolor{ok}{HTML}{047857}
\definecolor{leancol}{HTML}{3730A3}

\title{The grid of lifts and descents, and the crude rows}
\author{grh-compact}
\date{}

% ------------------------------------------------------------------------------------- macros
\newcommand{\leanfont}{\ttfamily\fontsize{8.6}{10.5}\selectfont\color{leancol}}
\newcommand{\lean}[1]{\csname LN@#1\endcsname}
\newcommand{\fillin}[1]{{\color{arith}\itshape[#1]}}  % a place of the template that is not written yet; the build reports every one that is left
\newcommand{\leanquote}[1]{\csname LQ@#1\endcsname}
% generated by rb.py note; do not edit
\newcommand{\ConcDecls}{125}
\newcommand{\Pin}{72fead10}
\newcommand{\PinText}{that of commit {\leanfont 72fead10}}
\newcommand{\ArtifactName}{grh-compact}
\expandafter\newcommand\csname LN@existsisTidyintegraleqgridPeriod\endcsname{\href{https://github.com/lsernau/grh-lean/blob/72fead102f663d5c7a407964a13dc0cf9d977207/grh-compact/src/New/Exchange/CrudeRows.lean\#L23-L40}{{\leanfont Data.\allowbreak{}exists\_isTidy\_integral\_eq\_gridPeriod}}}
\expandafter\newcommand\csname LN@HasRoute\endcsname{\href{https://github.com/lsernau/grh-lean/blob/72fead102f663d5c7a407964a13dc0cf9d977207/grh-compact/src/New/Exchange/CrudeRows.lean\#L44-L47}{{\leanfont Data.\allowbreak{}HasRoute}}}
\expandafter\newcommand\csname LQ@HasRoute\endcsname{\leanquoteline{0}{def Data.HasRoute \{i j\} (h : i ≤ D.m) (h' : j ≤ D.m) k := ∃ d : \_ → ℂˣ, D.CostsSquares k (d ·) ∧ ∀ n}\leanquoteline{2}{Λ, (∀ ϑ, IsTidy n (D.path i) (Λ ϑ)) → ∃ Φ, (∀ ϑ, IsTidy (n + k) (D.path j) (Φ ϑ)) ∧ ∀ T : \_ → ℂ,}\leanquoteline{2}{(∀ g, ∃ β, ∀ x, gridPeriod Φ x h' g = β * T x) → ∀ g, ∃ β, ∀ x, gridPeriod Λ x h g =}\leanquoteline{2}{β * (T x / d x.1)}}
\expandafter\newcommand\csname LN@hasRoutesucc\endcsname{\href{https://github.com/lsernau/grh-lean/blob/72fead102f663d5c7a407964a13dc0cf9d977207/grh-compact/src/New/Exchange/CrudeRows.lean\#L52-L104}{{\leanfont Data.\allowbreak{}hasRoute\_succ}}}
\expandafter\newcommand\csname LN@gridPeriodclimb\endcsname{\href{https://github.com/lsernau/grh-lean/blob/72fead102f663d5c7a407964a13dc0cf9d977207/grh-compact/src/New/Exchange/CrudeRows.lean\#L109-L116}{{\leanfont Data.\allowbreak{}gridPeriod\_climb}}}
\expandafter\newcommand\csname LN@existsmulendPeriodmemembValRing\endcsname{\href{https://github.com/lsernau/grh-lean/blob/72fead102f663d5c7a407964a13dc0cf9d977207/grh-compact/src/New/Exchange/CrudeRows.lean\#L120-L151}{{\leanfont Data.\allowbreak{}exists\_mul\_endPeriod\_mem\_embValRing}}}
\expandafter\newcommand\csname LN@endPeriodle\endcsname{\href{https://github.com/lsernau/grh-lean/blob/72fead102f663d5c7a407964a13dc0cf9d977207/grh-compact/src/New/Exchange/CrudeRows.lean\#L156-L184}{{\leanfont Data.\allowbreak{}endPeriod\_le}}}
\expandafter\newcommand\csname LN@cruderows\endcsname{\href{https://github.com/lsernau/grh-lean/blob/72fead102f663d5c7a407964a13dc0cf9d977207/grh-compact/src/New/Exchange/CrudeRows.lean\#L191-L247}{{\leanfont Data.\allowbreak{}crude\_rows}}}
\expandafter\newcommand\csname LQ@cruderows\endcsname{\leanquoteline{0}{lemma Data.crude\_rows : ∃ c : \_ → ℂˣ, D.CostsSquares (2 * D.m * D.m) (c ·) ∧ span ℂ \{y |}\leanquoteline{4}{(∀ ϑ l [Fact l.Prime] ι, l.Coprime ϑ.levelNorm →}\leanquoteline{2}{(c ϑ / c 1 : ℂ) • D.iotaAtLevel ϑ y ∈ LevelPlaces.LevelForm.integralForms (l := l) ι) ∧}\leanquoteline{2}{∀ (σ : ℂ ≃+* ℂ) g, weylNormSq (σ ∘ (D.ι .one y).1 g) ≤ exp (D.x₁ \textasciicircum{} 2 * D.b \textasciicircum{} 2)\} = ⊤}}
\expandafter\newcommand\csname LN@triangleOrdereqmullevelIndex\endcsname{\href{https://github.com/lsernau/grh-lean/blob/72fead102f663d5c7a407964a13dc0cf9d977207/grh-compact/src/New/Exchange/GridFirstColumns.lean\#L22-L30}{{\leanfont Data.\allowbreak{}triangleOrder\_eq\_mul\_levelIndex}}}
\expandafter\newcommand\csname LN@diagonalPeriodFormformRep\endcsname{\href{https://github.com/lsernau/grh-lean/blob/72fead102f663d5c7a407964a13dc0cf9d977207/grh-compact/src/New/Exchange/GridFirstColumns.lean\#L35-L42}{{\leanfont Data.\allowbreak{}diagonalPeriodForm\_formRep}}}
\expandafter\newcommand\csname LN@existsrestrictbordered\endcsname{\href{https://github.com/lsernau/grh-lean/blob/72fead102f663d5c7a407964a13dc0cf9d977207/grh-compact/src/New/Exchange/GridFirstColumns.lean\#L47-L70}{{\leanfont Data.\allowbreak{}exists\_restrict\_bordered}}}
\expandafter\newcommand\csname LN@formReplegDescentstupleAction\endcsname{\href{https://github.com/lsernau/grh-lean/blob/72fead102f663d5c7a407964a13dc0cf9d977207/grh-compact/src/New/Exchange/GridFirstColumns.lean\#L79-L92}{{\leanfont Data.\allowbreak{}formRep\_legDescents\_tupleAction}}}
\expandafter\newcommand\csname LN@diagonalPeriodFormtupleAction\endcsname{\href{https://github.com/lsernau/grh-lean/blob/72fead102f663d5c7a407964a13dc0cf9d977207/grh-compact/src/New/Exchange/GridFirstColumns.lean\#L97-L106}{{\leanfont Data.\allowbreak{}diagonalPeriodForm\_tupleAction}}}
\expandafter\newcommand\csname LN@gridPeriodeqendPeriod\endcsname{\href{https://github.com/lsernau/grh-lean/blob/72fead102f663d5c7a407964a13dc0cf9d977207/grh-compact/src/New/Exchange/GridFirstColumns.lean\#L110-L134}{{\leanfont Data.\allowbreak{}gridPeriod\_eq\_endPeriod}}}
\expandafter\newcommand\csname LN@endPeriodmem\endcsname{\href{https://github.com/lsernau/grh-lean/blob/72fead102f663d5c7a407964a13dc0cf9d977207/grh-compact/src/New/Exchange/GridFirstColumns.lean\#L138-L193}{{\leanfont Data.\allowbreak{}endPeriod\_mem}}}
\expandafter\newcommand\csname LN@existsisLevelTidyiotaPluseq\endcsname{\href{https://github.com/lsernau/grh-lean/blob/72fead102f663d5c7a407964a13dc0cf9d977207/grh-compact/src/New/Exchange/GridFirstColumns.lean\#L199-L237}{{\leanfont Data.\allowbreak{}exists\_isLevelTidy\_iotaPlus\_eq}}}
\expandafter\newcommand\csname LN@existsisTidydiagonalPeriodFormeq\endcsname{\href{https://github.com/lsernau/grh-lean/blob/72fead102f663d5c7a407964a13dc0cf9d977207/grh-compact/src/New/Exchange/GridFirstColumns.lean\#L242-L272}{{\leanfont Data.\allowbreak{}exists\_isTidy\_diagonalPeriodForm\_eq}}}
\expandafter\newcommand\csname LN@fitsParamsauxPaira\endcsname{\href{https://github.com/lsernau/grh-lean/blob/72fead102f663d5c7a407964a13dc0cf9d977207/grh-compact/src/New/Exchange/GridModuli.lean\#L14-L18}{{\leanfont LValueData.\allowbreak{}fitsParams\_auxPair}}}
\expandafter\newcommand\csname LN@heckePolyOnFormstestValueeqzero\endcsname{\href{https://github.com/lsernau/grh-lean/blob/72fead102f663d5c7a407964a13dc0cf9d977207/grh-compact/src/New/Exchange/GridModuli.lean\#L24-L28}{{\leanfont AuxPair.\allowbreak{}heckePolyOnForms\_testValue\_eq\_zero}}}
\expandafter\newcommand\csname LN@nStar\endcsname{\href{https://github.com/lsernau/grh-lean/blob/72fead102f663d5c7a407964a13dc0cf9d977207/grh-compact/src/New/Exchange/GridModuli.lean\#L33-L33}{{\leanfont Data.\allowbreak{}nStar}}}
\expandafter\newcommand\csname LN@nStarbounds\endcsname{\href{https://github.com/lsernau/grh-lean/blob/72fead102f663d5c7a407964a13dc0cf9d977207/grh-compact/src/New/Exchange/GridModuli.lean\#L35-L36}{{\leanfont Data.\allowbreak{}nStar\_bounds}}}
\expandafter\newcommand\csname LN@orderW\endcsname{\href{https://github.com/lsernau/grh-lean/blob/72fead102f663d5c7a407964a13dc0cf9d977207/grh-compact/src/New/Exchange/GridModuli.lean\#L42-L43}{{\leanfont Data.\allowbreak{}orderW}}}
\expandafter\newcommand\csname LN@conjHolWeight\endcsname{\href{https://github.com/lsernau/grh-lean/blob/72fead102f663d5c7a407964a13dc0cf9d977207/grh-compact/src/New/Exchange/GridModuli.lean\#L47-L47}{{\leanfont Data.\allowbreak{}conjHolWeight}}}
\expandafter\newcommand\csname LN@kernelWeight\endcsname{\href{https://github.com/lsernau/grh-lean/blob/72fead102f663d5c7a407964a13dc0cf9d977207/grh-compact/src/New/Exchange/GridModuli.lean\#L51-L52}{{\leanfont Data.\allowbreak{}kernelWeight}}}
\expandafter\newcommand\csname LN@rankboundscardbigIndex\endcsname{\href{https://github.com/lsernau/grh-lean/blob/72fead102f663d5c7a407964a13dc0cf9d977207/grh-compact/src/New/Exchange/GridModuli.lean\#L59-L61}{{\leanfont Data.\allowbreak{}rank\_bounds\_card\_bigIndex}}}
\expandafter\newcommand\csname LN@twomulcarddetIndexaddle\endcsname{\href{https://github.com/lsernau/grh-lean/blob/72fead102f663d5c7a407964a13dc0cf9d977207/grh-compact/src/New/Exchange/GridModuli.lean\#L63-L65}{{\leanfont Data.\allowbreak{}two\_mul\_card\_detIndex\_add\_le}}}
\expandafter\newcommand\csname LN@conjHolWeighteqaddkernelWeight\endcsname{\href{https://github.com/lsernau/grh-lean/blob/72fead102f663d5c7a407964a13dc0cf9d977207/grh-compact/src/New/Exchange/GridModuli.lean\#L70-L76}{{\leanfont Data.\allowbreak{}conjHolWeight\_eq\_add\_kernelWeight}}}
\expandafter\newcommand\csname LN@conjHolWeightkernelWeightbounds\endcsname{\href{https://github.com/lsernau/grh-lean/blob/72fead102f663d5c7a407964a13dc0cf9d977207/grh-compact/src/New/Exchange/GridModuli.lean\#L78-L87}{{\leanfont Data.\allowbreak{}conjHolWeight\_kernelWeight\_bounds}}}
\expandafter\newcommand\csname LN@isPosFirstOrderorderW\endcsname{\href{https://github.com/lsernau/grh-lean/blob/72fead102f663d5c7a407964a13dc0cf9d977207/grh-compact/src/New/Exchange/GridModuli.lean\#L91-L92}{{\leanfont Data.\allowbreak{}isPosFirstOrder\_orderW}}}
\expandafter\newcommand\csname LN@paramWltoforderWlt\endcsname{\href{https://github.com/lsernau/grh-lean/blob/72fead102f663d5c7a407964a13dc0cf9d977207/grh-compact/src/New/Exchange/GridModuli.lean\#L94-L95}{{\leanfont Data.\allowbreak{}paramW\_lt\_of\_orderW\_lt}}}
\expandafter\newcommand\csname LN@instFactNeRingOfIntegersEHPowGbOfNat\endcsname{\href{https://github.com/lsernau/grh-lean/blob/72fead102f663d5c7a407964a13dc0cf9d977207/grh-compact/src/New/Exchange/GridModuli.lean\#L97-L97}{{\leanfont Data.\allowbreak{}instFactNeRingOfIntegersEHPowGbOfNat}}}
\expandafter\newcommand\csname LN@placesmodulusequnion\endcsname{\href{https://github.com/lsernau/grh-lean/blob/72fead102f663d5c7a407964a13dc0cf9d977207/grh-compact/src/New/Exchange/GridModuli.lean\#L102-L103}{{\leanfont Data.\allowbreak{}places\_modulus\_eq\_union}}}
\expandafter\newcommand\csname LN@isTameForofhasWeight\endcsname{\href{https://github.com/lsernau/grh-lean/blob/72fead102f663d5c7a407964a13dc0cf9d977207/grh-compact/src/New/Exchange/GridModuli.lean\#L105-L108}{{\leanfont Data.\allowbreak{}isTameFor\_of\_hasWeight}}}
\expandafter\newcommand\csname LN@isTameFortwistedChi\endcsname{\href{https://github.com/lsernau/grh-lean/blob/72fead102f663d5c7a407964a13dc0cf9d977207/grh-compact/src/New/Exchange/GridModuli.lean\#L113-L115}{{\leanfont Data.\allowbreak{}isTameFor\_twistedChi}}}
\expandafter\newcommand\csname LN@instNeZeroNatB\endcsname{\href{https://github.com/lsernau/grh-lean/blob/72fead102f663d5c7a407964a13dc0cf9d977207/grh-compact/src/New/Exchange/GridModuli.lean\#L117-L117}{{\leanfont Data.\allowbreak{}instNeZeroNatB₁}}}
\expandafter\newcommand\csname LN@gridLevelGen\endcsname{\href{https://github.com/lsernau/grh-lean/blob/72fead102f663d5c7a407964a13dc0cf9d977207/grh-compact/src/New/Exchange/GridModuli.lean\#L120-L120}{{\leanfont Data.\allowbreak{}gridLevelGen}}}
\expandafter\newcommand\csname LN@gridModulus\endcsname{\href{https://github.com/lsernau/grh-lean/blob/72fead102f663d5c7a407964a13dc0cf9d977207/grh-compact/src/New/Exchange/GridModuli.lean\#L124-L124}{{\leanfont Data.\allowbreak{}gridModulus}}}
\expandafter\newcommand\csname LN@gridModulusoneideal\endcsname{\href{https://github.com/lsernau/grh-lean/blob/72fead102f663d5c7a407964a13dc0cf9d977207/grh-compact/src/New/Exchange/GridModuli.lean\#L128-L129}{{\leanfont Data.\allowbreak{}gridModulus\_one\_ideal}}}
\expandafter\newcommand\csname LN@triangleModulus\endcsname{\href{https://github.com/lsernau/grh-lean/blob/72fead102f663d5c7a407964a13dc0cf9d977207/grh-compact/src/New/Exchange/GridModuli.lean\#L133-L133}{{\leanfont Data.\allowbreak{}triangleModulus}}}
\expandafter\newcommand\csname LN@reduceToTriangle\endcsname{\href{https://github.com/lsernau/grh-lean/blob/72fead102f663d5c7a407964a13dc0cf9d977207/grh-compact/src/New/Exchange/GridModuli.lean\#L137-L140}{{\leanfont Data.\allowbreak{}reduceToTriangle}}}
\expandafter\newcommand\csname LN@piGBardvdgridLevelGen\endcsname{\href{https://github.com/lsernau/grh-lean/blob/72fead102f663d5c7a407964a13dc0cf9d977207/grh-compact/src/New/Exchange/GridModuli.lean\#L142-L142}{{\leanfont Data.\allowbreak{}piGBar\_dvd\_gridLevelGen}}}
\expandafter\newcommand\csname LN@EdgeTuple\endcsname{\href{https://github.com/lsernau/grh-lean/blob/72fead102f663d5c7a407964a13dc0cf9d977207/grh-compact/src/New/Exchange/GridMove.lean\#L21-L22}{{\leanfont Data.\allowbreak{}EdgeTuple}}}
\expandafter\newcommand\csname LN@GridDatum\endcsname{\href{https://github.com/lsernau/grh-lean/blob/72fead102f663d5c7a407964a13dc0cf9d977207/grh-compact/src/New/Exchange/GridMove.lean\#L26-L26}{{\leanfont Data.\allowbreak{}GridDatum}}}
\expandafter\newcommand\csname LN@IsMuEdge\endcsname{\href{https://github.com/lsernau/grh-lean/blob/72fead102f663d5c7a407964a13dc0cf9d977207/grh-compact/src/New/Exchange/GridMove.lean\#L30-L30}{{\leanfont Data.\allowbreak{}IsMuEdge}}}
\expandafter\newcommand\csname LN@edgeChar\endcsname{\href{https://github.com/lsernau/grh-lean/blob/72fead102f663d5c7a407964a13dc0cf9d977207/grh-compact/src/New/Exchange/GridMove.lean\#L34-L35}{{\leanfont Data.\allowbreak{}edgeChar}}}
\expandafter\newcommand\csname LN@IsAdmissibleFamily\endcsname{\href{https://github.com/lsernau/grh-lean/blob/72fead102f663d5c7a407964a13dc0cf9d977207/grh-compact/src/New/Exchange/GridMove.lean\#L39-L41}{{\leanfont Data.\allowbreak{}IsAdmissibleFamily}}}
\expandafter\newcommand\csname LN@tupleAction\endcsname{\href{https://github.com/lsernau/grh-lean/blob/72fead102f663d5c7a407964a13dc0cf9d977207/grh-compact/src/New/Exchange/GridMove.lean\#L44-L44}{{\leanfont Data.\allowbreak{}tupleAction}}}
\expandafter\newcommand\csname LN@edgeCharProd\endcsname{\href{https://github.com/lsernau/grh-lean/blob/72fead102f663d5c7a407964a13dc0cf9d977207/grh-compact/src/New/Exchange/GridMove.lean\#L48-L49}{{\leanfont Data.\allowbreak{}edgeCharProd}}}
\expandafter\newcommand\csname LN@tupleBound\endcsname{\href{https://github.com/lsernau/grh-lean/blob/72fead102f663d5c7a407964a13dc0cf9d977207/grh-compact/src/New/Exchange/GridMove.lean\#L53-L53}{{\leanfont Data.\allowbreak{}tupleBound}}}
\expandafter\newcommand\csname LN@LeftInvOnMuEdges\endcsname{\href{https://github.com/lsernau/grh-lean/blob/72fead102f663d5c7a407964a13dc0cf9d977207/grh-compact/src/New/Exchange/GridMove.lean\#L57-L58}{{\leanfont Data.\allowbreak{}LeftInvOnMuEdges}}}
\expandafter\newcommand\csname LN@IsShortSum\endcsname{\href{https://github.com/lsernau/grh-lean/blob/72fead102f663d5c7a407964a13dc0cf9d977207/grh-compact/src/New/Exchange/GridMove.lean\#L62-L64}{{\leanfont Data.\allowbreak{}IsShortSum}}}
\expandafter\newcommand\csname LN@bind\endcsname{\href{https://github.com/lsernau/grh-lean/blob/72fead102f663d5c7a407964a13dc0cf9d977207/grh-compact/src/New/Exchange/GridMove.lean\#L68-L69}{{\leanfont Data.\allowbreak{}GridDatum.\allowbreak{}bind}}}
\expandafter\newcommand\csname LN@IsTidy\endcsname{\href{https://github.com/lsernau/grh-lean/blob/72fead102f663d5c7a407964a13dc0cf9d977207/grh-compact/src/New/Exchange/GridMove.lean\#L73-L76}{{\leanfont Data.\allowbreak{}IsTidy}}}
\expandafter\newcommand\csname LN@IsLevelTidy\endcsname{\href{https://github.com/lsernau/grh-lean/blob/72fead102f663d5c7a407964a13dc0cf9d977207/grh-compact/src/New/Exchange/GridMove.lean\#L80-L81}{{\leanfont Data.\allowbreak{}IsLevelTidy}}}
\expandafter\newcommand\csname LN@verticalEdge\endcsname{\href{https://github.com/lsernau/grh-lean/blob/72fead102f663d5c7a407964a13dc0cf9d977207/grh-compact/src/New/Exchange/GridMove.lean\#L84-L84}{{\leanfont Data.\allowbreak{}verticalEdge}}}
\expandafter\newcommand\csname LN@rowEdge\endcsname{\href{https://github.com/lsernau/grh-lean/blob/72fead102f663d5c7a407964a13dc0cf9d977207/grh-compact/src/New/Exchange/GridMove.lean\#L88-L88}{{\leanfont Data.\allowbreak{}rowEdge}}}
\expandafter\newcommand\csname LN@legEdges\endcsname{\href{https://github.com/lsernau/grh-lean/blob/72fead102f663d5c7a407964a13dc0cf9d977207/grh-compact/src/New/Exchange/GridMove.lean\#L91-L91}{{\leanfont Data.\allowbreak{}legEdges}}}
\expandafter\newcommand\csname LN@gridPath\endcsname{\href{https://github.com/lsernau/grh-lean/blob/72fead102f663d5c7a407964a13dc0cf9d977207/grh-compact/src/New/Exchange/GridMove.lean\#L95-L95}{{\leanfont Data.\allowbreak{}gridPath}}}
\expandafter\newcommand\csname LN@path\endcsname{\href{https://github.com/lsernau/grh-lean/blob/72fead102f663d5c7a407964a13dc0cf9d977207/grh-compact/src/New/Exchange/GridMove.lean\#L99-L99}{{\leanfont Data.\allowbreak{}path}}}
\expandafter\newcommand\csname LN@legDescents\endcsname{\href{https://github.com/lsernau/grh-lean/blob/72fead102f663d5c7a407964a13dc0cf9d977207/grh-compact/src/New/Exchange/GridMove.lean\#L104-L106}{{\leanfont Data.\allowbreak{}legDescents}}}
\expandafter\newcommand\csname LN@gridPathMap\endcsname{\href{https://github.com/lsernau/grh-lean/blob/72fead102f663d5c7a407964a13dc0cf9d977207/grh-compact/src/New/Exchange/GridMove.lean\#L113-L114}{{\leanfont Data.\allowbreak{}gridPathMap}}}
\expandafter\newcommand\csname LN@diagonalPeriodForm\endcsname{\href{https://github.com/lsernau/grh-lean/blob/72fead102f663d5c7a407964a13dc0cf9d977207/grh-compact/src/New/Exchange/GridMove.lean\#L118-L121}{{\leanfont Data.\allowbreak{}diagonalPeriodForm}}}
\expandafter\newcommand\csname LN@gridPeriod\endcsname{\href{https://github.com/lsernau/grh-lean/blob/72fead102f663d5c7a407964a13dc0cf9d977207/grh-compact/src/New/Exchange/GridMove.lean\#L125-L127}{{\leanfont Data.\allowbreak{}gridPeriod}}}
\expandafter\newcommand\csname LQ@gridPeriod\endcsname{\leanquoteline{0}{@[simp] def gridPeriod (Λ : ∀ ϑ, D.GridDatum ϑ) (y : \_ × \_ × \_) h g := Λ y.1 fun S ↦}\leanquoteline{2}{diagonalPeriodForm y.1 h g}\leanquoteline{2}{(gridPathMap S \_ j j (levelIndex y.1 false 0 D.m • y.2.1)) (gridPathMap S \_ j j y.2.2)}}
\expandafter\newcommand\csname LN@edgeCharProdmul\endcsname{\href{https://github.com/lsernau/grh-lean/blob/72fead102f663d5c7a407964a13dc0cf9d977207/grh-compact/src/New/Exchange/GridMove.lean\#L129-L130}{{\leanfont Data.\allowbreak{}edgeCharProd\_mul}}}
\expandafter\newcommand\csname LN@legDescentscongr\endcsname{\href{https://github.com/lsernau/grh-lean/blob/72fead102f663d5c7a407964a13dc0cf9d977207/grh-compact/src/New/Exchange/GridMove.lean\#L134-L135}{{\leanfont Data.\allowbreak{}legDescents\_congr}}}
\expandafter\newcommand\csname LN@gridPathMapcongr\endcsname{\href{https://github.com/lsernau/grh-lean/blob/72fead102f663d5c7a407964a13dc0cf9d977207/grh-compact/src/New/Exchange/GridMove.lean\#L137-L138}{{\leanfont Data.\allowbreak{}gridPathMap\_congr}}}
\expandafter\newcommand\csname LN@binda\endcsname{\href{https://github.com/lsernau/grh-lean/blob/72fead102f663d5c7a407964a13dc0cf9d977207/grh-compact/src/New/Exchange/GridMove.lean\#L143-L147}{{\leanfont Data.\allowbreak{}IsShortSum.\allowbreak{}bind}}}
\expandafter\newcommand\csname LN@cardmatrixletupleBound\endcsname{\href{https://github.com/lsernau/grh-lean/blob/72fead102f663d5c7a407964a13dc0cf9d977207/grh-compact/src/New/Exchange/GridMove.lean\#L149-L152}{{\leanfont Data.\allowbreak{}card\_matrix\_le\_tupleBound}}}
\expandafter\newcommand\csname LN@EndIndex\endcsname{\href{https://github.com/lsernau/grh-lean/blob/72fead102f663d5c7a407964a13dc0cf9d977207/grh-compact/src/New/Exchange/GridMove.lean\#L156-L156}{{\leanfont Data.\allowbreak{}EndIndex}}}
\expandafter\newcommand\csname LN@endWeightData\endcsname{\href{https://github.com/lsernau/grh-lean/blob/72fead102f663d5c7a407964a13dc0cf9d977207/grh-compact/src/New/Exchange/GridMove.lean\#L161-L174}{{\leanfont Data.\allowbreak{}endWeightData}}}
\expandafter\newcommand\csname LN@endDetExp\endcsname{\href{https://github.com/lsernau/grh-lean/blob/72fead102f663d5c7a407964a13dc0cf9d977207/grh-compact/src/New/Exchange/GridMove.lean\#L178-L178}{{\leanfont Data.\allowbreak{}endDetExp}}}
\expandafter\newcommand\csname LN@endLevel\endcsname{\href{https://github.com/lsernau/grh-lean/blob/72fead102f663d5c7a407964a13dc0cf9d977207/grh-compact/src/New/Exchange/GridMove.lean\#L182-L183}{{\leanfont Data.\allowbreak{}endLevel}}}
\expandafter\newcommand\csname LN@endLevelleprincipalLevel\endcsname{\href{https://github.com/lsernau/grh-lean/blob/72fead102f663d5c7a407964a13dc0cf9d977207/grh-compact/src/New/Exchange/GridMove.lean\#L185-L188}{{\leanfont Data.\allowbreak{}endLevel\_le\_principalLevel}}}
\expandafter\newcommand\csname LN@endKernel\endcsname{\href{https://github.com/lsernau/grh-lean/blob/72fead102f663d5c7a407964a13dc0cf9d977207/grh-compact/src/New/Exchange/GridMove.lean\#L192-L193}{{\leanfont Data.\allowbreak{}endKernel}}}
\expandafter\newcommand\csname LN@boxesendWeightDatale\endcsname{\href{https://github.com/lsernau/grh-lean/blob/72fead102f663d5c7a407964a13dc0cf9d977207/grh-compact/src/New/Exchange/GridMove.lean\#L195-L200}{{\leanfont Data.\allowbreak{}boxes\_endWeightData\_le}}}
\expandafter\newcommand\csname LN@endBranching\endcsname{\href{https://github.com/lsernau/grh-lean/blob/72fead102f663d5c7a407964a13dc0cf9d977207/grh-compact/src/New/Exchange/GridMove.lean\#L204-L206}{{\leanfont Data.\allowbreak{}endBranching}}}
\expandafter\newcommand\csname LN@endPairing\endcsname{\href{https://github.com/lsernau/grh-lean/blob/72fead102f663d5c7a407964a13dc0cf9d977207/grh-compact/src/New/Exchange/GridMove.lean\#L209-L209}{{\leanfont Data.\allowbreak{}endPairing}}}
\expandafter\newcommand\csname LN@endBlockPoint\endcsname{\href{https://github.com/lsernau/grh-lean/blob/72fead102f663d5c7a407964a13dc0cf9d977207/grh-compact/src/New/Exchange/GridMove.lean\#L211-L211}{{\leanfont Data.\allowbreak{}endBlockPoint}}}
\expandafter\newcommand\csname LN@isStdgridSpace\endcsname{\href{https://github.com/lsernau/grh-lean/blob/72fead102f663d5c7a407964a13dc0cf9d977207/grh-compact/src/New/Exchange/GridMove.lean\#L213-L215}{{\leanfont Data.\allowbreak{}isStd\_gridSpace}}}
\expandafter\newcommand\csname LN@endToAlgForm\endcsname{\href{https://github.com/lsernau/grh-lean/blob/72fead102f663d5c7a407964a13dc0cf9d977207/grh-compact/src/New/Exchange/GridMove.lean\#L219-L219}{{\leanfont Data.\allowbreak{}endToAlgForm}}}
\expandafter\newcommand\csname LN@FixedByEndLevels\endcsname{\href{https://github.com/lsernau/grh-lean/blob/72fead102f663d5c7a407964a13dc0cf9d977207/grh-compact/src/New/Exchange/GridMove.lean\#L222-L223}{{\leanfont Data.\allowbreak{}FixedByEndLevels}}}
\expandafter\newcommand\csname LN@endDatum\endcsname{\href{https://github.com/lsernau/grh-lean/blob/72fead102f663d5c7a407964a13dc0cf9d977207/grh-compact/src/New/Exchange/GridMove.lean\#L227-L228}{{\leanfont Data.\allowbreak{}endDatum}}}
\expandafter\newcommand\csname LN@endChain\endcsname{\href{https://github.com/lsernau/grh-lean/blob/72fead102f663d5c7a407964a13dc0cf9d977207/grh-compact/src/New/Exchange/GridMove.lean\#L230-L230}{{\leanfont Data.\allowbreak{}endChain}}}
\expandafter\newcommand\csname LN@endPeriod\endcsname{\href{https://github.com/lsernau/grh-lean/blob/72fead102f663d5c7a407964a13dc0cf9d977207/grh-compact/src/New/Exchange/GridMove.lean\#L234-L235}{{\leanfont Data.\allowbreak{}endPeriod}}}
\expandafter\newcommand\csname LN@existssquaremove\endcsname{\href{https://github.com/lsernau/grh-lean/blob/72fead102f663d5c7a407964a13dc0cf9d977207/grh-compact/src/New/Exchange/GridMove.lean\#L242-L293}{{\leanfont Data.\allowbreak{}IsTidy.\allowbreak{}exists\_square\_move}}}
\expandafter\newcommand\csname LN@GridIndex\endcsname{\href{https://github.com/lsernau/grh-lean/blob/72fead102f663d5c7a407964a13dc0cf9d977207/grh-compact/src/New/Exchange/GridParameters.lean\#L18-L19}{{\leanfont Data.\allowbreak{}GridIndex}}}
\expandafter\newcommand\csname LN@gridNumber\endcsname{\href{https://github.com/lsernau/grh-lean/blob/72fead102f663d5c7a407964a13dc0cf9d977207/grh-compact/src/New/Exchange/GridParameters.lean\#L23-L24}{{\leanfont Data.\allowbreak{}gridNumber}}}
\expandafter\newcommand\csname LN@gridBlockEquiv\endcsname{\href{https://github.com/lsernau/grh-lean/blob/72fead102f663d5c7a407964a13dc0cf9d977207/grh-compact/src/New/Exchange/GridParameters.lean\#L28-L28}{{\leanfont Data.\allowbreak{}gridBlockEquiv}}}
\expandafter\newcommand\csname LN@entrySigneq\endcsname{\href{https://github.com/lsernau/grh-lean/blob/72fead102f663d5c7a407964a13dc0cf9d977207/grh-compact/src/New/Exchange/GridParameters.lean\#L32-L33}{{\leanfont Data.\allowbreak{}entrySign\_eq}}}
\expandafter\newcommand\csname LN@gridSigns\endcsname{\href{https://github.com/lsernau/grh-lean/blob/72fead102f663d5c7a407964a13dc0cf9d977207/grh-compact/src/New/Exchange/GridParameters.lean\#L37-L38}{{\leanfont Data.\allowbreak{}gridSigns}}}
\expandafter\newcommand\csname LN@gridParam\endcsname{\href{https://github.com/lsernau/grh-lean/blob/72fead102f663d5c7a407964a13dc0cf9d977207/grh-compact/src/New/Exchange/GridParameters.lean\#L43-L46}{{\leanfont Data.\allowbreak{}gridParam}}}
\expandafter\newcommand\csname LN@signsY\endcsname{\href{https://github.com/lsernau/grh-lean/blob/72fead102f663d5c7a407964a13dc0cf9d977207/grh-compact/src/New/Exchange/GridParameters.lean\#L50-L50}{{\leanfont Data.\allowbreak{}signsY}}}
\expandafter\newcommand\csname LN@paramY\endcsname{\href{https://github.com/lsernau/grh-lean/blob/72fead102f663d5c7a407964a13dc0cf9d977207/grh-compact/src/New/Exchange/GridParameters.lean\#L54-L54}{{\leanfont Data.\allowbreak{}paramY}}}
\expandafter\newcommand\csname LN@gridWeight\endcsname{\href{https://github.com/lsernau/grh-lean/blob/72fead102f663d5c7a407964a13dc0cf9d977207/grh-compact/src/New/Exchange/GridParameters.lean\#L60-L60}{{\leanfont Data.\allowbreak{}gridWeight}}}
\expandafter\newcommand\csname LN@gridHarmonic\endcsname{\href{https://github.com/lsernau/grh-lean/blob/72fead102f663d5c7a407964a13dc0cf9d977207/grh-compact/src/New/Exchange/GridParameters.lean\#L64-L66}{{\leanfont Data.\allowbreak{}gridHarmonic}}}
\expandafter\newcommand\csname LN@cardgridIndex\endcsname{\href{https://github.com/lsernau/grh-lean/blob/72fead102f663d5c7a407964a13dc0cf9d977207/grh-compact/src/New/Exchange/GridParameters.lean\#L70-L71}{{\leanfont Data.\allowbreak{}card\_gridIndex}}}
\expandafter\newcommand\csname LN@cardneggridSignseqsum\endcsname{\href{https://github.com/lsernau/grh-lean/blob/72fead102f663d5c7a407964a13dc0cf9d977207/grh-compact/src/New/Exchange/GridParameters.lean\#L73-L75}{{\leanfont Data.\allowbreak{}card\_neg\_gridSigns\_eq\_sum}}}
\expandafter\newcommand\csname LN@cardneggridSigns\endcsname{\href{https://github.com/lsernau/grh-lean/blob/72fead102f663d5c7a407964a13dc0cf9d977207/grh-compact/src/New/Exchange/GridParameters.lean\#L81-L83}{{\leanfont Data.\allowbreak{}card\_neg\_gridSigns}}}
\expandafter\newcommand\csname LN@gridSignszero\endcsname{\href{https://github.com/lsernau/grh-lean/blob/72fead102f663d5c7a407964a13dc0cf9d977207/grh-compact/src/New/Exchange/GridParameters.lean\#L86-L87}{{\leanfont Data.\allowbreak{}gridSigns\_zero}}}
\expandafter\newcommand\csname LN@gridDiagEquiv\endcsname{\href{https://github.com/lsernau/grh-lean/blob/72fead102f663d5c7a407964a13dc0cf9d977207/grh-compact/src/New/Exchange/GridParameters.lean\#L91-L95}{{\leanfont Data.\allowbreak{}gridDiagEquiv}}}
\expandafter\newcommand\csname LN@gridDiagEquivspec\endcsname{\href{https://github.com/lsernau/grh-lean/blob/72fead102f663d5c7a407964a13dc0cf9d977207/grh-compact/src/New/Exchange/GridParameters.lean\#L99-L101}{{\leanfont Data.\allowbreak{}gridDiagEquiv\_spec}}}
\expandafter\newcommand\csname LN@cardgridIndexmodtwo\endcsname{\href{https://github.com/lsernau/grh-lean/blob/72fead102f663d5c7a407964a13dc0cf9d977207/grh-compact/src/New/Exchange/GridParameters.lean\#L105-L106}{{\leanfont Data.\allowbreak{}card\_gridIndex\_mod\_two}}}
\expandafter\newcommand\csname LN@waddgaple\endcsname{\href{https://github.com/lsernau/grh-lean/blob/72fead102f663d5c7a407964a13dc0cf9d977207/grh-compact/src/New/Exchange/GridParameters.lean\#L108-L109}{{\leanfont Data.\allowbreak{}w\_add\_gap\_le}}}
\expandafter\newcommand\csname LN@gridSignseqone\endcsname{\href{https://github.com/lsernau/grh-lean/blob/72fead102f663d5c7a407964a13dc0cf9d977207/grh-compact/src/New/Exchange/GridParameters.lean\#L116-L117}{{\leanfont Data.\allowbreak{}gridSigns\_eq\_one}}}
\expandafter\newcommand\csname LN@hasGapgridParam\endcsname{\href{https://github.com/lsernau/grh-lean/blob/72fead102f663d5c7a407964a13dc0cf9d977207/grh-compact/src/New/Exchange/GridParameters.lean\#L121-L124}{{\leanfont Data.\allowbreak{}hasGap\_gridParam}}}
\expandafter\newcommand\csname LN@gridParaminterlace\endcsname{\href{https://github.com/lsernau/grh-lean/blob/72fead102f663d5c7a407964a13dc0cf9d977207/grh-compact/src/New/Exchange/GridParameters.lean\#L128-L140}{{\leanfont Data.\allowbreak{}gridParam\_interlace}}}
\expandafter\newcommand\csname LN@absgridParamle\endcsname{\href{https://github.com/lsernau/grh-lean/blob/72fead102f663d5c7a407964a13dc0cf9d977207/grh-compact/src/New/Exchange/GridParameters.lean\#L142-L144}{{\leanfont Data.\allowbreak{}abs\_gridParam\_le}}}
\expandafter\newcommand\csname LN@cardgridIndexle\endcsname{\href{https://github.com/lsernau/grh-lean/blob/72fead102f663d5c7a407964a13dc0cf9d977207/grh-compact/src/New/Exchange/GridParameters.lean\#L146-L147}{{\leanfont Data.\allowbreak{}card\_gridIndex\_le}}}
\expandafter\newcommand\csname LN@paramsCorrespondgrid\endcsname{\href{https://github.com/lsernau/grh-lean/blob/72fead102f663d5c7a407964a13dc0cf9d977207/grh-compact/src/New/Exchange/GridParameters.lean\#L154-L165}{{\leanfont Data.\allowbreak{}paramsCorrespond\_grid}}}
\expandafter\newcommand\csname LN@paramsCorrespondfirstarrow\endcsname{\href{https://github.com/lsernau/grh-lean/blob/72fead102f663d5c7a407964a13dc0cf9d977207/grh-compact/src/New/Exchange/GridParameters.lean\#L167-L176}{{\leanfont Data.\allowbreak{}paramsCorrespond\_first\_arrow}}}
\expandafter\newcommand\csname LN@finitepair\endcsname{\href{https://github.com/lsernau/grh-lean/blob/72fead102f663d5c7a407964a13dc0cf9d977207/grh-compact/src/New/Exchange/LValueRatio.lean\#L14-L18}{{\leanfont Data.\allowbreak{}finite\_pair}}}
\expandafter\newcommand\csname LN@pow\endcsname{\href{https://github.com/lsernau/grh-lean/blob/72fead102f663d5c7a407964a13dc0cf9d977207/grh-compact/src/New/Exchange/LValueRatio.lean\#L22-L23}{{\leanfont Data.\allowbreak{}TameChar.\allowbreak{}pow}}}
\expandafter\newcommand\csname LN@MFoldTwist\endcsname{\href{https://github.com/lsernau/grh-lean/blob/72fead102f663d5c7a407964a13dc0cf9d977207/grh-compact/src/New/Exchange/LValueRatio.lean\#L27-L28}{{\leanfont Data.\allowbreak{}MFoldTwist}}}
\expandafter\newcommand\csname LN@chainChar\endcsname{\href{https://github.com/lsernau/grh-lean/blob/72fead102f663d5c7a407964a13dc0cf9d977207/grh-compact/src/New/Exchange/LValueRatio.lean\#L32-L32}{{\leanfont Data.\allowbreak{}chainChar}}}
\expandafter\newcommand\csname LN@instOneTws\endcsname{\href{https://github.com/lsernau/grh-lean/blob/72fead102f663d5c7a407964a13dc0cf9d977207/grh-compact/src/New/Exchange/LValueRatio.lean\#L34-L34}{{\leanfont Data.\allowbreak{}instOneTws}}}
\expandafter\newcommand\csname LN@instCoeTwsTwist\endcsname{\href{https://github.com/lsernau/grh-lean/blob/72fead102f663d5c7a407964a13dc0cf9d977207/grh-compact/src/New/Exchange/LValueRatio.lean\#L36-L36}{{\leanfont Data.\allowbreak{}instCoeTwsTwist}}}
\expandafter\newcommand\csname LN@powTwist\endcsname{\href{https://github.com/lsernau/grh-lean/blob/72fead102f663d5c7a407964a13dc0cf9d977207/grh-compact/src/New/Exchange/LValueRatio.lean\#L39-L40}{{\leanfont Data.\allowbreak{}MFoldTwist.\allowbreak{}powTwist}}}
\expandafter\newcommand\csname LN@twistonespec\endcsname{\href{https://github.com/lsernau/grh-lean/blob/72fead102f663d5c7a407964a13dc0cf9d977207/grh-compact/src/New/Exchange/LValueRatio.lean\#L42-L43}{{\leanfont Data.\allowbreak{}twist\_one\_spec}}}
\expandafter\newcommand\csname LN@ratio\endcsname{\href{https://github.com/lsernau/grh-lean/blob/72fead102f663d5c7a407964a13dc0cf9d977207/grh-compact/src/New/Exchange/LValueRatio.lean\#L47-L48}{{\leanfont Data.\allowbreak{}Pair.\allowbreak{}ratio}}}
\expandafter\newcommand\csname LN@existssquareLawgrid\endcsname{\href{https://github.com/lsernau/grh-lean/blob/72fead102f663d5c7a407964a13dc0cf9d977207/grh-compact/src/New/Exchange/RowLadder.lean\#L16-L34}{{\leanfont Data.\allowbreak{}exists\_squareLaw\_grid}}}
\expandafter\newcommand\csname LN@descentpastgridPathMap\endcsname{\href{https://github.com/lsernau/grh-lean/blob/72fead102f663d5c7a407964a13dc0cf9d977207/grh-compact/src/New/Exchange/RowLadder.lean\#L40-L62}{{\leanfont Data.\allowbreak{}descent\_past\_gridPathMap}}}
\expandafter\newcommand\csname LN@iotaPlusadjointweightedChain\endcsname{\href{https://github.com/lsernau/grh-lean/blob/72fead102f663d5c7a407964a13dc0cf9d977207/grh-compact/src/New/Exchange/RowsReached.lean\#L21-L34}{{\leanfont Data.\allowbreak{}iotaPlus\_adjoint\_weightedChain}}}
\expandafter\newcommand\csname LN@gridLevelFormseqmapiotaPlus\endcsname{\href{https://github.com/lsernau/grh-lean/blob/72fead102f663d5c7a407964a13dc0cf9d977207/grh-compact/src/New/Exchange/RowsReached.lean\#L39-L56}{{\leanfont Data.\allowbreak{}gridLevelForms\_eq\_map\_iotaPlus}}}
\expandafter\newcommand\csname LN@eqzeroofpeterssonrestricteqzero\endcsname{\href{https://github.com/lsernau/grh-lean/blob/72fead102f663d5c7a407964a13dc0cf9d977207/grh-compact/src/New/Exchange/RowsReached.lean\#L67-L135}{{\leanfont eq\_zero\_of\_petersson\_restrict\_eq\_zero}}}
\expandafter\newcommand\csname LN@eqzerooforthogonalgridLevelForms\endcsname{\href{https://github.com/lsernau/grh-lean/blob/72fead102f663d5c7a407964a13dc0cf9d977207/grh-compact/src/New/Exchange/RowsReached.lean\#L142-L162}{{\leanfont Data.\allowbreak{}eq\_zero\_of\_orthogonal\_gridLevelForms}}}
\expandafter\newcommand\csname LN@existsisSteep\endcsname{\href{https://github.com/lsernau/grh-lean/blob/72fead102f663d5c7a407964a13dc0cf9d977207/grh-compact/src/New/Exchange/SteepParameters.lean\#L13-L29}{{\leanfont Data.\allowbreak{}exists\_isSteep}}}
\newcommand{\ConcordanceRows}{\crow{existsisTidyintegraleqgridPeriod}{\href{https://github.com/lsernau/grh-lean/blob/72fead102f663d5c7a407964a13dc0cf9d977207/grh-compact/src/New/Exchange/CrudeRows.lean\#L23-L40}{{\leanfont Data.\allowbreak{}exists\_isTidy\_integral\_eq\_gridPeriod}}}{CrudeRows.lean}{23--40}
\crow{HasRoute}{\href{https://github.com/lsernau/grh-lean/blob/72fead102f663d5c7a407964a13dc0cf9d977207/grh-compact/src/New/Exchange/CrudeRows.lean\#L44-L47}{{\leanfont Data.\allowbreak{}HasRoute}}}{CrudeRows.lean}{44--47}
\crow{hasRoutesucc}{\href{https://github.com/lsernau/grh-lean/blob/72fead102f663d5c7a407964a13dc0cf9d977207/grh-compact/src/New/Exchange/CrudeRows.lean\#L52-L104}{{\leanfont Data.\allowbreak{}hasRoute\_succ}}}{CrudeRows.lean}{52--104}
\crow{gridPeriodclimb}{\href{https://github.com/lsernau/grh-lean/blob/72fead102f663d5c7a407964a13dc0cf9d977207/grh-compact/src/New/Exchange/CrudeRows.lean\#L109-L116}{{\leanfont Data.\allowbreak{}gridPeriod\_climb}}}{CrudeRows.lean}{109--116}
\crow{existsmulendPeriodmemembValRing}{\href{https://github.com/lsernau/grh-lean/blob/72fead102f663d5c7a407964a13dc0cf9d977207/grh-compact/src/New/Exchange/CrudeRows.lean\#L120-L151}{{\leanfont Data.\allowbreak{}exists\_mul\_endPeriod\_mem\_embValRing}}}{CrudeRows.lean}{120--151}
\crow{endPeriodle}{\href{https://github.com/lsernau/grh-lean/blob/72fead102f663d5c7a407964a13dc0cf9d977207/grh-compact/src/New/Exchange/CrudeRows.lean\#L156-L184}{{\leanfont Data.\allowbreak{}endPeriod\_le}}}{CrudeRows.lean}{156--184}
\crow{cruderows}{\href{https://github.com/lsernau/grh-lean/blob/72fead102f663d5c7a407964a13dc0cf9d977207/grh-compact/src/New/Exchange/CrudeRows.lean\#L191-L247}{{\leanfont Data.\allowbreak{}crude\_rows}}}{CrudeRows.lean}{191--247}
\crow{triangleOrdereqmullevelIndex}{\href{https://github.com/lsernau/grh-lean/blob/72fead102f663d5c7a407964a13dc0cf9d977207/grh-compact/src/New/Exchange/GridFirstColumns.lean\#L22-L30}{{\leanfont Data.\allowbreak{}triangleOrder\_eq\_mul\_levelIndex}}}{GridFirstColumns.lean}{22--30}
\crow{diagonalPeriodFormformRep}{\href{https://github.com/lsernau/grh-lean/blob/72fead102f663d5c7a407964a13dc0cf9d977207/grh-compact/src/New/Exchange/GridFirstColumns.lean\#L35-L42}{{\leanfont Data.\allowbreak{}diagonalPeriodForm\_formRep}}}{GridFirstColumns.lean}{35--42}
\crow{existsrestrictbordered}{\href{https://github.com/lsernau/grh-lean/blob/72fead102f663d5c7a407964a13dc0cf9d977207/grh-compact/src/New/Exchange/GridFirstColumns.lean\#L47-L70}{{\leanfont Data.\allowbreak{}exists\_restrict\_bordered}}}{GridFirstColumns.lean}{47--70}
\crow{formReplegDescentstupleAction}{\href{https://github.com/lsernau/grh-lean/blob/72fead102f663d5c7a407964a13dc0cf9d977207/grh-compact/src/New/Exchange/GridFirstColumns.lean\#L79-L92}{{\leanfont Data.\allowbreak{}formRep\_legDescents\_tupleAction}}}{GridFirstColumns.lean}{79--92}
\crow{diagonalPeriodFormtupleAction}{\href{https://github.com/lsernau/grh-lean/blob/72fead102f663d5c7a407964a13dc0cf9d977207/grh-compact/src/New/Exchange/GridFirstColumns.lean\#L97-L106}{{\leanfont Data.\allowbreak{}diagonalPeriodForm\_tupleAction}}}{GridFirstColumns.lean}{97--106}
\crow{gridPeriodeqendPeriod}{\href{https://github.com/lsernau/grh-lean/blob/72fead102f663d5c7a407964a13dc0cf9d977207/grh-compact/src/New/Exchange/GridFirstColumns.lean\#L110-L134}{{\leanfont Data.\allowbreak{}gridPeriod\_eq\_endPeriod}}}{GridFirstColumns.lean}{110--134}
\crow{endPeriodmem}{\href{https://github.com/lsernau/grh-lean/blob/72fead102f663d5c7a407964a13dc0cf9d977207/grh-compact/src/New/Exchange/GridFirstColumns.lean\#L138-L193}{{\leanfont Data.\allowbreak{}endPeriod\_mem}}}{GridFirstColumns.lean}{138--193}
\crow{existsisLevelTidyiotaPluseq}{\href{https://github.com/lsernau/grh-lean/blob/72fead102f663d5c7a407964a13dc0cf9d977207/grh-compact/src/New/Exchange/GridFirstColumns.lean\#L199-L237}{{\leanfont Data.\allowbreak{}exists\_isLevelTidy\_iotaPlus\_eq}}}{GridFirstColumns.lean}{199--237}
\crow{existsisTidydiagonalPeriodFormeq}{\href{https://github.com/lsernau/grh-lean/blob/72fead102f663d5c7a407964a13dc0cf9d977207/grh-compact/src/New/Exchange/GridFirstColumns.lean\#L242-L272}{{\leanfont Data.\allowbreak{}exists\_isTidy\_diagonalPeriodForm\_eq}}}{GridFirstColumns.lean}{242--272}
\crow{fitsParamsauxPaira}{\href{https://github.com/lsernau/grh-lean/blob/72fead102f663d5c7a407964a13dc0cf9d977207/grh-compact/src/New/Exchange/GridModuli.lean\#L14-L18}{{\leanfont LValueData.\allowbreak{}fitsParams\_auxPair}}}{GridModuli.lean}{14--18}
\crow{heckePolyOnFormstestValueeqzero}{\href{https://github.com/lsernau/grh-lean/blob/72fead102f663d5c7a407964a13dc0cf9d977207/grh-compact/src/New/Exchange/GridModuli.lean\#L24-L28}{{\leanfont AuxPair.\allowbreak{}heckePolyOnForms\_testValue\_eq\_zero}}}{GridModuli.lean}{24--28}
\crow{nStar}{\href{https://github.com/lsernau/grh-lean/blob/72fead102f663d5c7a407964a13dc0cf9d977207/grh-compact/src/New/Exchange/GridModuli.lean\#L33-L33}{{\leanfont Data.\allowbreak{}nStar}}}{GridModuli.lean}{33--33}
\crow{nStarbounds}{\href{https://github.com/lsernau/grh-lean/blob/72fead102f663d5c7a407964a13dc0cf9d977207/grh-compact/src/New/Exchange/GridModuli.lean\#L35-L36}{{\leanfont Data.\allowbreak{}nStar\_bounds}}}{GridModuli.lean}{35--36}
\crow{orderW}{\href{https://github.com/lsernau/grh-lean/blob/72fead102f663d5c7a407964a13dc0cf9d977207/grh-compact/src/New/Exchange/GridModuli.lean\#L42-L43}{{\leanfont Data.\allowbreak{}orderW}}}{GridModuli.lean}{42--43}
\crow{conjHolWeight}{\href{https://github.com/lsernau/grh-lean/blob/72fead102f663d5c7a407964a13dc0cf9d977207/grh-compact/src/New/Exchange/GridModuli.lean\#L47-L47}{{\leanfont Data.\allowbreak{}conjHolWeight}}}{GridModuli.lean}{47--47}
\crow{kernelWeight}{\href{https://github.com/lsernau/grh-lean/blob/72fead102f663d5c7a407964a13dc0cf9d977207/grh-compact/src/New/Exchange/GridModuli.lean\#L51-L52}{{\leanfont Data.\allowbreak{}kernelWeight}}}{GridModuli.lean}{51--52}
\crow{rankboundscardbigIndex}{\href{https://github.com/lsernau/grh-lean/blob/72fead102f663d5c7a407964a13dc0cf9d977207/grh-compact/src/New/Exchange/GridModuli.lean\#L59-L61}{{\leanfont Data.\allowbreak{}rank\_bounds\_card\_bigIndex}}}{GridModuli.lean}{59--61}
\crow{twomulcarddetIndexaddle}{\href{https://github.com/lsernau/grh-lean/blob/72fead102f663d5c7a407964a13dc0cf9d977207/grh-compact/src/New/Exchange/GridModuli.lean\#L63-L65}{{\leanfont Data.\allowbreak{}two\_mul\_card\_detIndex\_add\_le}}}{GridModuli.lean}{63--65}
\crow{conjHolWeighteqaddkernelWeight}{\href{https://github.com/lsernau/grh-lean/blob/72fead102f663d5c7a407964a13dc0cf9d977207/grh-compact/src/New/Exchange/GridModuli.lean\#L70-L76}{{\leanfont Data.\allowbreak{}conjHolWeight\_eq\_add\_kernelWeight}}}{GridModuli.lean}{70--76}
\crow{conjHolWeightkernelWeightbounds}{\href{https://github.com/lsernau/grh-lean/blob/72fead102f663d5c7a407964a13dc0cf9d977207/grh-compact/src/New/Exchange/GridModuli.lean\#L78-L87}{{\leanfont Data.\allowbreak{}conjHolWeight\_kernelWeight\_bounds}}}{GridModuli.lean}{78--87}
\crow{isPosFirstOrderorderW}{\href{https://github.com/lsernau/grh-lean/blob/72fead102f663d5c7a407964a13dc0cf9d977207/grh-compact/src/New/Exchange/GridModuli.lean\#L91-L92}{{\leanfont Data.\allowbreak{}isPosFirstOrder\_orderW}}}{GridModuli.lean}{91--92}
\crow{paramWltoforderWlt}{\href{https://github.com/lsernau/grh-lean/blob/72fead102f663d5c7a407964a13dc0cf9d977207/grh-compact/src/New/Exchange/GridModuli.lean\#L94-L95}{{\leanfont Data.\allowbreak{}paramW\_lt\_of\_orderW\_lt}}}{GridModuli.lean}{94--95}
\crow{instFactNeRingOfIntegersEHPowGbOfNat}{\href{https://github.com/lsernau/grh-lean/blob/72fead102f663d5c7a407964a13dc0cf9d977207/grh-compact/src/New/Exchange/GridModuli.lean\#L97-L97}{{\leanfont Data.\allowbreak{}instFactNeRingOfIntegersEHPowGbOfNat}}}{GridModuli.lean}{97--97}
\crow{placesmodulusequnion}{\href{https://github.com/lsernau/grh-lean/blob/72fead102f663d5c7a407964a13dc0cf9d977207/grh-compact/src/New/Exchange/GridModuli.lean\#L102-L103}{{\leanfont Data.\allowbreak{}places\_modulus\_eq\_union}}}{GridModuli.lean}{102--103}
\crow{isTameForofhasWeight}{\href{https://github.com/lsernau/grh-lean/blob/72fead102f663d5c7a407964a13dc0cf9d977207/grh-compact/src/New/Exchange/GridModuli.lean\#L105-L108}{{\leanfont Data.\allowbreak{}isTameFor\_of\_hasWeight}}}{GridModuli.lean}{105--108}
\crow{isTameFortwistedChi}{\href{https://github.com/lsernau/grh-lean/blob/72fead102f663d5c7a407964a13dc0cf9d977207/grh-compact/src/New/Exchange/GridModuli.lean\#L113-L115}{{\leanfont Data.\allowbreak{}isTameFor\_twistedChi}}}{GridModuli.lean}{113--115}
\crow{instNeZeroNatB}{\href{https://github.com/lsernau/grh-lean/blob/72fead102f663d5c7a407964a13dc0cf9d977207/grh-compact/src/New/Exchange/GridModuli.lean\#L117-L117}{{\leanfont Data.\allowbreak{}instNeZeroNatB₁}}}{GridModuli.lean}{117--117}
\crow{gridLevelGen}{\href{https://github.com/lsernau/grh-lean/blob/72fead102f663d5c7a407964a13dc0cf9d977207/grh-compact/src/New/Exchange/GridModuli.lean\#L120-L120}{{\leanfont Data.\allowbreak{}gridLevelGen}}}{GridModuli.lean}{120--120}
\crow{gridModulus}{\href{https://github.com/lsernau/grh-lean/blob/72fead102f663d5c7a407964a13dc0cf9d977207/grh-compact/src/New/Exchange/GridModuli.lean\#L124-L124}{{\leanfont Data.\allowbreak{}gridModulus}}}{GridModuli.lean}{124--124}
\crow{gridModulusoneideal}{\href{https://github.com/lsernau/grh-lean/blob/72fead102f663d5c7a407964a13dc0cf9d977207/grh-compact/src/New/Exchange/GridModuli.lean\#L128-L129}{{\leanfont Data.\allowbreak{}gridModulus\_one\_ideal}}}{GridModuli.lean}{128--129}
\crow{triangleModulus}{\href{https://github.com/lsernau/grh-lean/blob/72fead102f663d5c7a407964a13dc0cf9d977207/grh-compact/src/New/Exchange/GridModuli.lean\#L133-L133}{{\leanfont Data.\allowbreak{}triangleModulus}}}{GridModuli.lean}{133--133}
\crow{reduceToTriangle}{\href{https://github.com/lsernau/grh-lean/blob/72fead102f663d5c7a407964a13dc0cf9d977207/grh-compact/src/New/Exchange/GridModuli.lean\#L137-L140}{{\leanfont Data.\allowbreak{}reduceToTriangle}}}{GridModuli.lean}{137--140}
\crow{piGBardvdgridLevelGen}{\href{https://github.com/lsernau/grh-lean/blob/72fead102f663d5c7a407964a13dc0cf9d977207/grh-compact/src/New/Exchange/GridModuli.lean\#L142-L142}{{\leanfont Data.\allowbreak{}piGBar\_dvd\_gridLevelGen}}}{GridModuli.lean}{142--142}
\crow{EdgeTuple}{\href{https://github.com/lsernau/grh-lean/blob/72fead102f663d5c7a407964a13dc0cf9d977207/grh-compact/src/New/Exchange/GridMove.lean\#L21-L22}{{\leanfont Data.\allowbreak{}EdgeTuple}}}{GridMove.lean}{21--22}
\crow{GridDatum}{\href{https://github.com/lsernau/grh-lean/blob/72fead102f663d5c7a407964a13dc0cf9d977207/grh-compact/src/New/Exchange/GridMove.lean\#L26-L26}{{\leanfont Data.\allowbreak{}GridDatum}}}{GridMove.lean}{26--26}
\crow{IsMuEdge}{\href{https://github.com/lsernau/grh-lean/blob/72fead102f663d5c7a407964a13dc0cf9d977207/grh-compact/src/New/Exchange/GridMove.lean\#L30-L30}{{\leanfont Data.\allowbreak{}IsMuEdge}}}{GridMove.lean}{30--30}
\crow{edgeChar}{\href{https://github.com/lsernau/grh-lean/blob/72fead102f663d5c7a407964a13dc0cf9d977207/grh-compact/src/New/Exchange/GridMove.lean\#L34-L35}{{\leanfont Data.\allowbreak{}edgeChar}}}{GridMove.lean}{34--35}
\crow{IsAdmissibleFamily}{\href{https://github.com/lsernau/grh-lean/blob/72fead102f663d5c7a407964a13dc0cf9d977207/grh-compact/src/New/Exchange/GridMove.lean\#L39-L41}{{\leanfont Data.\allowbreak{}IsAdmissibleFamily}}}{GridMove.lean}{39--41}
\crow{tupleAction}{\href{https://github.com/lsernau/grh-lean/blob/72fead102f663d5c7a407964a13dc0cf9d977207/grh-compact/src/New/Exchange/GridMove.lean\#L44-L44}{{\leanfont Data.\allowbreak{}tupleAction}}}{GridMove.lean}{44--44}
\crow{edgeCharProd}{\href{https://github.com/lsernau/grh-lean/blob/72fead102f663d5c7a407964a13dc0cf9d977207/grh-compact/src/New/Exchange/GridMove.lean\#L48-L49}{{\leanfont Data.\allowbreak{}edgeCharProd}}}{GridMove.lean}{48--49}
\crow{tupleBound}{\href{https://github.com/lsernau/grh-lean/blob/72fead102f663d5c7a407964a13dc0cf9d977207/grh-compact/src/New/Exchange/GridMove.lean\#L53-L53}{{\leanfont Data.\allowbreak{}tupleBound}}}{GridMove.lean}{53--53}
\crow{LeftInvOnMuEdges}{\href{https://github.com/lsernau/grh-lean/blob/72fead102f663d5c7a407964a13dc0cf9d977207/grh-compact/src/New/Exchange/GridMove.lean\#L57-L58}{{\leanfont Data.\allowbreak{}LeftInvOnMuEdges}}}{GridMove.lean}{57--58}
\crow{IsShortSum}{\href{https://github.com/lsernau/grh-lean/blob/72fead102f663d5c7a407964a13dc0cf9d977207/grh-compact/src/New/Exchange/GridMove.lean\#L62-L64}{{\leanfont Data.\allowbreak{}IsShortSum}}}{GridMove.lean}{62--64}
\crow{bind}{\href{https://github.com/lsernau/grh-lean/blob/72fead102f663d5c7a407964a13dc0cf9d977207/grh-compact/src/New/Exchange/GridMove.lean\#L68-L69}{{\leanfont Data.\allowbreak{}GridDatum.\allowbreak{}bind}}}{GridMove.lean}{68--69}
\crow{IsTidy}{\href{https://github.com/lsernau/grh-lean/blob/72fead102f663d5c7a407964a13dc0cf9d977207/grh-compact/src/New/Exchange/GridMove.lean\#L73-L76}{{\leanfont Data.\allowbreak{}IsTidy}}}{GridMove.lean}{73--76}
\crow{IsLevelTidy}{\href{https://github.com/lsernau/grh-lean/blob/72fead102f663d5c7a407964a13dc0cf9d977207/grh-compact/src/New/Exchange/GridMove.lean\#L80-L81}{{\leanfont Data.\allowbreak{}IsLevelTidy}}}{GridMove.lean}{80--81}
\crow{verticalEdge}{\href{https://github.com/lsernau/grh-lean/blob/72fead102f663d5c7a407964a13dc0cf9d977207/grh-compact/src/New/Exchange/GridMove.lean\#L84-L84}{{\leanfont Data.\allowbreak{}verticalEdge}}}{GridMove.lean}{84--84}
\crow{rowEdge}{\href{https://github.com/lsernau/grh-lean/blob/72fead102f663d5c7a407964a13dc0cf9d977207/grh-compact/src/New/Exchange/GridMove.lean\#L88-L88}{{\leanfont Data.\allowbreak{}rowEdge}}}{GridMove.lean}{88--88}
\crow{legEdges}{\href{https://github.com/lsernau/grh-lean/blob/72fead102f663d5c7a407964a13dc0cf9d977207/grh-compact/src/New/Exchange/GridMove.lean\#L91-L91}{{\leanfont Data.\allowbreak{}legEdges}}}{GridMove.lean}{91--91}
\crow{gridPath}{\href{https://github.com/lsernau/grh-lean/blob/72fead102f663d5c7a407964a13dc0cf9d977207/grh-compact/src/New/Exchange/GridMove.lean\#L95-L95}{{\leanfont Data.\allowbreak{}gridPath}}}{GridMove.lean}{95--95}
\crow{path}{\href{https://github.com/lsernau/grh-lean/blob/72fead102f663d5c7a407964a13dc0cf9d977207/grh-compact/src/New/Exchange/GridMove.lean\#L99-L99}{{\leanfont Data.\allowbreak{}path}}}{GridMove.lean}{99--99}
\crow{legDescents}{\href{https://github.com/lsernau/grh-lean/blob/72fead102f663d5c7a407964a13dc0cf9d977207/grh-compact/src/New/Exchange/GridMove.lean\#L104-L106}{{\leanfont Data.\allowbreak{}legDescents}}}{GridMove.lean}{104--106}
\crow{gridPathMap}{\href{https://github.com/lsernau/grh-lean/blob/72fead102f663d5c7a407964a13dc0cf9d977207/grh-compact/src/New/Exchange/GridMove.lean\#L113-L114}{{\leanfont Data.\allowbreak{}gridPathMap}}}{GridMove.lean}{113--114}
\crow{diagonalPeriodForm}{\href{https://github.com/lsernau/grh-lean/blob/72fead102f663d5c7a407964a13dc0cf9d977207/grh-compact/src/New/Exchange/GridMove.lean\#L118-L121}{{\leanfont Data.\allowbreak{}diagonalPeriodForm}}}{GridMove.lean}{118--121}
\crow{gridPeriod}{\href{https://github.com/lsernau/grh-lean/blob/72fead102f663d5c7a407964a13dc0cf9d977207/grh-compact/src/New/Exchange/GridMove.lean\#L125-L127}{{\leanfont Data.\allowbreak{}gridPeriod}}}{GridMove.lean}{125--127}
\crow{edgeCharProdmul}{\href{https://github.com/lsernau/grh-lean/blob/72fead102f663d5c7a407964a13dc0cf9d977207/grh-compact/src/New/Exchange/GridMove.lean\#L129-L130}{{\leanfont Data.\allowbreak{}edgeCharProd\_mul}}}{GridMove.lean}{129--130}
\crow{legDescentscongr}{\href{https://github.com/lsernau/grh-lean/blob/72fead102f663d5c7a407964a13dc0cf9d977207/grh-compact/src/New/Exchange/GridMove.lean\#L134-L135}{{\leanfont Data.\allowbreak{}legDescents\_congr}}}{GridMove.lean}{134--135}
\crow{gridPathMapcongr}{\href{https://github.com/lsernau/grh-lean/blob/72fead102f663d5c7a407964a13dc0cf9d977207/grh-compact/src/New/Exchange/GridMove.lean\#L137-L138}{{\leanfont Data.\allowbreak{}gridPathMap\_congr}}}{GridMove.lean}{137--138}
\crow{binda}{\href{https://github.com/lsernau/grh-lean/blob/72fead102f663d5c7a407964a13dc0cf9d977207/grh-compact/src/New/Exchange/GridMove.lean\#L143-L147}{{\leanfont Data.\allowbreak{}IsShortSum.\allowbreak{}bind}}}{GridMove.lean}{143--147}
\crow{cardmatrixletupleBound}{\href{https://github.com/lsernau/grh-lean/blob/72fead102f663d5c7a407964a13dc0cf9d977207/grh-compact/src/New/Exchange/GridMove.lean\#L149-L152}{{\leanfont Data.\allowbreak{}card\_matrix\_le\_tupleBound}}}{GridMove.lean}{149--152}
\crow{EndIndex}{\href{https://github.com/lsernau/grh-lean/blob/72fead102f663d5c7a407964a13dc0cf9d977207/grh-compact/src/New/Exchange/GridMove.lean\#L156-L156}{{\leanfont Data.\allowbreak{}EndIndex}}}{GridMove.lean}{156--156}
\crow{endWeightData}{\href{https://github.com/lsernau/grh-lean/blob/72fead102f663d5c7a407964a13dc0cf9d977207/grh-compact/src/New/Exchange/GridMove.lean\#L161-L174}{{\leanfont Data.\allowbreak{}endWeightData}}}{GridMove.lean}{161--174}
\crow{endDetExp}{\href{https://github.com/lsernau/grh-lean/blob/72fead102f663d5c7a407964a13dc0cf9d977207/grh-compact/src/New/Exchange/GridMove.lean\#L178-L178}{{\leanfont Data.\allowbreak{}endDetExp}}}{GridMove.lean}{178--178}
\crow{endLevel}{\href{https://github.com/lsernau/grh-lean/blob/72fead102f663d5c7a407964a13dc0cf9d977207/grh-compact/src/New/Exchange/GridMove.lean\#L182-L183}{{\leanfont Data.\allowbreak{}endLevel}}}{GridMove.lean}{182--183}
\crow{endLevelleprincipalLevel}{\href{https://github.com/lsernau/grh-lean/blob/72fead102f663d5c7a407964a13dc0cf9d977207/grh-compact/src/New/Exchange/GridMove.lean\#L185-L188}{{\leanfont Data.\allowbreak{}endLevel\_le\_principalLevel}}}{GridMove.lean}{185--188}
\crow{endKernel}{\href{https://github.com/lsernau/grh-lean/blob/72fead102f663d5c7a407964a13dc0cf9d977207/grh-compact/src/New/Exchange/GridMove.lean\#L192-L193}{{\leanfont Data.\allowbreak{}endKernel}}}{GridMove.lean}{192--193}
\crow{boxesendWeightDatale}{\href{https://github.com/lsernau/grh-lean/blob/72fead102f663d5c7a407964a13dc0cf9d977207/grh-compact/src/New/Exchange/GridMove.lean\#L195-L200}{{\leanfont Data.\allowbreak{}boxes\_endWeightData\_le}}}{GridMove.lean}{195--200}
\crow{endBranching}{\href{https://github.com/lsernau/grh-lean/blob/72fead102f663d5c7a407964a13dc0cf9d977207/grh-compact/src/New/Exchange/GridMove.lean\#L204-L206}{{\leanfont Data.\allowbreak{}endBranching}}}{GridMove.lean}{204--206}
\crow{endPairing}{\href{https://github.com/lsernau/grh-lean/blob/72fead102f663d5c7a407964a13dc0cf9d977207/grh-compact/src/New/Exchange/GridMove.lean\#L209-L209}{{\leanfont Data.\allowbreak{}endPairing}}}{GridMove.lean}{209--209}
\crow{endBlockPoint}{\href{https://github.com/lsernau/grh-lean/blob/72fead102f663d5c7a407964a13dc0cf9d977207/grh-compact/src/New/Exchange/GridMove.lean\#L211-L211}{{\leanfont Data.\allowbreak{}endBlockPoint}}}{GridMove.lean}{211--211}
\crow{isStdgridSpace}{\href{https://github.com/lsernau/grh-lean/blob/72fead102f663d5c7a407964a13dc0cf9d977207/grh-compact/src/New/Exchange/GridMove.lean\#L213-L215}{{\leanfont Data.\allowbreak{}isStd\_gridSpace}}}{GridMove.lean}{213--215}
\crow{endToAlgForm}{\href{https://github.com/lsernau/grh-lean/blob/72fead102f663d5c7a407964a13dc0cf9d977207/grh-compact/src/New/Exchange/GridMove.lean\#L219-L219}{{\leanfont Data.\allowbreak{}endToAlgForm}}}{GridMove.lean}{219--219}
\crow{FixedByEndLevels}{\href{https://github.com/lsernau/grh-lean/blob/72fead102f663d5c7a407964a13dc0cf9d977207/grh-compact/src/New/Exchange/GridMove.lean\#L222-L223}{{\leanfont Data.\allowbreak{}FixedByEndLevels}}}{GridMove.lean}{222--223}
\crow{endDatum}{\href{https://github.com/lsernau/grh-lean/blob/72fead102f663d5c7a407964a13dc0cf9d977207/grh-compact/src/New/Exchange/GridMove.lean\#L227-L228}{{\leanfont Data.\allowbreak{}endDatum}}}{GridMove.lean}{227--228}
\crow{endChain}{\href{https://github.com/lsernau/grh-lean/blob/72fead102f663d5c7a407964a13dc0cf9d977207/grh-compact/src/New/Exchange/GridMove.lean\#L230-L230}{{\leanfont Data.\allowbreak{}endChain}}}{GridMove.lean}{230--230}
\crow{endPeriod}{\href{https://github.com/lsernau/grh-lean/blob/72fead102f663d5c7a407964a13dc0cf9d977207/grh-compact/src/New/Exchange/GridMove.lean\#L234-L235}{{\leanfont Data.\allowbreak{}endPeriod}}}{GridMove.lean}{234--235}
\crow{existssquaremove}{\href{https://github.com/lsernau/grh-lean/blob/72fead102f663d5c7a407964a13dc0cf9d977207/grh-compact/src/New/Exchange/GridMove.lean\#L242-L293}{{\leanfont Data.\allowbreak{}IsTidy.\allowbreak{}exists\_square\_move}}}{GridMove.lean}{242--293}
\crow{GridIndex}{\href{https://github.com/lsernau/grh-lean/blob/72fead102f663d5c7a407964a13dc0cf9d977207/grh-compact/src/New/Exchange/GridParameters.lean\#L18-L19}{{\leanfont Data.\allowbreak{}GridIndex}}}{GridParameters.lean}{18--19}
\crow{gridNumber}{\href{https://github.com/lsernau/grh-lean/blob/72fead102f663d5c7a407964a13dc0cf9d977207/grh-compact/src/New/Exchange/GridParameters.lean\#L23-L24}{{\leanfont Data.\allowbreak{}gridNumber}}}{GridParameters.lean}{23--24}
\crow{gridBlockEquiv}{\href{https://github.com/lsernau/grh-lean/blob/72fead102f663d5c7a407964a13dc0cf9d977207/grh-compact/src/New/Exchange/GridParameters.lean\#L28-L28}{{\leanfont Data.\allowbreak{}gridBlockEquiv}}}{GridParameters.lean}{28--28}
\crow{entrySigneq}{\href{https://github.com/lsernau/grh-lean/blob/72fead102f663d5c7a407964a13dc0cf9d977207/grh-compact/src/New/Exchange/GridParameters.lean\#L32-L33}{{\leanfont Data.\allowbreak{}entrySign\_eq}}}{GridParameters.lean}{32--33}
\crow{gridSigns}{\href{https://github.com/lsernau/grh-lean/blob/72fead102f663d5c7a407964a13dc0cf9d977207/grh-compact/src/New/Exchange/GridParameters.lean\#L37-L38}{{\leanfont Data.\allowbreak{}gridSigns}}}{GridParameters.lean}{37--38}
\crow{gridParam}{\href{https://github.com/lsernau/grh-lean/blob/72fead102f663d5c7a407964a13dc0cf9d977207/grh-compact/src/New/Exchange/GridParameters.lean\#L43-L46}{{\leanfont Data.\allowbreak{}gridParam}}}{GridParameters.lean}{43--46}
\crow{signsY}{\href{https://github.com/lsernau/grh-lean/blob/72fead102f663d5c7a407964a13dc0cf9d977207/grh-compact/src/New/Exchange/GridParameters.lean\#L50-L50}{{\leanfont Data.\allowbreak{}signsY}}}{GridParameters.lean}{50--50}
\crow{paramY}{\href{https://github.com/lsernau/grh-lean/blob/72fead102f663d5c7a407964a13dc0cf9d977207/grh-compact/src/New/Exchange/GridParameters.lean\#L54-L54}{{\leanfont Data.\allowbreak{}paramY}}}{GridParameters.lean}{54--54}
\crow{gridWeight}{\href{https://github.com/lsernau/grh-lean/blob/72fead102f663d5c7a407964a13dc0cf9d977207/grh-compact/src/New/Exchange/GridParameters.lean\#L60-L60}{{\leanfont Data.\allowbreak{}gridWeight}}}{GridParameters.lean}{60--60}
\crow{gridHarmonic}{\href{https://github.com/lsernau/grh-lean/blob/72fead102f663d5c7a407964a13dc0cf9d977207/grh-compact/src/New/Exchange/GridParameters.lean\#L64-L66}{{\leanfont Data.\allowbreak{}gridHarmonic}}}{GridParameters.lean}{64--66}
\crow{cardgridIndex}{\href{https://github.com/lsernau/grh-lean/blob/72fead102f663d5c7a407964a13dc0cf9d977207/grh-compact/src/New/Exchange/GridParameters.lean\#L70-L71}{{\leanfont Data.\allowbreak{}card\_gridIndex}}}{GridParameters.lean}{70--71}
\crow{cardneggridSignseqsum}{\href{https://github.com/lsernau/grh-lean/blob/72fead102f663d5c7a407964a13dc0cf9d977207/grh-compact/src/New/Exchange/GridParameters.lean\#L73-L75}{{\leanfont Data.\allowbreak{}card\_neg\_gridSigns\_eq\_sum}}}{GridParameters.lean}{73--75}
\crow{cardneggridSigns}{\href{https://github.com/lsernau/grh-lean/blob/72fead102f663d5c7a407964a13dc0cf9d977207/grh-compact/src/New/Exchange/GridParameters.lean\#L81-L83}{{\leanfont Data.\allowbreak{}card\_neg\_gridSigns}}}{GridParameters.lean}{81--83}
\crow{gridSignszero}{\href{https://github.com/lsernau/grh-lean/blob/72fead102f663d5c7a407964a13dc0cf9d977207/grh-compact/src/New/Exchange/GridParameters.lean\#L86-L87}{{\leanfont Data.\allowbreak{}gridSigns\_zero}}}{GridParameters.lean}{86--87}
\crow{gridDiagEquiv}{\href{https://github.com/lsernau/grh-lean/blob/72fead102f663d5c7a407964a13dc0cf9d977207/grh-compact/src/New/Exchange/GridParameters.lean\#L91-L95}{{\leanfont Data.\allowbreak{}gridDiagEquiv}}}{GridParameters.lean}{91--95}
\crow{gridDiagEquivspec}{\href{https://github.com/lsernau/grh-lean/blob/72fead102f663d5c7a407964a13dc0cf9d977207/grh-compact/src/New/Exchange/GridParameters.lean\#L99-L101}{{\leanfont Data.\allowbreak{}gridDiagEquiv\_spec}}}{GridParameters.lean}{99--101}
\crow{cardgridIndexmodtwo}{\href{https://github.com/lsernau/grh-lean/blob/72fead102f663d5c7a407964a13dc0cf9d977207/grh-compact/src/New/Exchange/GridParameters.lean\#L105-L106}{{\leanfont Data.\allowbreak{}card\_gridIndex\_mod\_two}}}{GridParameters.lean}{105--106}
\crow{waddgaple}{\href{https://github.com/lsernau/grh-lean/blob/72fead102f663d5c7a407964a13dc0cf9d977207/grh-compact/src/New/Exchange/GridParameters.lean\#L108-L109}{{\leanfont Data.\allowbreak{}w\_add\_gap\_le}}}{GridParameters.lean}{108--109}
\crow{gridSignseqone}{\href{https://github.com/lsernau/grh-lean/blob/72fead102f663d5c7a407964a13dc0cf9d977207/grh-compact/src/New/Exchange/GridParameters.lean\#L116-L117}{{\leanfont Data.\allowbreak{}gridSigns\_eq\_one}}}{GridParameters.lean}{116--117}
\crow{hasGapgridParam}{\href{https://github.com/lsernau/grh-lean/blob/72fead102f663d5c7a407964a13dc0cf9d977207/grh-compact/src/New/Exchange/GridParameters.lean\#L121-L124}{{\leanfont Data.\allowbreak{}hasGap\_gridParam}}}{GridParameters.lean}{121--124}
\crow{gridParaminterlace}{\href{https://github.com/lsernau/grh-lean/blob/72fead102f663d5c7a407964a13dc0cf9d977207/grh-compact/src/New/Exchange/GridParameters.lean\#L128-L140}{{\leanfont Data.\allowbreak{}gridParam\_interlace}}}{GridParameters.lean}{128--140}
\crow{absgridParamle}{\href{https://github.com/lsernau/grh-lean/blob/72fead102f663d5c7a407964a13dc0cf9d977207/grh-compact/src/New/Exchange/GridParameters.lean\#L142-L144}{{\leanfont Data.\allowbreak{}abs\_gridParam\_le}}}{GridParameters.lean}{142--144}
\crow{cardgridIndexle}{\href{https://github.com/lsernau/grh-lean/blob/72fead102f663d5c7a407964a13dc0cf9d977207/grh-compact/src/New/Exchange/GridParameters.lean\#L146-L147}{{\leanfont Data.\allowbreak{}card\_gridIndex\_le}}}{GridParameters.lean}{146--147}
\crow{paramsCorrespondgrid}{\href{https://github.com/lsernau/grh-lean/blob/72fead102f663d5c7a407964a13dc0cf9d977207/grh-compact/src/New/Exchange/GridParameters.lean\#L154-L165}{{\leanfont Data.\allowbreak{}paramsCorrespond\_grid}}}{GridParameters.lean}{154--165}
\crow{paramsCorrespondfirstarrow}{\href{https://github.com/lsernau/grh-lean/blob/72fead102f663d5c7a407964a13dc0cf9d977207/grh-compact/src/New/Exchange/GridParameters.lean\#L167-L176}{{\leanfont Data.\allowbreak{}paramsCorrespond\_first\_arrow}}}{GridParameters.lean}{167--176}
\crow{finitepair}{\href{https://github.com/lsernau/grh-lean/blob/72fead102f663d5c7a407964a13dc0cf9d977207/grh-compact/src/New/Exchange/LValueRatio.lean\#L14-L18}{{\leanfont Data.\allowbreak{}finite\_pair}}}{LValueRatio.lean}{14--18}
\crow{pow}{\href{https://github.com/lsernau/grh-lean/blob/72fead102f663d5c7a407964a13dc0cf9d977207/grh-compact/src/New/Exchange/LValueRatio.lean\#L22-L23}{{\leanfont Data.\allowbreak{}TameChar.\allowbreak{}pow}}}{LValueRatio.lean}{22--23}
\crow{MFoldTwist}{\href{https://github.com/lsernau/grh-lean/blob/72fead102f663d5c7a407964a13dc0cf9d977207/grh-compact/src/New/Exchange/LValueRatio.lean\#L27-L28}{{\leanfont Data.\allowbreak{}MFoldTwist}}}{LValueRatio.lean}{27--28}
\crow{chainChar}{\href{https://github.com/lsernau/grh-lean/blob/72fead102f663d5c7a407964a13dc0cf9d977207/grh-compact/src/New/Exchange/LValueRatio.lean\#L32-L32}{{\leanfont Data.\allowbreak{}chainChar}}}{LValueRatio.lean}{32--32}
\crow{instOneTws}{\href{https://github.com/lsernau/grh-lean/blob/72fead102f663d5c7a407964a13dc0cf9d977207/grh-compact/src/New/Exchange/LValueRatio.lean\#L34-L34}{{\leanfont Data.\allowbreak{}instOneTws}}}{LValueRatio.lean}{34--34}
\crow{instCoeTwsTwist}{\href{https://github.com/lsernau/grh-lean/blob/72fead102f663d5c7a407964a13dc0cf9d977207/grh-compact/src/New/Exchange/LValueRatio.lean\#L36-L36}{{\leanfont Data.\allowbreak{}instCoeTwsTwist}}}{LValueRatio.lean}{36--36}
\crow{powTwist}{\href{https://github.com/lsernau/grh-lean/blob/72fead102f663d5c7a407964a13dc0cf9d977207/grh-compact/src/New/Exchange/LValueRatio.lean\#L39-L40}{{\leanfont Data.\allowbreak{}MFoldTwist.\allowbreak{}powTwist}}}{LValueRatio.lean}{39--40}
\crow{twistonespec}{\href{https://github.com/lsernau/grh-lean/blob/72fead102f663d5c7a407964a13dc0cf9d977207/grh-compact/src/New/Exchange/LValueRatio.lean\#L42-L43}{{\leanfont Data.\allowbreak{}twist\_one\_spec}}}{LValueRatio.lean}{42--43}
\crow{ratio}{\href{https://github.com/lsernau/grh-lean/blob/72fead102f663d5c7a407964a13dc0cf9d977207/grh-compact/src/New/Exchange/LValueRatio.lean\#L47-L48}{{\leanfont Data.\allowbreak{}Pair.\allowbreak{}ratio}}}{LValueRatio.lean}{47--48}
\crow{existssquareLawgrid}{\href{https://github.com/lsernau/grh-lean/blob/72fead102f663d5c7a407964a13dc0cf9d977207/grh-compact/src/New/Exchange/RowLadder.lean\#L16-L34}{{\leanfont Data.\allowbreak{}exists\_squareLaw\_grid}}}{RowLadder.lean}{16--34}
\crow{descentpastgridPathMap}{\href{https://github.com/lsernau/grh-lean/blob/72fead102f663d5c7a407964a13dc0cf9d977207/grh-compact/src/New/Exchange/RowLadder.lean\#L40-L62}{{\leanfont Data.\allowbreak{}descent\_past\_gridPathMap}}}{RowLadder.lean}{40--62}
\crow{iotaPlusadjointweightedChain}{\href{https://github.com/lsernau/grh-lean/blob/72fead102f663d5c7a407964a13dc0cf9d977207/grh-compact/src/New/Exchange/RowsReached.lean\#L21-L34}{{\leanfont Data.\allowbreak{}iotaPlus\_adjoint\_weightedChain}}}{RowsReached.lean}{21--34}
\crow{gridLevelFormseqmapiotaPlus}{\href{https://github.com/lsernau/grh-lean/blob/72fead102f663d5c7a407964a13dc0cf9d977207/grh-compact/src/New/Exchange/RowsReached.lean\#L39-L56}{{\leanfont Data.\allowbreak{}gridLevelForms\_eq\_map\_iotaPlus}}}{RowsReached.lean}{39--56}
\crow{eqzeroofpeterssonrestricteqzero}{\href{https://github.com/lsernau/grh-lean/blob/72fead102f663d5c7a407964a13dc0cf9d977207/grh-compact/src/New/Exchange/RowsReached.lean\#L67-L135}{{\leanfont eq\_zero\_of\_petersson\_restrict\_eq\_zero}}}{RowsReached.lean}{67--135}
\crow{eqzerooforthogonalgridLevelForms}{\href{https://github.com/lsernau/grh-lean/blob/72fead102f663d5c7a407964a13dc0cf9d977207/grh-compact/src/New/Exchange/RowsReached.lean\#L142-L162}{{\leanfont Data.\allowbreak{}eq\_zero\_of\_orthogonal\_gridLevelForms}}}{RowsReached.lean}{142--162}
\crow{existsisSteep}{\href{https://github.com/lsernau/grh-lean/blob/72fead102f663d5c7a407964a13dc0cf9d977207/grh-compact/src/New/Exchange/SteepParameters.lean\#L13-L29}{{\leanfont Data.\allowbreak{}exists\_isSteep}}}{SteepParameters.lean}{13--29}}
% the notations of the defined functions: one slot for each explicit argument of the Lean declaration
\newcommand{\HasRoute}[4]{\operatorname{Route}_{#1}(#2,#3;#4)}
\newcommand{\nStar}[1]{\operatorname{nst}(#1)}
\newcommand{\orderW}[2]{\operatorname{ordW}(#1,#2)}
\newcommand{\conjHolWeight}[2]{\operatorname{chw}(#1,#2)}
\newcommand{\kernelWeight}[2]{\operatorname{kw}(#1,#2)}
\newcommand{\gridModulus}[1]{\mathfrak{m}^{\mathrm{gr}}(#1)}
\newcommand{\triangleModulus}[2]{\mathfrak{m}^{\triangle}_{#1}(#2)}
\newcommand{\reduceToTriangle}[2]{\operatorname{red}_{#1,#2}}
\newcommand{\EdgeTuple}[1]{\operatorname{Tup}(#1)}
\newcommand{\GridDatum}[1]{\operatorname{Dat}(#1)}
\newcommand{\IsMuEdge}[2]{\operatorname{MuEdge}(#1,#2)}
\newcommand{\edgeChar}[2]{\chi^{\mathrm{ed}}_{#1}(#2)}
\newcommand{\IsAdmissibleFamily}[1]{\operatorname{Adm}(#1)}
\newcommand{\tupleAction}[3]{\operatorname{act}(#1,#2)(#3)}
\newcommand{\edgeCharProd}[2]{\operatorname{ecp}(#1,#2)}
\newcommand{\tupleBound}[1]{\operatorname{tb}(#1)}
\newcommand{\LeftInvOnMuEdges}[2]{\operatorname{LInv}(#1,#2)}
\newcommand{\IsShortSum}[3]{\operatorname{Short}(#1,#2;#3)}
\newcommand{\bind}[3]{\operatorname{bind}(#1;#2,#3)}
\newcommand{\IsTidy}[3]{\operatorname{Tidy}(#1,#2;#3)}
\newcommand{\IsLevelTidy}[3]{\operatorname{LTidy}(#1,#2;#3)}
\newcommand{\verticalEdge}[3]{\operatorname{ve}(#1,#2,#3)}
\newcommand{\rowEdge}[3]{\operatorname{re}(#1,#2,#3)}
\newcommand{\legEdges}[2]{\operatorname{leg}(#1,#2)}
\newcommand{\gridPath}[3]{\operatorname{gp}(#1,#2,#3)}
\newcommand{\stagePath}[1]{\operatorname{path}(#1)}
\newcommand{\legDescents}[4]{\operatorname{ld}_{#1}(#2,#3)(#4)}
\newcommand{\gridPathMap}[5]{\operatorname{pm}_{#1}(#2;#3,#4)(#5)}
\newcommand{\diagonalPeriodForm}[3]{\operatorname{dp}_{#1}(#2;#3)}
\newcommand{\gridPeriod}[4]{\operatorname{per}(#1;#2;#3,#4)}
\newcommand{\EndIndex}[2]{\operatorname{EI}(#1,#2)}
\newcommand{\finWeightData}[1]{\operatorname{ewd}(#1)}
\newcommand{\finDetExp}[1]{\operatorname{ede}(#1)}
\newcommand{\finLevel}[3]{\operatorname{elev}(#1,#2,#3)}
\newcommand{\finKernel}[3]{\operatorname{eker}(#1,#2,#3)}
\newcommand{\finBlockPoint}[2]{\operatorname{ebp}(#1,#2)}
\newcommand{\finToAlgForm}[1]{\operatorname{alg}_{#1}}
\newcommand{\FixedByEndLevels}[2]{\operatorname{Fix}(#1,#2)}
\newcommand{\finDatum}[3]{\operatorname{edat}(#1,#2,#3)}
\newcommand{\finChain}[4]{\operatorname{ech}(#1,#2,#3)(#4)}
\newcommand{\finPeriod}[2]{\operatorname{eper}(#1;#2)}
\newcommand{\GridIndex}[4]{\operatorname{GI}_{#1}(#2;#3,#4)}
\newcommand{\gridNumber}[2]{\operatorname{num}_{#1}(#2)}
\newcommand{\gridBlockEquiv}[1]{\operatorname{gbe}_{#1}}
\newcommand{\gridSigns}[6]{\operatorname{sg}_{#1}(#2,#3;#4,#5)(#6)}
\newcommand{\gridParam}[6]{\operatorname{par}_{#1}(#2,#3;#4,#5)(#6)}
\newcommand{\signsY}[3]{\operatorname{sgY}_{#1}(#2)(#3)}
\newcommand{\paramY}[4]{\operatorname{parY}_{#1}(#2)(#3,#4)}
\newcommand{\gridWeight}[5]{\operatorname{gw}_{#1}(#2,#3;#4)(#5)}
\newcommand{\gridDiagEquiv}[2]{\operatorname{gde}_{#1}(#2)}
\newcommand{\pow}[2]{#1^{#2}}
\newcommand{\MFoldTwist}[1]{\operatorname{Tw}^{m}(#1)}
\newcommand{\chainChar}[3]{\operatorname{cch}_{#1}(#2,#3)}
\newcommand{\powTwist}[2]{#1^{[#2]}}
\newcommand{\pairRatio}[3]{\operatorname{rat}(#1,#2,#3)}
\newcommand{\UnitMembers}{305}
\newcommand{\UnitInNote}{125}
\newcommand{\UnitProofs}{122}
\newcommand{\UnitSteps}{879}
\newcommand{\UnitAllSteps}{933}


\linespread{1.08}                       % 11 pt on 14.7 pt
\setlength{\parindent}{0pt}
\setlength{\parskip}{5.5pt plus 1.5pt minus 0.5pt}
\setlength{\abovedisplayskip}{7.5pt plus 2pt minus 3pt}
\setlength{\belowdisplayskip}{7.5pt plus 2pt minus 3pt}
\setlength{\abovedisplayshortskip}{3pt plus 2pt}
\setlength{\belowdisplayshortskip}{6pt plus 2pt minus 2pt}
\emergencystretch=1.5em
\clubpenalty=10000 \widowpenalty=10000 \displaywidowpenalty=10000
\relpenalty=7000 \binoppenalty=9000
\raggedbottom
\pagestyle{fancy}\fancyhf{}\renewcommand{\headrulewidth}{0pt}
\fancyfoot[C]{\sffamily\fontsize{7.6}{9}\selectfont\color{inkfaint}\thepage\,/\,\pageref*{LastPage}}
\setlist{topsep=2pt,partopsep=0pt,parsep=0pt,itemsep=4pt,leftmargin=2.1em,labelsep=0.55em}

\newcommand{\Q}{\mathbb{Q}}
\newcommand{\Z}{\mathbb{Z}}
\newcommand{\F}{\mathbb{F}}
\newcommand{\len}{\operatorname{len}}
\newcommand{\cok}{\operatorname{cok}}
\newcommand{\End}{\operatorname{End}}
\newcommand{\Hom}{\operatorname{Hom}}
\newcommand{\Gal}{\operatorname{Gal}}
\newcommand{\diag}{\operatorname{diag}}
\newcommand{\kw}[1]{\emph{#1}}
\newcommand{\m}[1]{\mbox{$#1$}}        % a short formula that must not be broken across lines
\newcommand{\head}[1]{\par\addvspace{17pt plus 4pt minus 2pt}\Needspace{6\baselineskip}%
  \noindent{\fontsize{13.2}{16}\selectfont\bfseries\color{ink}#1}\par\nopagebreak\vspace{-1pt}}
\newcommand{\sect}[2]{\refstepcounter{section}\label{#2}\par\addvspace{19pt plus 4pt minus 2pt}\Needspace{7\baselineskip}%
  \noindent\llap{{\fontsize{13.2}{16}\selectfont\bfseries\color{inkfaint}\thesection}\hspace{9pt}}%
  {\fontsize{13.2}{16}\selectfont\bfseries\color{ink}#1}\par\nopagebreak\vspace{-1pt}}
\newcommand{\subhead}[1]{\par\addvspace{9pt plus 2pt}\Needspace{4\baselineskip}\noindent\textbf{#1}\enspace\ignorespaces}
% where a statement lives in the Lean archive
\newcommand{\leantag}{{\sffamily\bfseries\fontsize{6.3}{8}\selectfont\color{leancol}\addfontfeatures{LetterSpace=7}LEAN}}
\newcommand{\sep}{\unskip\nobreak\hspace{5pt}{\color{inkfaint}\textperiodcentered}\hspace{5pt plus 2pt}\ignorespaces}
\newcommand{\leanline}[1]{\par\nopagebreak\vspace{-2.5pt}{\fontsize{8.9}{12.4}\selectfont\color{inksoft}\raggedright
  \noindent\llap{\leantag\hspace{9pt}}#1\par}\vspace{1pt}}
% a displayed statement in a tinted box: #1 colour, #2 text
\newcommand{\statement}[2]{\par\addvspace{7pt}\begingroup\setlength{\fboxsep}{8.5pt}\setlength{\fboxrule}{0.6pt}%
  \noindent\fcolorbox{#1!30!white}{#1!3!white}{\parbox{\dimexpr\linewidth-2\fboxsep-2\fboxrule\relax}{#2}}%
  \endgroup\par\addvspace{7pt}}
\newcommand{\nt}[2]{\par\noindent\hangindent=1.1em\hangafter=1 #1\enspace{\color{inksoft}#2}\par\vspace{1pt}}
\newcommand{\legdot}[1]{\tikz[baseline=-0.62ex]\fill[#1] (0,0) circle (2.5pt);\hspace{3pt}}
\newcommand{\crow}[4]{\csname CS@#1\endcsname & #2 & \textcolor{inksoft}{#3} & \textcolor{inksoft}{#4}\\}
\newcommand{\leanquoteline}[2]{\par\noindent\hangindent=\dimexpr#1\fontdimen2\font+2.2em\relax\hangafter=1
  \hspace*{#1\fontdimen2\font}#2}
\newcommand{\quoted}[1]{\par\vspace{3pt}{\ttfamily\fontsize{7.5}{10.4}\selectfont\color{leancol}\raggedright
  \setlength{\parskip}{0pt}\leanquote{#1}\par}\vspace{4pt}}
\newcommand{\bx}[3]{{\sffamily\bfseries\fontsize{7.3}{9.2}\selectfont\color{#1!82!black}#2}\\[1.5pt]#3}
\newcommand{\cs}[2]{\expandafter\newcommand\csname CS@#1\endcsname{#2}}

% theorems
\newtheoremstyle{plainx}{7pt plus 2pt}{7pt plus 2pt}{\itshape}{}{\bfseries}{.}{0.5em}{}
\newtheoremstyle{uprightx}{7pt plus 2pt}{7pt plus 2pt}{}{}{\bfseries}{.}{0.5em}{}
\theoremstyle{plainx}
\newtheorem{theorem}{Theorem}[section]
\newtheorem{lemma}[theorem]{Lemma}
\newtheorem{proposition}[theorem]{Proposition}
\newtheorem{corollary}[theorem]{Corollary}
\theoremstyle{uprightx}
\newtheorem{definition}[theorem]{Definition}
\newtheorem{setting}[theorem]{Setting}          % hypotheses shared by several statements; each of them cites it by number
\newtheorem{remark}[theorem]{Remark}
\newtheorem{example}[theorem]{Example}
\renewcommand{\qedsymbol}{{\setlength{\fboxsep}{0pt}\setlength{\fboxrule}{0.5pt}\fbox{\rule{0pt}{0.46em}\hspace{0.46em}}}}
% what holds throughout the note: the ambient objects only (implicit and instance arguments), never a hypothesis
\newenvironment{conventions}{\par\addvspace{5pt}\noindent\textbf{Conventions.}\enspace\ignorespaces}{\par\addvspace{3pt}}
\makeatletter
\renewenvironment{proof}[1][Proof]{\par\pushQED{\qed}\normalfont\topsep0pt\partopsep0pt
  \trivlist\item[\hskip\labelsep\itshape #1.]\ignorespaces}{\popQED\endtrivlist\@endpefalse}
\makeatother

% where each declaration of the concordance is treated: \cs{key}{place}, the place naming a \label or a lettered item
\cs{cruderows}{Theorem 1}
\cs{gridPeriodclimb}{Theorem 2}
\cs{hasRoutesucc}{Theorem 2$'$}
\cs{existsisTidyintegraleqgridPeriod}{Theorem 3}
\cs{existsmulendPeriodmemembValRing}{Theorem 4}
\cs{endPeriodle}{Theorem 4}
\cs{GridIndex}{Definition~\ref{d:grid}}
\cs{gridNumber}{Definition~\ref{d:grid}}
\cs{gridBlockEquiv}{Definition~\ref{d:grid}}
\cs{gridSigns}{Definition~\ref{d:grid}}
\cs{gridParam}{Definition~\ref{d:grid}}
\cs{signsY}{Definition~\ref{d:grid}}
\cs{paramY}{Definition~\ref{d:grid}}
\cs{gridWeight}{Definition~\ref{d:grid}}
\cs{gridHarmonic}{Definition~\ref{d:grid}}
\cs{gridDiagEquiv}{Definition~\ref{d:grid}}
\cs{gridLevelGen}{Definition~\ref{d:grid}}
\cs{gridModulus}{Definition~\ref{d:grid}}
\cs{triangleModulus}{Definition~\ref{d:grid}}
\cs{reduceToTriangle}{Definition~\ref{d:grid}}
\cs{instFactNeRingOfIntegersEHPowGbOfNat}{Definition~\ref{d:grid}}
\cs{instNeZeroNatB}{Definition~\ref{d:grid}}
\cs{entrySigneq}{Definition~\ref{d:grid}}
\cs{cardgridIndex}{Definition~\ref{d:grid}}
\cs{cardgridIndexle}{Definition~\ref{d:grid}}
\cs{cardgridIndexmodtwo}{Definition~\ref{d:grid}}
\cs{cardneggridSigns}{Definition~\ref{d:grid}}
\cs{cardneggridSignseqsum}{Definition~\ref{d:grid}}
\cs{gridSignseqone}{Definition~\ref{d:grid}}
\cs{gridSignszero}{Definition~\ref{d:grid}}
\cs{gridDiagEquivspec}{Definition~\ref{d:grid}}
\cs{waddgaple}{Definition~\ref{d:grid}}
\cs{hasGapgridParam}{Definition~\ref{d:grid}}
\cs{absgridParamle}{Definition~\ref{d:grid}}
\cs{gridParaminterlace}{Definition~\ref{d:grid}}
\cs{nStarbounds}{Definition~\ref{d:grid}}
\cs{rankboundscardbigIndex}{Definition~\ref{d:grid}}
\cs{twomulcarddetIndexaddle}{Definition~\ref{d:grid}}
\cs{isPosFirstOrderorderW}{Definition~\ref{d:grid}}
\cs{paramWltoforderWlt}{Definition~\ref{d:grid}}
\cs{placesmodulusequnion}{Definition~\ref{d:grid}}
\cs{isTameFortwistedChi}{Definition~\ref{d:grid}}
\cs{gridModulusoneideal}{Definition~\ref{d:grid}}
\cs{piGBardvdgridLevelGen}{Definition~\ref{d:grid}}
\cs{twistonespec}{Definition~\ref{d:grid}}
\cs{edgeCharProdmul}{Definition~\ref{d:grid}}
\cs{legDescentscongr}{Definition~\ref{d:grid}}
\cs{gridPathMapcongr}{Definition~\ref{d:grid}}
\cs{cardmatrixletupleBound}{Definition~\ref{d:grid}}
\cs{boxesendWeightDatale}{Definition~\ref{d:grid}}
\cs{paramsCorrespondgrid}{Lemma~\ref{l:corr}}
\cs{paramsCorrespondfirstarrow}{Lemma~\ref{l:corr}}
\cs{pow}{Definition~\ref{d:twists}}
\cs{MFoldTwist}{Definition~\ref{d:twists}}
\cs{powTwist}{Definition~\ref{d:twists}}
\cs{chainChar}{Definition~\ref{d:twists}}
\cs{ratio}{Definition~\ref{d:twists}}
\cs{finitepair}{Definition~\ref{d:twists}}
\cs{instOneTws}{Definition~\ref{d:twists}}
\cs{instCoeTwsTwist}{Definition~\ref{d:twists}}
\cs{nStar}{Definition~\ref{d:wts}}
\cs{orderW}{Definition~\ref{d:wts}}
\cs{conjHolWeight}{Definition~\ref{d:wts}}
\cs{kernelWeight}{Definition~\ref{d:wts}}
\cs{conjHolWeighteqaddkernelWeight}{Lemma~\ref{l:wts}}
\cs{conjHolWeightkernelWeightbounds}{Lemma~\ref{l:wts}}
\cs{fitsParamsauxPaira}{Lemma~\ref{l:wts}}
\cs{heckePolyOnFormstestValueeqzero}{Lemma~\ref{l:wts}}
\cs{isTameForofhasWeight}{Lemma~\ref{l:wts}}
\cs{EdgeTuple}{Definition~\ref{d:tuples}}
\cs{GridDatum}{Definition~\ref{d:tuples}}
\cs{IsMuEdge}{Definition~\ref{d:tuples}}
\cs{edgeChar}{Definition~\ref{d:tuples}}
\cs{IsAdmissibleFamily}{Definition~\ref{d:tuples}}
\cs{tupleAction}{Definition~\ref{d:tuples}}
\cs{edgeCharProd}{Definition~\ref{d:tuples}}
\cs{tupleBound}{Definition~\ref{d:tuples}}
\cs{LeftInvOnMuEdges}{Definition~\ref{d:tuples}}
\cs{IsShortSum}{Definition~\ref{d:tuples}}
\cs{bind}{Definition~\ref{d:tuples}}
\cs{IsTidy}{Definition~\ref{d:tuples}}
\cs{IsLevelTidy}{Definition~\ref{d:tuples}}
\cs{verticalEdge}{Definition~\ref{d:paths}}
\cs{rowEdge}{Definition~\ref{d:paths}}
\cs{legEdges}{Definition~\ref{d:paths}}
\cs{gridPath}{Definition~\ref{d:paths}}
\cs{path}{Definition~\ref{d:paths}}
\cs{legDescents}{Definition~\ref{d:paths}}
\cs{gridPathMap}{Definition~\ref{d:paths}}
\cs{diagonalPeriodForm}{Definition~\ref{d:paths}}
\cs{gridPeriod}{Definition~\ref{d:paths}}
\cs{HasRoute}{Definition~\ref{d:route}}
\cs{binda}{Lemma~\ref{l:bind}}
\cs{existssquaremove}{Proposition~\ref{p:move}}
\cs{existssquareLawgrid}{Lemma~\ref{l:sqgrid}}
\cs{descentpastgridPathMap}{Proposition~\ref{p:past}}
\cs{triangleOrdereqmullevelIndex}{Lemma~\ref{l:order}}
\cs{diagonalPeriodFormformRep}{Lemma~\ref{l:transl}}
\cs{formReplegDescentstupleAction}{Lemma~\ref{l:transl}}
\cs{diagonalPeriodFormtupleAction}{Lemma~\ref{l:transl}}
\cs{existsisLevelTidyiotaPluseq}{Proposition~\ref{p:tidy0}}
\cs{existsisTidydiagonalPeriodFormeq}{Proposition~\ref{p:tidy0}}
\cs{existsrestrictbordered}{Lemma~\ref{l:restrict}}
\cs{iotaPlusadjointweightedChain}{Lemma~\ref{l:adjchain}}
\cs{gridLevelFormseqmapiotaPlus}{Lemma~\ref{l:adjchain}}
\cs{existsisSteep}{Lemma~\ref{l:steep}}
\cs{eqzeroofpeterssonrestricteqzero}{Proposition~\ref{p:vanish}}
\cs{eqzerooforthogonalgridLevelForms}{Corollary~\ref{c:rows}}
\cs{EndIndex}{Definition~\ref{d:end}}
\cs{endWeightData}{Definition~\ref{d:end}}
\cs{endDetExp}{Definition~\ref{d:end}}
\cs{endLevel}{Definition~\ref{d:end}}
\cs{endKernel}{Definition~\ref{d:end}}
\cs{endBranching}{Definition~\ref{d:end}}
\cs{endPairing}{Definition~\ref{d:end}}
\cs{endBlockPoint}{Definition~\ref{d:end}}
\cs{endToAlgForm}{Definition~\ref{d:end}}
\cs{FixedByEndLevels}{Definition~\ref{d:end}}
\cs{endDatum}{Definition~\ref{d:end}}
\cs{endChain}{Definition~\ref{d:end}}
\cs{endPeriod}{Definition~\ref{d:end}}
\cs{isStdgridSpace}{Definition~\ref{d:end}}
\cs{endLevelleprincipalLevel}{Lemma~\ref{l:endeq}}
\cs{gridPeriodeqendPeriod}{Lemma~\ref{l:endeq}}
\cs{endPeriodmem}{Proposition~\ref{p:endmem}}
\newcommand{\sR}{\mathsf{R}}
\newcommand{\bn}{\mathbf{n}}
\newcommand{\br}{\mathbf{r}}
\newcommand{\sS}{\mathsf{S}}
\newcommand{\sV}{\mathsf{V}}
\newcommand{\sZ}{\mathsf{Z}}
\newcommand{\sY}{\mathsf{Y}}
\newcommand{\sW}{\mathsf{W}}
\newcommand{\sI}{\mathsf{I}}
\newcommand{\sg}{\mathsf{g}}
\newcommand{\sh}{\mathsf{h}}
\newcommand{\sk}{\mathsf{k}}
\newcommand{\sx}{\mathsf{x}}
\newcommand{\su}{\mathsf{u}}
\newcommand{\sq}{\mathsf{q}}
\newcommand{\pp}{\mathsf{p}}
\newcommand{\fp}{\mathfrak{p}}
\newcommand{\fm}{\mathfrak{m}}
\newcommand{\fa}{\mathfrak{a}}
\newcommand{\fl}{\mathfrak{l}}
\newcommand{\rc}{\mathrm{c}}
\newcommand{\rC}{\mathrm{C}}
\newcommand{\A}{\mathbb{A}}
\newcommand{\CC}{\mathbb{C}}
\newcommand{\GL}{\operatorname{GL}}
\newcommand{\U}{\operatorname{U}}
\newcommand{\tr}{\operatorname{tr}}
\newcommand{\cL}{\mathcal{L}}
\newcommand{\cE}{\mathcal{E}}
\newcommand{\cO}{\mathcal{O}}
\newcommand{\cG}{\mathcal{G}}
\newcommand{\Np}{\operatorname{N}}
\newcommand{\negc}{\operatorname{neg}}
\newcommand{\pet}[2]{\langle #1,#2\rangle_{\mathrm{Pet}}}
\newcommand{\SqLaw}[5]{\operatorname{SqLaw}(#1,#2;#3;#4,#5)}
\newcommand{\Cost}[3]{\operatorname{Cost}_{#1}(#2;#3)}
\newcommand{\Tlev}[1]{\Theta^{\mathrm{lev}}_{D,#1}}
\newcommand{\Tall}[1]{\Theta^{\circ}_{D,#1}}
\newcommand{\dn}[2]{\operatorname{des}^{#1}_{#2}}
\newcommand{\up}[2]{\operatorname{lift}^{#1}_{#2}}

\begin{document}
\color{ink}

{\raggedright
{\sffamily\bfseries\fontsize{7.2}{9}\selectfont\color{leancol}\addfontfeatures{LetterSpace=8}READBACK NOTE 18\par}
\vspace{5pt}
{\fontsize{21.5}{25.5}\selectfont\bfseries The grid of lifts and descents, and the crude rows\par}
\vspace{5pt}
{\fontsize{12.4}{16}\selectfont\itshape\color{inksoft} One unit of the Lean artifact \ArtifactName, read back into
mathematics\par}}
\vspace{8pt}
{\color{rulegrey}\hrule height 0.6pt}
\vspace{7pt}
\thispagestyle{empty}

\textbf{Summary.} Theorem~C of the handout completes a family of rows to a basis by \emph{crude rows}: rows whose
twisted chains of lifts are integral after multiplication by a constant that costs few squares, and whose
untwisted lifts have bounded conjugates. This note proves that such rows span (Theorem~1). The pairing of a
chain of lifts with a test form is written as a period over a grid of unitary groups; the exchange theorem, in
the form of the square law, moves the period across the grid, each square at the cost of one square of
$L$-values (Theorem~2); at the far end the period is a pairing of two chains of lifts between definite groups,
which is integral and small for elementary reasons (Theorem~4); and at the start the periods see every row
(Theorem~3).

\statement{leancol}{\fontsize{9.6}{13.2}\selectfont\textbf{What this note is.} The artifact \ArtifactName{} is a
body of Lean code: the compact edition of a formal proof, attributed by its authors to L.~Alpöge with the AI system
Claude, whose final theorems are the Riemann hypothesis and its generalisation to Dirichlet $L$-functions. That
headline is a claim. The authors report that every file compiles and that Lean's kernel accepted the proofs; no
mathematician has read the argument through, and nothing here says that the Riemann hypothesis is proved. This
note is the third of three on one unit of the artifact, the unit of the exchange identity, which belongs to the
proof of Theorem~C of the handout. It treats \UnitInNote{} of the unit's \UnitMembers{} declarations; the notes
\emph{Partner data, vanishing on standard data, and the exchange identity at the top of the tower} (the first
note) and \emph{The exchange theorem, one square of lifts, and the triangle of lifts and descents} (the second
note) treat the others. The proofs are written from the structured proofs of the unit's dossier, which a program
compared with the Lean text; what they take from other units and from the two other notes is listed as Inputs in
§\ref{s:setting} and is not proved here. The statements of Theorems 1 to~4 and the definitions of routes, of
grid periods and of the maps along a path were compared with their Lean statements by the writer;
§\ref{s:notes} says what was not. A second reader compared all statements of the unit and a passage of every structured proof with the Lean text after the dossier was made; where her reading differs from the dossier's, this note follows hers (§\ref{s:notes}). The Lean text read is \PinText; Appendix~A is the concordance of \ConcDecls{}
declarations.}

\sect{The statements}{s:intro}
Let $P$ be parameters and $D$ admissible data, with scales $x$, $x_1$, $b$ and number of lifts $m$. The forms of
\emph{principal level} on the first space $\sY_0(D)$ of the chain of lifts are the forms of the type of the rows
fixed by the principal level at the chosen places; they contain the rows of the overview. For an $m$-fold twist
$\vartheta$, $\Tlev{\vartheta}$ is the composite of the $m$ theta lifts from $\sY_0(D)$ to the definite space
$\sY_m(D)$ on these forms, with values in algebraic forms, and $\Tall{1}$ the chain for the trivial twist.
The statement $\Cost{D}{n}{\rc}$ (``$\rc$ costs $n$ squares'') is defined in the second note: for every fit family
$\theta=(\theta_{\ell,\iota})$ of $m$-fold twists, indexed by all pairs of a prime $\ell$ and an isomorphism $\iota$
of $\CC$ with an algebraic closure of the $\ell$-adic numbers, the family
$(\ell,\iota)\mapsto|\rc(\theta_{\ell,\iota})/\rc(1)|^{-2}$ has height at most $n\cdot2\,t(D)\,x^{999}$. For a prime
$\ell$, $\CC_\ell$ is the completion of an algebraic closure of the $\ell$-adic numbers; a ring homomorphism
$j\colon\CC\to\CC_\ell$ is what the formal text calls an $\ell$-adic place, and a number, a vector or a form is
\emph{$j$-integral} if its image under $j$ has absolute value, respectively gauge, at most $1$ (for a form: at
every point).

\statement{arith}{\phantomsection\label{t:1}\textbf{Theorem 1} (crude rows).\enspace \emph{Let $P$ be parameters and
$D$ admissible data for $P$. There is a function $\rc$ from the $m$-fold twists of $D$ to the non-zero complex
numbers such that $\Cost{D}{2m^2}{\rc}$ holds and such that the space of all forms of principal level on
$\sY_0(D)$ is spanned by the forms $y$ with the following two properties:}
\begin{enumerate}[label=\textup{(\arabic*)},leftmargin=2.3em,itemsep=1pt]
\item \emph{for every $m$-fold twist $\vartheta$, every prime $\ell$ prime to $\Np Q(\vartheta)$ and every ring
homomorphism $j\colon\CC\to\CC_\ell$ (every $\ell$-adic place): the form
$(\rc(\vartheta)/\rc(1))\cdot\Tlev{\vartheta}(y)$ is $j$-integral;}
\item \emph{for every automorphism $\sigma$ of $\CC$ and every point $\sg$: the squared Weyl norm
(§\ref{s:setting}) of the vector $\sigma\bigl(\Tall{1}(y)(\sg)\bigr)$ is at most $\exp(x_1^2\,b^2)$.}
\end{enumerate}}
\leanline{\lean{cruderows}}

\statement{arith}{\phantomsection\label{t:2}\textbf{Theorem 2} (the climb).\enspace \emph{Let $P$ be parameters, $D$
admissible data and $i\le m$ a natural number. Then $\HasRoute{D}{0}{i}{2i\cdot i}$ holds: there is a function
$\mathrm{d}$ from the $m$-fold twists to the non-zero complex numbers with $\Cost{D}{2i^2}{\mathrm{d}}$ such that
for every $n$ and every family $\cL$ of data with $\cL_\vartheta$ tidy of order $n$ on the path of stage $0$ for
all $\vartheta$, there is a family $\cL'$ with $\cL'_\vartheta$ tidy of order $n+2i^2$ on the path of stage $i$ for
all $\vartheta$, with this property: for every complex function $T$ of the triples $(\vartheta,F,F')$, if for
every real point $\sh$ the stage-$i$ periods of $\cL'$ at $\sh$ are a constant multiple of $T$, then for every real
point $\sh$ the stage-$0$ periods of $\cL$ at $\sh$ are a constant multiple of $T/\mathrm{d}(\vartheta)$.}

\vspace{4pt}\phantomsection\label{t:2b}\textbf{Theorem 2$'$} (one step).\enspace \emph{Let $P$ be parameters, $D$
admissible data and $i<m$ a natural number. Then $\HasRoute{D}{i}{i+1}{4i+2}$ holds.}}
\leanline{\lean{gridPeriodclimb}\sep \lean{hasRoutesucc}}

\statement{arith}{\phantomsection\label{t:3}\textbf{Theorem 3} (the period at the start sees every row).\enspace
\emph{Let $P$ be parameters and $D$ admissible data. There are a family $\cL$ of data with
$\IsTidy{4m}{\stagePath{0}}{\cL_\vartheta}$ for every $m$-fold twist $\vartheta$, a real point $\sh$, a complex number
$\mathrm{a}$, and a linear map $\mathsf{r}$ from the forms of principal level $\pi_{\mathrm{gr}}$ on the large definite space
$\sZ_{0,m}(D)$ to the forms of principal level on $\sY_0(D)$, surjective and independent of $\vartheta$, such that
for every $m$-fold twist $\vartheta$, every form $F$ on $\sY_m(D)$ and every form $F'$ in the domain of
$\mathsf{r}$}
\[\mathrm{i}_D(\vartheta)\int_{[\U(\sY_m(D))]}F(\sx)\cdot\operatorname{ch}_{\vartheta}\bigl(\mathsf{r}(F')\bigr)(\sx)\,d\sx
=\mathrm{a}\cdot\gridPeriod{\cL}{(\vartheta,F,F')}{0}{\sh},\]
\emph{where $\operatorname{ch}_{\vartheta}$ is the chain of the $m$ lifts for $\vartheta$, with values in functions on
$\U(\sY_m(D))(\A)$, and $\mathrm{i}_D(\vartheta)$ is the index of the last level of the chain: the index, in the
stabilizer of the integral lattice of $\sY_m(D)$, of the principal level of the modulus $\fm_D(\vartheta,1)$ (the
product of the chosen places and of the places of $Q(\vartheta)$), that is,
$\mathrm{i}_D(\vartheta)=\#\GL_{\GridIndex{D}{\mathrm{false}}{0}{m}}\bigl(\cO_E/\fm_D(\vartheta,1)\bigr)=\mathrm{ord}_{D,\mathrm{false}}(\vartheta,m)$.
It is not the number $\mathrm{i}^{\mathrm{gr}}_D(\vartheta)$ that enters the period (Definition~\ref{d:paths}).}}
\leanline{\lean{existsisTidyintegraleqgridPeriod}}

\statement{arith}{\phantomsection\label{t:4}\textbf{Theorem 4} (the period at the end is integral and
small).\enspace \emph{Let $P$ be parameters and $D$ admissible data.}

\emph{\textup{(a)} There is a natural number $\mathrm{d}_{\mathrm{end}}\ne0$ with $\mathrm{d}_{\mathrm{end}}\le\exp(x^{99}b^2)$ such that for
every $m$-fold twist $\vartheta$, every $n$, every datum $\cL_0$ with $\IsTidy{n}{\stagePath{m}}{\cL_0}$, every pair
$y$ of forms with $\FixedByEndLevels{\vartheta}{y}$, and every prime $\ell$ with a homomorphism $j$ from $\CC$ to
the completed algebraic closure of the $\ell$-adic numbers such that $\ell$ is prime to $\Np Q(\vartheta)$: if the
values of the algebraic forms of the two members of $y$ are $j$-integral, then
$\mathrm{d}_{\mathrm{end}}\cdot\finPeriod{\cL_0}{y}$ is $j$-integral.}

\emph{\textup{(b)} Let $\cL_0$ be a datum for the trivial twist with $\IsTidy{4m+2m\cdot m}{\stagePath{m}}{\cL_0}$, let
$y$ be a pair of forms with $\FixedByEndLevels{1}{y}$, $\sigma$ an automorphism of $\CC$, and $B_0$ a real
function of the two grids such that the squared Weyl norms of the $\sigma$-conjugates of the values of the
algebraic form of the member of $y$ in the grid $o$ are at most $B_0(o)$. Then}
\[\bigl|\sigma\bigl(\finPeriod{\cL_0}{y}\bigr)\bigr|\le\exp(x^{200}b^2)\cdot\sqrt{B_0(\mathrm{false})}\cdot\sqrt{B_0(\mathrm{true})}.\]}
\leanline{\lean{existsmulendPeriodmemembValRing}\sep \lean{endPeriodle}}

\begin{center}
\begin{tikzpicture}[x=9mm,y=7mm,>={Stealth[length=1.5mm]},
  nd/.style={circle,fill=inksoft,inner sep=1pt},
  lb/.style={font=\fontsize{7.2}{8.8}\selectfont\color{inksoft}}]
\foreach \i in {0,...,3} \foreach \n in {0,...,3} \node[nd] (p\i\n) at (\i,\n) {};
\foreach \n/\nn in {0/1,1/2,2/3} \draw[->,rulegrey] (p0\n) -- (p0\nn);
\draw[->,arith,line width=0.7pt] (p03) -- (p02); \draw[->,arith,line width=0.7pt] (p02) -- (p01);
\draw[->,arith,line width=0.7pt] (p01) -- (p11);
\draw[->,ana,line width=0.7pt] (p03) to[bend left=18] (p02);
\draw[->,ana,line width=0.7pt] (p02) -- (p12); \draw[->,ana,line width=0.7pt] (p12) -- (p22);
\node[lb,left] at (0,3) {$(0,m)$}; \node[lb,left] at (0,0) {$(0,0)$}; \node[lb,below] at (3,0) {$(m,0)$};
\node[lb,right] at (3,3) {$(m,m)$};
\node[lb,right] at (4.2,2.3) {\textcolor{arith}{stage 1}: descend to $(0,1)$, lift to $(1,1)$};
\node[lb,right] at (4.2,1.5) {\textcolor{ana}{stage 2}: descend to $(0,2)$, lift to $(2,2)$};
\node[lb,right] at (4.2,0.7) {first index $i$ to the right, second index $n$ upwards};
\end{tikzpicture}
\end{center}
\vspace{-4pt}
{\fontsize{8.9}{12}\selectfont\color{inksoft}\noindent One of the two grids, for $m=3$. A form at the node $(0,m)$
is carried along the path of stage $i$ to the diagonal node $(i,i)$; the period pairs the results in the two grids.
At stage $0$ the path consists of $m$ descents, at stage $m$ of $m$ lifts along the top row.\par}

\sect{The setting}{s:setting}
\begin{conventions}\phantomsection\label{c:conventions}
The natural numbers are $0,1,2,\dots$. The letters of the handout keep their meaning: $P$ are parameters, $D$
admissible data, with scales $x$, $x_1$, $b$, number of lifts $m$, rank of the rows $r$, step $t(D)$; $\vartheta$
is a tame character or a twist of $D$, $\tau$ an infinite place, $\sigma$ an automorphism of $\CC$, $\ell$ a prime,
$\iota$ an isomorphism of $\CC$ with an algebraic closure of the $\ell$-adic numbers and $j$ a ring homomorphism from
$\CC$ to the completion of that closure. As in the other notes of the unit, finite index sets are bold,
sans-serif letters are rings, hermitian spaces, matrices, tuples of matrices, points of groups and linear maps,
fraktur letters are places, moduli and records, and upright letters are constants and functions with constant
values. A \kw{flag} $o$ is one of the two values false and true; it names a grid, and $[o]$ is $0$ for false and
$1$ for true. The implicit arguments of the declarations of the unit are of the following kinds, and each
statement introduces the ones it uses: parameters and admissible data with their tame characters, twists and
pairs; flags and natural numbers; nodes and finite sets of edges; tuples, data and families of matrices at the
nodes; forms and linear functionals; finite index sets and bijections between them; characters; a CM field with a
CM type and Weil data, hermitian spaces, signs and parameters; a prime with an embedding of $\CC$; real numbers
and constants.
\end{conventions}

\subhead{What is recalled.} From the second note: the square law $\SqLaw{\chi}{\mathsf{e}}{\rc}{\cG}{\cG'}$ for
functions $\cG$, $\cG'$ on matrices over a finite ring, with its partner identity; the graph kernel and the
block embedding of general linear groups; costs in squares, their product rule, and the cost of a constant whose
squared absolute value is a product of $L$-values; one square of lifts and descents; triangles, weighted chains
and the triangle of the data, with its orders $\mathrm{ord}_{D,o}(\vartheta,n)$; $L$-value data $\fl$ and their
auxiliary pairs; pairs $\pp$ of the data with character $\omega_{\pp}$ and point $e_{\pp}$; standard data and
admissible combinations. From other units: the index set of the rows, with $r$ elements, numbered; the weights
$w_{D,i}(\tau)$ of the lifts and the parameter of the rows; the chosen places and the integer $\bar\pi_g$ of the
data, with a non-zero exponent $b_1(D)$; theta lifts $\up{}{}$ and descents $\dn{}{}$ between the forms at
two neighbouring nodes of a grid, attached to a matrix modulo the modulus of the grid (unit 141); the level of a
node and its index; the untwisted and twisted descents from the node $(0,m)$ to the node $(0,n)$ along a grid
(unit 141); algebraic forms on definite spaces, their gauges and norms; theta kernels between definite spaces
and chains of them (unit 134). The values of an algebraic form are coordinate vectors $\xi$ of the representation
of its weight in the basis of products of minors, and the \kw{squared Weyl norm} of such a vector is
$\xi^{*}\mathrm{G}\,\xi$, where $\mathrm{G}$ is the Gram matrix of the representation in this basis, a positive
definite matrix with integer entries (background library). It is not the sum of the squares of the coordinates.
The moduli $\fm_D(\vartheta,e)$ are those of the second note: the product of the chosen places to the power $e$
times the places of $Q(\vartheta)$.

\statement{inksoft}{\textbf{Inputs.} The proofs below take the following for granted. Numbers are the units of the
reading; references are those of the handout's catalogue, with the relation to the literature, and the place
where one is given, that the second reader found by comparing the formal statement with the source.

\vspace{3pt}\textbf{Input 1} (the second note of this unit).\enspace One square (its Theorem~2); the triangle of
the data (its Theorem~3) and the injectivity of weighted chains; the square law at the inclusion and after a
Fourier transform; the graph kernel reproduces the small side of a square law; costs: the product rule, the cost
of a product of $L$-values of a pair ramified at the chosen places; tameness needs only the weights; the
standard data of a function modulo a modulus are admissible combinations of translates of the integral vacuum.

\vspace{3pt}\textbf{Input 2} (unit 141, Seesaw; units 27 and 25 of the background library).\enspace
\textup{(a)} The classical \emph{seesaw identity}: for a see-saw pair of unitary groups, the period of a lift
against a form is the period of the form against the lift in the other direction (S.~Kudla, Seesaw dual
reductive pairs, Progr.\ Math.\ 46 (1984), 244--268, §1: the adjointness (1.9) and the seesaw identity (1.14)).
\textup{(b)} The \emph{seesaw identity for periods} of the formal text, which is what the proofs use: a
proportionality of finite sums of weighted periods is transferred from lifts to descents, under a spanning
hypothesis. It is an inference from \textup{(a)} with additional hypotheses and not a statement of the source (a
weaker form). \textup{(c)} The lifts and descents along the grids and their transformation under the levels;
adjointness of a lift and a descent through an edge family.

\vspace{3pt}\textbf{Input 3} (units 134, 133, 116, 102 and 143, forms).\enspace A dual pairing of the weights of two
neighbouring definite groups and the branching pairing; delta forms; short algebraic forms span the forms of a
level (unit 143); bounds for forms near a cusp and in terms of the first minimum (reduction theory: I.~Satake and R.~Godement, Sém.\ H.~Cartan 10 (1957/58); the formal statement is specific to this proof); the Laplace method for matrix coefficients of discrete series with steep parameters (reference to
be confirmed); Poincaré series as kernels of periods (relation to the literature unresolved: R.~Godement, Sém.\
H.~Cartan 10 (1957/58), exp.~10, §4, Thm~5 and 5~bis, and G.~Muić, arXiv:1610.05483, Thm~6-4 and Thm~7-6,
construct such series under additional hypotheses, and the formal statement allows cases that they exclude).

\vspace{3pt}\textbf{Input 4} (units 150 Curve and 145 Columns).\enspace Jets of admissible combinations of
translates of the integral vacuum: their integrality at a prime and the bounds for the norms of their conjugates
under all automorphisms of $\CC$, uniformly along a chain of kernels.

\vspace{3pt}\textbf{Input 5} (units 96 and 104).\enspace Hecke polynomials and tame characters for a modulus; the
places of the moduli of the data.}

\sect{The two grids}{s:grid}
\begin{definition}[index sets, signs, parameters and moduli of the grids]\label{d:grid}
Let $P$ be parameters and $D$ admissible data.
\begin{enumerate}[label=\textup{(\alph*)},leftmargin=2.3em,itemsep=1pt]
\item For a flag $o$ and natural numbers $i$, $n$, the index set $\GridIndex{D}{o}{i}{n}$ of the node $(i,n)$ of
the grid $o$ consists of $r+[o]$ base indices (the index set of the rows and, if $o$ is true, one more point),
of $i$ blocks of size $t(D)+[o]$ and of $n$ blocks of size $t(D)-[o]$. The node carries a hermitian space with
this index set: $\sV_{i,n}(D)$ of rank $r+(i+n)\,t(D)$ in the grid false, and a space $\sZ_{i,n}(D)$ of rank
$r+1+i(t(D)+1)+n(t(D)-1)$ in the grid true. $\gridNumber{D}{o}$ is the numbering of the base indices (the given
numbering of the rows, and $r$ for the extra point), and $\gridBlockEquiv{D}$ the bijection that lists the indices
of a node as the base indices followed by the blocks.
\item For an offset $n'$ and an index $\mathsf{x}$ of $\GridIndex{D}{o}{i}{n}$, the sign
$\gridSigns{D}{o}{n'}{i}{n}{\mathsf{x}}$ is $+1$ on the added blocks and, on the base index number $i'$, the sign
$-1$ or $+1$ according to the parity of an explicit function of $n'$, $n$ and $i'+1-k(D)$.
\item For an offset $n'$ and an infinite place $\tau$, $\gridParam{D}{o}{n'}{i}{n}{\tau}$ is the doubled
Harish-Chandra parameter at the node $(i,n)$ for these signs: on the base indices twice the parameter of the rows,
and on the added blocks arithmetic progressions of step $2$ whose position is given by the weights
$w_{D,i'}(\tau)$ of the lifts, shifted by $n'$ for the blocks of the second kind and by the parity of $i+n$.
\item The signs and the doubled parameter of the space $\sY_n(D)=\sV_{0,n}(D)$ of the chain are
\begin{align*}
\signsY{D}{n}{\mathsf{x}}&=\gridSigns{D}{\mathrm{false}}{0}{0}{n}{\mathsf{x}},\\
\paramY{D}{n}{\tau}{\mathsf{x}}&=\gridParam{D}{\mathrm{false}}{0}{0}{n}{\tau}(\mathsf{x}).
\end{align*}
\item For a flag $o'$, a natural number $i'$, a finite index set $\bn$ and $\tau$:
$\gridWeight{D}{o'}{i'}{\bn}{\tau}=\lfloor(|\bn|+[\lnot o'])/2\rfloor-w_{D,i'+1}(\tau)$, the weight at $\tau$ of the
character of the grid. The \kw{harmonic of the grid} for two nodes of a grid, an offset, two flags and an index is
the product over the infinite places of the chosen joint harmonics for these weights, the parameters of the two
nodes (negated or not) and their signs.
\item $\gridDiagEquiv{D}{i}$ is a bijection from the index set of $\sV_{i,i}(D)$ plus one point to the index set
of the diagonal node $(i,i)$ of the grid true that keeps the signs, the extra point having the sign
$(-1)^{\min(i,m)}$; it is the identity for $i=0$.
\item The \kw{generator of the grid} is $\pi_{\mathrm{gr}}=\bar\pi_g^{\,b_1(D)}$. For an $m$-fold twist $\vartheta$,
$\gridModulus{\vartheta}$ is the modulus that is the product of the chosen places to the power $b_1(D)$ times
$Q(\vartheta)$; $\triangleModulus{o}{\vartheta}$ is the product of the chosen places times $Q(\vartheta)$ if $o$
is false, and the product of the chosen places to the power $b_1(D)$ if $o$ is true; and
$\reduceToTriangle{\vartheta}{o}$ is the reduction from the residue ring of $\gridModulus{\vartheta}$ to that of
$\triangleModulus{o}{\vartheta}$.
\end{enumerate}
\end{definition}
\leanline{\lean{GridIndex}\sep \lean{gridNumber}\sep \lean{gridBlockEquiv}\sep \lean{gridSigns}\sep
\lean{gridParam}\sep \lean{signsY}\sep \lean{paramY}\sep \lean{gridWeight}\sep \lean{gridHarmonic}\sep
\lean{gridDiagEquiv}\sep \lean{gridLevelGen}\sep \lean{gridModulus}\sep \lean{triangleModulus}\sep
\lean{reduceToTriangle}\sep two instances record that the powers of $\bar\pi_g$ and $b_1(D)$ are non-zero:
\lean{instFactNeRingOfIntegersEHPowGbOfNat}, \lean{instNeZeroNatB}}

\subhead{Small facts on the grids.} They are single steps or computations in the formal text, and are used
inside the proofs below. The sign function of the base indices has a closed form; the index set of a node has
$r+[o]+i(t(D)+[o])+n(t(D)-[o])$ elements, at most twice the rank of $\sY_m(D)$ when $i,n\le m$, and its parity is
even in the grid true and that of $i+n+1$ in the grid false; with offset $0$ the node $(i,n)$ has
$\lfloor(m-n+[o])/2\rfloor$ negative signs (a count written first as a sum), none if $m\le n$; the signs of
the node $(0,0)$ of the grid true are those of $\sY_0(D)$ followed by $+1$, and those of the diagonal nodes
correspond through $\gridDiagEquiv{D}{i}$; the weights of the lifts increase by at least the gap of the data from
one index to the next, and by twice the gap further on; the parameter of a node has that gap, provided the offset
is $0$ or $i=0$; its entries are at most $\lceil20b\rceil$ in absolute value for $i,n\le m$; and at the node
$(m,m)$ the negated parameter of the grid true interlaces the parameter of the grid false. For the weights of
$\sW_i(D)$: three size comparisons of the data; the order of Definition~\ref{d:wts} puts the positive vectors
first and is compatible with the parameter. For the moduli: the places of the moduli of the data are the chosen
places and those of $Q(\vartheta)$; for a lift index $i$ ($1\le i\le m$), the $i$-th power of the basic character times a twist is tame for them; the modulus of the
grid for the trivial twist is generated by $\pi_{\mathrm{gr}}$, which is divisible by $\bar\pi_g$; the trivial twist has the
trivial character and an empty set of places. For tuples: the product of edge characters is multiplicative in the
family; the maps of Definition~\ref{d:paths} depend only on the matrices on their edges; the number of matrices
on an edge between nodes with indices at most $m$ is at most the tuple bound; the end nodes have at most twice
the rank of $\sY_m(D)$ indices, and their weights and exponents are bounded in terms of $b_1(D)\,r^4$ and
$\lceil20b\rceil$.
\leanline{\lean{entrySigneq}\sep \lean{cardgridIndex}\sep \lean{cardgridIndexle}\sep \lean{cardgridIndexmodtwo}\sep
\lean{cardneggridSigns}\sep \lean{cardneggridSignseqsum}\sep \lean{gridSignseqone}\sep \lean{gridSignszero}\sep
\lean{gridDiagEquivspec}\sep \lean{waddgaple}\sep \lean{hasGapgridParam}\sep \lean{absgridParamle}\sep
\lean{gridParaminterlace}\sep \lean{nStarbounds}\sep \lean{rankboundscardbigIndex}\sep
\lean{twomulcarddetIndexaddle}\sep \lean{isPosFirstOrderorderW}\sep \lean{paramWltoforderWlt}\sep
\lean{placesmodulusequnion}\sep \lean{isTameFortwistedChi}\sep \lean{gridModulusoneideal}\sep
\lean{piGBardvdgridLevelGen}\sep \lean{twistonespec}\sep \lean{edgeCharProdmul}\sep \lean{legDescentscongr}\sep
\lean{gridPathMapcongr}\sep \lean{cardmatrixletupleBound}\sep \lean{boxesendWeightDatale}}

\begin{lemma}[the steps of the grids are edges]\label{l:corr}
Let $P$ be parameters and $D$ admissible data.
\begin{enumerate}[label=\textup{(\alph*)},leftmargin=2.3em,itemsep=1pt]
\item Let $\mathrm{d}$ be a flag (a direction), $o$ a flag, $n'$, $i$, $n$ natural numbers, and let
$(o',i',i_1,n_1)$ be $(o,i,i+1,n)$ if $\mathrm{d}$ is true and $(\lnot o,n+n',i,n+1)$ if it is false. Suppose
$i_1\le m$, $n_1\le m$, and $n'=0$ or $i_1=0$. Then the parameters of the nodes $(i,n)$ and $(i_1,n_1)$ of the grid
$o$ with offset $n'$, both negated if $\mathrm{d}$ is true, correspond for the weights
$\gridWeight{D}{o'}{i'}{\GridIndex{D}{o}{i}{n}}{\tau}$ and
$\gridWeight{D}{o'}{i'}{\GridIndex{D}{o}{i_1}{n_1}}{\tau}$ and the signs of the two nodes: the step is an edge.
\item Let $o$ be a flag and $i$, $n$ natural numbers with $i+n<m$. For the complementary weights
$|\GridIndex{D}{o}{0}{n}|-\gridWeight{D}{\lnot o}{i}{\GridIndex{D}{o}{0}{n}}{\tau}$, the negated parameters of the
nodes $(0,n)$ with offset $i+1$ and $(0,n+1)$ with offset $i$ correspond, and so do those of $(0,n+1)$ with
offset $i$ and $(0,n)$ with offset $i+1$.
\end{enumerate}
\end{lemma}
\begin{proof}
(a) At each infinite place the two parameters have the gap of the data (a small fact above), and the entries of
the two nodes are matched one by one by integer arithmetic. (b) follows from (a) with the offsets: the gaps carry
over to the negated parameters, and the matching of the entries is checked again.
\end{proof}
\leanline{\lean{paramsCorrespondgrid}\sep \lean{paramsCorrespondfirstarrow}}

\begin{definition}[twists, their powers, and ratios of $L$-values]\label{d:twists}
Let $P$ be parameters and $D$ admissible data.
\begin{enumerate}[label=\textup{(\alph*)},leftmargin=2.3em,itemsep=1pt]
\item For a tame character $\vartheta$ and a natural number $n$, $\pow{\vartheta}{n}$ is the tame character with
the same set $Q(\vartheta)$ and the $n$-th power of the tame part.
\item $\MFoldTwist{D}$ is the set of the \kw{$m$-fold twists}: the twists $\vartheta$ whose local tame parts
satisfy $(\vartheta^{\mathrm{t}}_{\fp})^{n}\ne1$ for every $\fp\in Q(\vartheta)$ and every $1\le n\le m$. The trivial twist $1$
is one, and an $m$-fold twist is read as a twist. For a lift index $n$, $\powTwist{\vartheta}{n}$ is the power
$\pow{\vartheta}{n}$ as a twist.
\item For a tame character $\vartheta$ and a natural number $i$, $\chainChar{D}{\vartheta}{i}$ is the character of
$\U(1)(\A)$ of the lift number $i$ of a chain, the basic character of that lift times $\pow{\vartheta}{i}$.
\item For a pair $\pp$ of $D$, a tame character $\vartheta$ and a sign $\pm$,
\[\pairRatio{\pp}{\vartheta}{\pm}=\frac{\bigl|L^{\Sigma\cup\bar\Sigma}\bigl(e_{\pp},((\omega_{\pp}\vartheta^{\pm1})^{-1})^{\mathrm{bc}}\bigr)\bigr|}{\bigl|L^{\Sigma_D\cup\bar\Sigma_D}\bigl(e_{\pp},(\omega_{\pp}^{-1})^{\mathrm{bc}}\bigr)\bigr|},
\qquad\Sigma=\Sigma_D\cup Q(\vartheta),\]
where, as in the two other notes, $L^{\Sigma\cup\bar\Sigma}(s,\omega^{\mathrm{bc}})$ is the Hecke $L$-function of the
base change $\omega^{\mathrm{bc}}$ of a character $\omega$ to $E(D)$ without the Euler factors at a set $\Sigma$ of
places and at the set $\bar\Sigma$ of their conjugates, and $\Sigma_D$ is the set of the chosen places: the Euler
factors at the chosen places and at their conjugates are removed in both values, and those at $Q(\vartheta)$ and
its conjugates in the numerator.
\end{enumerate}
\end{definition}
\leanline{\lean{pow}\sep \lean{MFoldTwist}\sep \lean{powTwist}\sep \lean{chainChar}\sep \lean{ratio}\sep the set
of pairs is finite, the trivial twist and the reading of an $m$-fold twist as a twist are instances:
\lean{finitepair}, \lean{instOneTws}, \lean{instCoeTwsTwist}}

\begin{definition}[weights on the spaces $\sW_i(D)$]\label{d:wts}
Let $P$ be parameters and $D$ admissible data. $\nStar{D}=f(D)\cdot2\,t(D)^3$. For a lift index $i$ and an index
$\mathsf{x}$ of the basis of $\sW_i(D)$, $\orderW{i}{\mathsf{x}}$ is an integer that orders the basis, the positive
vectors first and in reversed order. For an infinite place $\tau$, $\conjHolWeight{i}{\tau}$ is the lowest
$K$-type of the Harish-Chandra parameter that is minus the parameter of $\sW_i(D)$ at $\tau$, and
$\kernelWeight{i}{\tau}$ is the staircase weight of the kernel of the $i$-th untwisted auxiliary pair on its
small definite space.
\end{definition}
\leanline{\lean{nStar}\sep \lean{orderW}\sep \lean{conjHolWeight}\sep \lean{kernelWeight}}

\begin{lemma}[weights and auxiliary pairs]\label{l:wts}
\textup{(a)} Let $P$ be parameters, $D$ admissible data, $i$ a lift index and $\tau$ an infinite place. Then
$\conjHolWeight{i}{\tau}$ is the sum of the staircase weight of the big definite space of the $i$-th auxiliary
pair and of $\kernelWeight{i}{\tau}$.

\textup{(b)} With the same $i$ and $\tau$: the entries of $\conjHolWeight{i}{\tau}$ are at most $5b$ in absolute
value; $\conjHolWeight{i}{\tau}$ and $\kernelWeight{i}{\tau}$ are dominant for the signs of $\sW_i(D)$ and the order
$\orderW{i}{\cdot}$; and both have holomorphic gaps $300\,\nStar{D}$.

\textup{(c)} Let $E$ be a CM field with a CM type, $\bn$ a finite index set and $\fl$ an $L$-value datum for $E$
and $\bn$. Then the parameter $0$ fits the auxiliary pair of $\fl$.

\textup{(d)} Let $E$ be a CM field with a CM type and Weil data, $\fa$ an auxiliary pair, and $\fp$ a finite place
that is not its own conjugate and is off the places of $\fa$ and their conjugates. If an algebraic form $y$ of
$\fa$ lifts a family of eigenvalues, then the Hecke polynomial operator of $\fa$ at $\fp$, evaluated at the test
value of $\fa$ at $\fp$, kills $y$.

\textup{(e)} Let $P$ be parameters, $D$ admissible data, $e\ne0$ a natural number, $\vartheta$ a twist and
$\omega$ a character of $\U(1)(\A)$ that has weights $\mu$ and tame part $\omega^{\mathrm{t}}\vartheta^{\mathrm{t}}$, where $\omega^{\mathrm{t}}$ is
non-trivial exactly at the chosen places. Then $\omega$ is tame for the modulus that is the product of the chosen
places to the power $e$ times $Q(\vartheta)$, with the weights $\mu$ and the character of this tame part.
\end{lemma}
\begin{proof}
(a) Entry by entry: by the joint harmonic of the fitting auxiliary pair (second note) the lowest $K$-type is the
staircase weight of the pair; split by the sign of the index. (b) The size from the bound for lowest $K$-types of
a parameter with a gap; dominance and gaps by the sign of each index and sums of column lengths, with the small
facts. (c) The four conditions of fitting are computed for the one-row weight of $\fl$. (d) The Hecke polynomial
acts on $y$ by its value at the test value, where the first factor vanishes because the rank difference is
positive. (e) The four conditions of tameness: the modulus has places, the weights are given, the tame part is
non-trivial exactly at its places, and the character is that of the tame part.
\end{proof}
\leanline{\lean{conjHolWeighteqaddkernelWeight}\sep \lean{conjHolWeightkernelWeightbounds}\sep
\lean{fitsParamsauxPaira}\sep \lean{heckePolyOnFormstestValueeqzero}\sep \lean{isTameForofhasWeight}}

\sect{Tuples, tidy data and periods}{s:tidy}
A \kw{node} is a triple $(o,i,n)$ of a flag and two natural numbers; it stands for the node $(i,n)$ of the grid
$o$, with its index set and its space of forms. An \kw{edge} is an ordered pair of nodes.

\begin{definition}[tuples and data]\label{d:tuples}
Let $P$ be parameters, $D$ admissible data and $\vartheta$ an $m$-fold twist; let $\sR_\vartheta$ be the residue
ring of $\gridModulus{\vartheta}$.
\begin{enumerate}[label=\textup{(\alph*)},leftmargin=2.3em,itemsep=1pt]
\item $\EdgeTuple{\vartheta}$ is the set of \kw{tuples}: a tuple $\sS$ assigns to every edge $(\pp,\pp')$ a
matrix over $\sR_\vartheta$ with columns indexed by the index set of $\pp$ and rows by that of $\pp'$.
$\GridDatum{\vartheta}$ is the space of \kw{data}: the linear functionals on the complex functions of tuples.
\item For nodes $\pp$, $\sq$: $\IsMuEdge{\pp}{\sq}$ holds if $\pp$ lies in the grid true exactly when the first
indices of $\pp$ and $\sq$ differ; these are the edges whose character carries the twist. For an edge
$\mathsf{b}=(\pp,\pp')$, $\edgeChar{\vartheta}{\mathsf{b}}$ is the multiplicative character of $\sR_\vartheta$ made
from the tame part of the character of the lift along $\mathsf{b}$, twisted if $\mathsf{b}$ carries the twist.
\item For a family $\sg$ of invertible matrices over $\sR_\vartheta$, one at each node: $\IsAdmissibleFamily{\sg}$
holds if for both grids the value of $\sg$ at the node $(0,m)$ lies in the level of that node, and for every
$i$ the value at the diagonal node $(i,i)$ of the grid true is the block extension by $1$, through
$\gridDiagEquiv{D}{i}$, of the value at $(i,i)$ of the grid false. For a tuple $\sS$ and an edge
$\mathsf{b}=(\pp,\pp')$, $\tupleAction{\sg}{\sS}{\mathsf{b}}=\sg(\pp')^{-1}\,\sS(\mathsf{b})\,\sg(\pp)$; and for a
finite set $\cE$ of edges,
$\edgeCharProd{\sg}{\cE}=\prod_{\mathsf{b}=(\pp,\pp')\in\cE}\edgeChar{\vartheta}{\mathsf{b}}(\det\sg(\pp')/\det\sg(\pp))$.
\item $\tupleBound{\vartheta}=(\#\sR_\vartheta)^{(2\operatorname{rk}\sY_m(D))^2}$. For a tuple $\sS$ and a finite
set $\cE$ of edges, $\LeftInvOnMuEdges{\sS}{\cE}$ holds if $\sS(\mathsf{b})$ has a left inverse for every edge
$\mathsf{b}\in\cE$ that carries the twist.
\item For a natural number $n$, a finite set $\cE$ of edges and a datum $\cL_0$: $\IsShortSum{n}{\cE}{\cL_0}$
holds if $\cL_0(G)=\sum_{\mathsf{x}\in\mathcal{W}}\rc_{\mathsf{x}}\,G(\sS_{\mathsf{x}})$ for all functions $G$ of
tuples, with a finite set $\mathcal{W}$ of at most $\tupleBound{\vartheta}^{n}$ elements, tuples
$\sS_{\mathsf{x}}$ and coefficients $\rc_{\mathsf{x}}$, each of which is $0$ or a root of unity with
$\LeftInvOnMuEdges{\sS_{\mathsf{x}}}{\cE}$.
\item For a datum $\cL_0$, a finite set $\mathcal{X}$, weights $\mathrm{k}(\sS,\mathsf{x})\in\CC$ and tuples
$\mathrm{t}(\sS,\mathsf{x})$: $\bind{\cL_0}{\mathrm{k}}{\mathrm{t}}$ is the datum
$G\mapsto\cL_0\bigl(\sS\mapsto\sum_{\mathsf{x}\in\mathcal{X}}\mathrm{k}(\sS,\mathsf{x})\,G(\mathrm{t}(\sS,\mathsf{x}))\bigr)$.
\item $\IsTidy{n}{\cE}{\cL_0}$ (``$\cL_0$ is tidy of order $n$ on $\cE$'') holds if $\IsShortSum{n}{\cE}{\cL_0}$
holds and
$\cL_0\bigl(\sS\mapsto G(\tupleAction{\sg}{\sS}{\cdot})\bigr)=\edgeCharProd{\sg}{\cE}\cdot\cL_0(G)$ for every
family $\sg$ with $\IsAdmissibleFamily{\sg}$ and every $G$. $\IsLevelTidy{n}{\cE}{\cL_0}$ holds if
$\IsShortSum{n}{\cE}{\cL_0}$ holds and the same identity holds for every family $\sg$ whose values at the nodes
$(0,i)$, $i\le m$, of both grids lie in the levels of these nodes.
\end{enumerate}
\end{definition}
\leanline{\lean{EdgeTuple}\sep \lean{GridDatum}\sep \lean{IsMuEdge}\sep \lean{edgeChar}\sep
\lean{IsAdmissibleFamily}\sep \lean{tupleAction}\sep \lean{edgeCharProd}\sep \lean{tupleBound}\sep
\lean{LeftInvOnMuEdges}\sep \lean{IsShortSum}\sep \lean{bind}\sep \lean{IsTidy}\sep \lean{IsLevelTidy}}

\begin{definition}[paths, the maps along them, and periods]\label{d:paths}
Let $P$ be parameters, $D$ admissible data and $\vartheta$ an $m$-fold twist.
\begin{enumerate}[label=\textup{(\alph*)},leftmargin=2.3em,itemsep=1pt]
\item For a flag $o$ and natural numbers $i$, $n$: $\verticalEdge{o}{i}{n}$ is the edge from $(o,i,n)$ to
$(o,i,n+1)$ and $\rowEdge{o}{i}{n}$ the edge from $(o,i,n)$ to $(o,i+1,n)$. $\legEdges{o}{n}$ is the set of the
edges $\verticalEdge{o}{0}{n'}$ with $n\le n'<m$; $\gridPath{o}{i}{n}$ is the union of $\legEdges{o}{n}$ and of the
edges $\rowEdge{o}{i'}{n}$ with $i'<i$; and
$\stagePath{n}=\gridPath{\mathrm{false}}{n}{n}\cup\gridPath{\mathrm{true}}{n}{n}$ is the \kw{path of stage $n$}.
\item For a tuple $\sS$, a flag $o$, a natural number $n$ and a form $F$ at the node $(0,m)$ of the grid $o$:
$\legDescents{\sS}{o}{n}{F}$ is the form at $(0,n)$ obtained from $F$ by the descents along the edges
$\verticalEdge{o}{0}{n'}$, $n'=m-1,\dots,n$, each with the matrix that $\sS$ puts on its edge (it is $F$ for $n=m$
and $0$ for $n>m$).
\item For natural numbers $i$, $n$: $\gridPathMap{\sS}{o}{i}{n}{F}$ is the form at $(i,n)$ obtained from
$\legDescents{\sS}{o}{n}{F}$ by the lifts along the edges $\rowEdge{o}{i'}{n}$, $i'=0,\dots,i-1$, with the matrices
of $\sS$. Both maps are linear in $F$.
\item For an $m$-fold twist $\vartheta_1$, a natural number $i\le m$ and a real point $\sh$,
$\diagonalPeriodForm{\vartheta_1}{i}{\sh}$ is the bilinear form on the forms at the diagonal nodes $(i,i)$ of the
two grids
\[\diagonalPeriodForm{\vartheta_1}{i}{\sh}(y,y')=\int_{[\U(\sV_{i,i}(D))]}y(\sx)\,y'\bigl((\sx\oplus\su^{\,i})\,\sh\bigr)\,d\sx,\]
where $\sx\oplus\su^{\,i}$ is read in the group of the grid true through $\gridDiagEquiv{D}{i}$ and $\su$ is the
scalar of norm one attached to $\gridModulus{\vartheta_1}$. In the formal definition the second argument is the
proof of $i\le m$.
\item For a family $\cL$ of data ($\cL_\vartheta\in\GridDatum{\vartheta}$ for every $m$-fold twist), a triple
$(\vartheta,F,F')$ of an $m$-fold twist and forms at the nodes $(0,m)$ of the grids false and true, a natural
number $i\le m$ and a real point $\sh$, the \kw{period of stage $i$} is
\begin{multline*}
\gridPeriod{\cL}{(\vartheta,F,F')}{i}{\sh}\\
=\cL_\vartheta\Bigl(\sS\mapsto
\diagonalPeriodForm{\vartheta}{i}{\sh}\bigl(\gridPathMap{\sS}{\mathrm{false}}{i}{i}{\mathrm{i}^{\mathrm{gr}}_D(\vartheta)F},\ \gridPathMap{\sS}{\mathrm{true}}{i}{i}{F'}\bigr)\Bigr),
\end{multline*}
where $\mathrm{i}^{\mathrm{gr}}_D(\vartheta)=\mathrm{i}^{\mathrm{gr}}_{D,m}(\vartheta)$ and, for a natural number $n$,
\[\mathrm{i}^{\mathrm{gr}}_{D,n}(\vartheta)=\#\GL_{\GridIndex{D}{\mathrm{false}}{0}{n}}\bigl(\cO_E/\gridModulus{\vartheta}\bigr)\]
is the index of the level of the node $(0,n)$ of the grid false; this level is the trivial subgroup of the
invertible matrices over the residue ring of $\gridModulus{\vartheta}$. The modulus $\gridModulus{\vartheta}$ has
the exponent $b_1(D)$ at the chosen places, so $\mathrm{i}^{\mathrm{gr}}_D(\vartheta)$ is in general not the number
$\mathrm{i}_D(\vartheta)=\mathrm{ord}_{D,\mathrm{false}}(\vartheta,m)$ of Theorem~3, which belongs to the modulus
$\fm_D(\vartheta,1)$; the two are related by Lemma~\ref{l:order}.
\end{enumerate}
\end{definition}
\leanline{\lean{verticalEdge}\sep \lean{rowEdge}\sep \lean{legEdges}\sep \lean{gridPath}\sep \lean{path}\sep
\lean{legDescents}\sep \lean{gridPathMap}\sep \lean{diagonalPeriodForm}\sep \lean{gridPeriod}}

\begin{definition}[routes]\label{d:route}
Let $P$ be parameters, $D$ admissible data, $i\le m$ and $i'\le m$ natural numbers and $e$ a natural number.
$\HasRoute{D}{i}{i'}{e}$ holds if there is a function $\mathrm{d}$ from the $m$-fold twists to the non-zero
complex numbers with $\Cost{D}{e}{\mathrm{d}}$ such that for every $n$ and every family $\cL$ of data with
$\IsTidy{n}{\stagePath{i}}{\cL_\vartheta}$ for all $\vartheta$ there is a family $\cL'$ with
$\IsTidy{n+e}{\stagePath{i'}}{\cL'_\vartheta}$ for all $\vartheta$ and with the following property. For every complex
function $T$ of the triples: if for every real point $\sh$ there is a constant $\mathrm{b}$ with
$\gridPeriod{\cL'}{(\vartheta,F,F')}{i'}{\sh}=\mathrm{b}\cdot T(\vartheta,F,F')$ for all triples, then for every
real point $\sh$ there is a constant $\mathrm{b}$ with
$\gridPeriod{\cL}{(\vartheta,F,F')}{i}{\sh}=\mathrm{b}\cdot T(\vartheta,F,F')/\mathrm{d}(\vartheta)$ for all triples.
In the formal definition the second and third arguments are the proofs of $i\le m$ and $i'\le m$.
\end{definition}
\leanline{\lean{HasRoute}}

\begin{lemma}[binding keeps short sums short]\label{l:bind}
Let $P$, $D$ and an $m$-fold twist $\vartheta$ be given, $n$ and $e$ natural numbers, $\cE$ and $\cE'$ finite
sets of edges, $\mathcal{X}$ a finite set, $\mathrm{k}$ and $\mathrm{t}$ weights and tuples as in
Definition~\ref{d:tuples}\,\textup{(f)}, and $\cL_0$ a datum. Suppose that
$\#\mathcal{X}\le\tupleBound{\vartheta}^{e}$, and that for all $\sS$ and $\mathsf{x}$ with
$\mathrm{k}(\sS,\mathsf{x})\ne0$ and $\LeftInvOnMuEdges{\sS}{\cE}$ the number $\mathrm{k}(\sS,\mathsf{x})$ is a
root of unity and $\LeftInvOnMuEdges{\mathrm{t}(\sS,\mathsf{x})}{\cE'}$ holds. If $\IsShortSum{n}{\cE}{\cL_0}$ holds,
then $\IsShortSum{n+e}{\cE'}{\bind{\cL_0}{\mathrm{k}}{\mathrm{t}}}$ holds.
\end{lemma}
\begin{proof}
Substitute the short sum of $\cL_0$: the new datum is the sum over the pairs of a term of $\cL_0$ and an element
of $\mathcal{X}$, with coefficients the products, which are $0$ or roots of unity, and their number is at most
$\tupleBound{\vartheta}^{n}\cdot\tupleBound{\vartheta}^{e}$.
\end{proof}
\leanline{\lean{binda}}

\sect{Crossing the grid}{s:cross}
\begin{proposition}[crossing a square]\label{p:move}
Let $P$, $D$, an $m$-fold twist $\vartheta$, a flag $o$, natural numbers $i$, $n'$, $n$, a finite set $\cE$ of
edges and a datum $\cL_0$ be given. Let $\pp$, $\mathsf{v}$, $\mathsf{Q}$, $\mathsf{U}$, $\mathsf{R}$ and
$\mathsf{r}$ be nodes, $\mathcal{I}_0$ a finite index set and $\mathcal{J}$ a set; let $\mathsf{e}_1$ be a
bijection from $\GridIndex{D}{o}{i}{n'}$ to the index set of $\mathsf{U}$, $\mathsf{e}_2$ a bijection from the
index set of $\pp$ plus $\mathcal{I}_0$ to that of $\mathsf{Q}$, and $\mathsf{e}_3$ a bijection from the index
set of $\mathsf{r}$ plus $\mathcal{I}_0$ to that of $\mathsf{R}$; let $\mathrm{e}_{\sR}$ be an additive character
of $\sR_\vartheta$; and let $\chi$ be a multiplicative character of $\sR_\vartheta$, $\mathsf{e}$ an identification,
$\rC$ a complex number and $\cG_{\mathsf{j}}$, $\cG'_{\mathsf{j}}$ ($\mathsf{j}\in\mathcal{J}$) functions with
$\SqLaw{\chi}{\mathsf{e}}{\rC}{\cG_{\mathsf{j}}}{\cG'_{\mathsf{j}}}$ for all $\mathsf{j}$. Assume:
\begin{enumerate}[label=\textup{(\roman*)},leftmargin=2.3em,itemsep=1pt]
\item for every admissible family $\sg$, the value $\sg(\mathsf{R})$ is the block extension by $1$ of
$\sg(\mathsf{r})$ through $\mathsf{e}_3$, and for every invertible $\sk$ the family obtained from $\sg$ by putting
$\sk$ at $\pp$ and its block extension through $\mathsf{e}_2$ at $\mathsf{Q}$ is admissible;
\item $\mathsf{U}\ne\mathsf{v}$; either $(\mathsf{Q},\mathsf{U})$ carries the twist and $(\mathsf{U},\mathsf{R})$
does not, or $(\pp,\mathsf{v})$ does and $(\mathsf{v},\mathsf{r})$ does not; the nodes $\mathsf{U}$, $\mathsf{R}$,
$\mathsf{v}$, $\mathsf{r}$ have indices at most $m$ and are different from $\pp$ and $\mathsf{Q}$; and every edge
of $\cE$ has its source outside $\{\pp,\mathsf{Q}\}$ and its target outside
$\{\pp,\mathsf{Q},\mathsf{R},\mathsf{r}\}$;
\item $\IsTidy{n}{\cE\cup\{(\pp,\mathsf{v}),(\mathsf{Q},\mathsf{U})\}}{\cL_0}$ holds;
\item $\chi=\bigl(\edgeChar{\vartheta}{(\pp,\mathsf{v})}\cdot\edgeChar{\vartheta}{(\mathsf{Q},\mathsf{U})}\bigr)^{-1}$,
$\edgeChar{\vartheta}{(\mathsf{U},\mathsf{R})}=\edgeChar{\vartheta}{(\pp,\mathsf{v})}$ and
$\edgeChar{\vartheta}{(\mathsf{v},\mathsf{r})}=\edgeChar{\vartheta}{(\mathsf{Q},\mathsf{U})}$.
\end{enumerate}
Then there is a datum $\cL_1$ with
$\IsTidy{n+2}{\cE\cup\{(\mathsf{U},\mathsf{R}),(\mathsf{v},\mathsf{r})\}}{\cL_1}$ such that for every function
$\mathsf{j}(\sS,\sk)$ with values in $\mathcal{J}$ that depends only on the values of $\sS$ on $\cE$:
\[\cL_0\Bigl(\sS\mapsto\sum_{\sk}\mathrm{e}_{\sR}(\tr(\sk\,\sx_2))\,\cG'_{\mathsf{j}(\sS,\sk)}\tbinom{\sS(\pp,\mathsf{v})}{\sx_1}\Bigr)
=\rC\cdot\cL_1\Bigl(\sS\mapsto\cG_{\mathsf{j}(\sS,\sx'_2)}\bigl(\sS(\mathsf{v},\mathsf{r})\mid\sx'_1\bigr)\Bigr),\]
where $\sx_1$, $\sx_2$ are the column blocks of $\sS(\mathsf{Q},\mathsf{U})$ reindexed by $\mathsf{e}_1$,
$\mathsf{e}_2$, and $\sx'_1$, $\sx'_2$ the row blocks of $\sS(\mathsf{U},\mathsf{R})$ reindexed by $\mathsf{e}_3$,
$\mathsf{e}_1$.
\end{proposition}
\begin{proof}
Let $\cL_1$ be the datum obtained from $\cL_0$ by binding (Definition~\ref{d:tuples}\,(f)) with the graph kernel
of $\chi$ and $\mathsf{e}$ as weight, corrected by the function of orbits of the second note (Input~1) and by the
Fourier factor $\mathrm{e}_{\sR}$: a tuple is replaced by the tuples in which the matrices on the two old edges
are exchanged for matrices on the two new edges. The kernel is a root of unity or zero, and its non-vanishing
transfers left inverses between the blocks (Input~1), so Lemma~\ref{l:bind} gives a short sum of order $n+2$, the
two new matrices being counted by the tuple bound. The equivariance under admissible families follows from the
equivariance of the graph kernel, from (i) and (iv), and from the averaging property of the correcting function.
The identity is the reproducing property of the graph kernel for a square law, $\cG'=\rC\sum\,(\text{kernel})\,\cG$
(Input~1), applied term by term in the short sum of $\cL_0$; the left inverses that it requires are those of
tidiness, by (ii).
\end{proof}
\leanline{\lean{existssquaremove}}

\begin{lemma}[the square law of a square of the grid]\label{l:sqgrid}
Let $P$ be parameters, $D$ admissible data, $o$ a flag and $i$, $n$ natural numbers with $i<m$ and $n<m$. There
are a complex number $\varkappa$, an identification $\mathsf{e}$ and a function $\rc$ from the $m$-fold twists to
the non-zero complex numbers with $\Cost{D}{1}{\rc}$ such that for every $m$-fold twist $\vartheta$, every form
$y$ at the node $(i,n+1)$ of the grid $o$ and every linear functional $\mathsf{l}$ on the forms at $(i+1,n)$:
the functions
$\sx\mapsto\mathsf{l}\bigl(\dn{}{\sx_2}\up{}{\sx_1}y\bigr)$ (lift, then descent, for the column blocks of
$\sx$) and $\mathsf{y}\mapsto\mathsf{l}\bigl(\up{}{\mathsf{y}_2}\dn{}{\mathsf{y}_1}y\bigr)$ (descent, then lift,
for the row blocks of $\mathsf{y}$) satisfy the square law for the character
$\bigl(\edgeChar{\vartheta}{\verticalEdge{o}{i}{n}}\cdot\edgeChar{\vartheta}{\rowEdge{o}{i}{n}}\bigr)^{-1}$, the
identification $\mathsf{e}$ and the constant $\varkappa/\rc(\vartheta)$.
\end{lemma}
\begin{proof}
Apply the theorem on one square (Input~1) to the four spaces at the corners $(i,n)$, $(i+1,n)$, $(i,n+1)$,
$(i+1,n+1)$, with the family of all $m$-fold twists; the edge families are provided by Lemma~\ref{l:corr}, the
characters are tame by Lemma~\ref{l:wts}\,(e), and the numerical hypotheses hold by the small facts. Its constant
has squared absolute value a fixed number times a product of $L$-values of a pair ramified at the chosen places,
so it costs one square (Input~1).
\end{proof}
\leanline{\lean{existssquareLawgrid}}

\begin{proposition}[a descent moves past the lifts of a path]\label{p:past}
Let $P$ be parameters, $D$ admissible data, $o$ a flag and $i$, $n$ natural numbers with $n<m$ and $i\le n$.
There are a complex number $\varkappa$ and a function $\rc$ from the $m$-fold twists to the non-zero complex
numbers with $\Cost{D}{i}{\rc}$ such that for every $m$-fold twist $\vartheta$, every $n_0$, every finite set
$\cE$ of edges and every datum $\cL_0$: if every edge of $\cE$ avoids the grid $o$ or is $\rowEdge{o}{i'}{n}$
with $i'\ge i$, and $\IsTidy{n_0}{\cE\cup\gridPath{o}{i}{n}}{\cL_0}$ holds, then there is a datum $\cL_1$ with
$\IsTidy{n_0+2i}{\{\verticalEdge{o}{i}{n}\}\cup\gridPath{o}{i}{n+1}\cup\cE}{\cL_1}$ and
\[\cL_0\bigl(\sS\mapsto\mathsf{l}_{\sS}(\gridPathMap{\sS}{o}{i}{n}{F})\bigr)
=\frac{\varkappa}{\rc(\vartheta)}\;\cL_1\Bigl(\sS\mapsto\mathsf{l}_{\sS}\bigl(\dn{}{\sS(\verticalEdge{o}{i}{n})}\,\gridPathMap{\sS}{o}{i}{n+1}{F}\bigr)\Bigr)\]
for every form $F$ at $(0,m)$ and every family of linear functionals $\mathsf{l}_{\sS}$ that depends only on the
values of $\sS$ on $\cE$.
\end{proposition}
\begin{proof}
Induction on $i$. For $i=0$ the two sides agree by the definition of the maps along a path. For $i+1$, apply the
induction hypothesis, and then cross the square with corners $(i,n)$, $(i+1,n)$, $(i,n+1)$, $(i+1,n+1)$ by
Proposition~\ref{p:move}, with the square law of Lemma~\ref{l:sqgrid}; each crossing raises the order by $2$ and
multiplies the constant by one that costs one square, and costs add in products (Input~1).
\end{proof}
\leanline{\lean{descentpastgridPathMap}}

\begin{proof}[Proof of Theorem 2$'$ and Theorem 2]
\emph{One step.} Let $\cL$ be tidy of order $n$ on $\stagePath{i}$. In each of the two grids move the last descent of
the path, the one from $(0,i+1)$ to $(0,i)$, past the $i$ lifts of the path by Proposition~\ref{p:past}: each of
the two moves has a constant that costs $i$ squares and raises the order by $2i$. The data are now tidy on the path
of stage $i+1$ in its first $i$ lifts, followed by a descent from $(i,i+1)$ to $(i,i)$ in each grid. Cross the
square at the diagonal by Proposition~\ref{p:move}, with the Fourier form of the square law (Input~1) for the
flipped kernel of the line that the grid true has in addition: its constant costs one square and the order rises
by $2$. The periods are compared by the seesaw identity for the pair formed by the diagonal space and the
subdiagonal one (Input~2), which transfers a proportionality of periods from the new path to the old; the flip
contributes a power of the number of residues, which costs no square. So the product of the constants costs
$2i+1\le4i+2$ squares and the order has risen by $4i+2$; a route records one bound for both numbers. \emph{The climb.} Induction on $i$: the route of cost
$0$ from stage $0$ to itself is the identity, and the routes compose, the functions $\mathrm{d}$ multiplying and
the costs adding, with $2i^2+4i+2=2(i+1)^2$.
\end{proof}

\sect{The start: stage zero}{s:start}
At stage $0$ the path consists of the $m$ descents of each grid from $(0,m)$ to $(0,0)$. Write
$\dn{\vartheta}{o,n}$ for the descent of unit 141 from the forms at $(0,m)$ to the forms at $(0,n)$ of the grid
$o$ for the twist $\vartheta$ (untwisted for $\vartheta=1$), and $\mathcal{A}^{\pi_{\mathrm{gr}}}_n(D)$ for the forms of
principal level $\pi_{\mathrm{gr}}$ at the node $(0,n)$ of the grid true.

\begin{lemma}[orders and indices]\label{l:order}
Let $P$ be parameters, $D$ admissible data and $n$ a natural number. There is a complex number $\rC$ such that for
every $m$-fold twist $\vartheta$:
$\mathrm{ord}_{D,\mathrm{false}}(\vartheta,n)=\rC\cdot\mathrm{i}^{\mathrm{gr}}_{D,n}(\vartheta)$, where
$\mathrm{i}^{\mathrm{gr}}_{D,n}(\vartheta)$ is the index of the level of the node $(0,n)$ of the grid false for
$\vartheta$ (Definition~\ref{d:paths}\,(e)).
\end{lemma}
\begin{proof}
Both numbers count invertible matrices. The index of the level is a mean over the residue ring of
$\gridModulus{\vartheta}$, and the reduction $\reduceToTriangle{\vartheta}{\mathrm{false}}$, which is surjective,
multiplies the counts by a fixed power of the norms of the chosen places.
\end{proof}
\leanline{\lean{triangleOrdereqmullevelIndex}}

\begin{lemma}[translations]\label{l:transl}
Let $P$ be parameters, $D$ admissible data and $\vartheta$ an $m$-fold twist.
\begin{enumerate}[label=\textup{(\alph*)},leftmargin=2.3em,itemsep=1pt]
\item Let $i\le m$ and $\sh$ a real point. For every finite point $\su$ of $\U(\sV_{i,i}(D))$, every form $y$ at
the node $(i,i)$ of the grid false and every form $y'$ at the node $(i,i)$ of the grid true: the form
$\diagonalPeriodForm{\vartheta}{i}{\sh}$ takes the same value at $(y,y')$ and at the translates of $y$ by $\su$
and of $y'$ by $\su\oplus1$.
\item Let $\sg$ be a family of invertible matrices at the nodes with value $1$ at the nodes $(0,m)$ of both
grids. Let $o$ be a flag, $n$ a natural number, $\sS$ a tuple, $F$ a form at $(0,m)$ of the grid $o$, and $\sk$
an element of the maximal compact subgroup whose reduction modulo $\gridModulus{\vartheta}$ is the value of $\sg$
at $(o,0,n)$, with $\sk=1$ if $n=m$. Then the translate by $\sk$ of
$\legDescents{\tupleAction{\sg}{\sS}{\cdot}}{o}{n}{F}$ is
$\edgeCharProd{\sg}{\legEdges{o}{n}}\cdot\legDescents{\sS}{o}{n}{F}$.
\item Let $\sg$ be as in \textup{(b)} and suppose moreover that its value at $(\mathrm{true},0,0)$ is the block
extension through $\gridDiagEquiv{D}{0}$ of its value at $(\mathrm{false},0,0)$. For a real point $\sh$, forms
$F$, $F'$ at the nodes $(0,m)$ of the grids false and true and a tuple $\sS$:
\begin{multline*}
\diagonalPeriodForm{\vartheta}{0}{\sh}\bigl(\legDescents{\tupleAction{\sg}{\sS}{\cdot}}{\mathrm{false}}{0}{F},\legDescents{\tupleAction{\sg}{\sS}{\cdot}}{\mathrm{true}}{0}{F'}\bigr)\\
=\edgeCharProd{\sg}{\stagePath{0}}\cdot\diagonalPeriodForm{\vartheta}{0}{\sh}\bigl(\legDescents{\sS}{\mathrm{false}}{0}{F},\legDescents{\sS}{\mathrm{true}}{0}{F'}\bigr).
\end{multline*}
\end{enumerate}
\end{lemma}
\begin{proof}
(a) Substitute $\sx\mapsto\sx\su$ in the integral; the quotient measure is invariant, the point in the larger
group moves by $\su\oplus1$, and the forms are invariant under rational points. (b) Induction along the descents:
at each step the descent with the moved matrix is the translate of the descent with the old one, by the
transformation law of the datum modulo the modulus, and the character of the edge is collected. (c) Translate
both descents by a lift of the value of $\sg$ at $(0,0)$ and its block extension: (a) leaves the period unchanged
and (b) turns each moved descent into the old one times the product of the edge characters over its edges.
\end{proof}
\leanline{\lean{diagonalPeriodFormformRep}\sep \lean{formReplegDescentstupleAction}\sep
\lean{diagonalPeriodFormtupleAction}}

\begin{proposition}[the descents of unit 141 as tidy data]\label{p:tidy0}
Let $P$ be parameters and $D$ admissible data.
\begin{enumerate}[label=\textup{(\alph*)},leftmargin=2.3em,itemsep=1pt]
\item Let $\vartheta$ be an $m$-fold twist and $n\le m$. There is a datum $\cL_1\in\GridDatum{\vartheta}$ with
$\IsLevelTidy{2(m-n)}{\legEdges{\mathrm{false}}{n}\cup\legEdges{\mathrm{true}}{n}}{\cL_1}$ such that for every
bilinear form $B_1$ and all forms $F$, $F'$ at the nodes $(0,m)$ of the two grids:
\begin{multline*}
B_1\bigl(\dn{\vartheta}{\mathrm{false},n}F,\dn{1}{\mathrm{true},n}F'\bigr)
=\prod_{n\le n'<m}\mathrm{ord}_{D,\mathrm{false}}(\vartheta,n')\,\mathrm{ord}_{D,\mathrm{true}}(\vartheta,n')\\
\cdot\ \cL_1\bigl(\sS\mapsto B_1(\legDescents{\sS}{\mathrm{false}}{n}{F},\legDescents{\sS}{\mathrm{true}}{n}{F'})\bigr).
\end{multline*}
\item There are a complex number $\rC$ and a family $\cL$ of data such that for every $m$-fold twist $\vartheta$:
$\IsTidy{4m}{\stagePath{0}}{\cL_\vartheta}$ holds, and for every real point $\sh$ and all forms $F$, $F'$ at the nodes
$(0,m)$ of the two grids
\[\mathrm{ord}_{D,\mathrm{false}}(\vartheta,m)\cdot\diagonalPeriodForm{\vartheta}{0}{\sh}\bigl(\dn{\vartheta}{\mathrm{false},0}F,\dn{1}{\mathrm{true},0}F'\bigr)
=\rC\cdot\gridPeriod{\cL}{(\vartheta,F,F')}{0}{\sh}.\]
\end{enumerate}
\end{proposition}
\begin{proof}
(a) Decreasing induction on $n$ from $m$. A descent with the point datum of the inclusion is a sum of descents
with indicator data over the matrices that reduce to the inclusion; binding the datum with these matrices for
the two new edges keeps level-tidiness, with one more edge and order two more per step (Lemma~\ref{l:bind}).
(b) Average the datum of (a) for $n=0$ over representatives of the cosets of the levels of the nodes, weighted by
their edge characters. This makes it equivariant for all admissible families, hence tidy, of order $4m$. The
identity follows from Lemma~\ref{l:transl}\,(c) and from the ratios of indices of Lemma~\ref{l:order}.
\end{proof}
\leanline{\lean{existsisLevelTidyiotaPluseq}\sep \lean{existsisTidydiagonalPeriodFormeq}}

\begin{lemma}[restriction to the rows]\label{l:restrict}
Let $P$ be parameters, $D$ admissible data and $\sh$ a real point. There is a linear map $\mathsf{r}_0$ from
$\mathcal{A}^{\pi_{\mathrm{gr}}}_0(D)$ to the forms on $\sY_0(D)$ of the level of the chosen places such that for every
$F'\in\mathcal{A}^{\pi_{\mathrm{gr}}}_0(D)$, every $m$-fold twist $\vartheta$ and every form $y$ on $\sY_0(D)$ of principal
level at the level places of $\vartheta$:
$\pet{y}{\mathsf{r}_0(F')}=\diagonalPeriodForm{\vartheta}{0}{\sh}(\bar y,F')$.
\end{lemma}
\begin{proof}
On the finite-dimensional space of forms of the level of the chosen places, $\mathsf{r}_0$ is the adjoint of the
period functional for the trivial twist. For a general $\vartheta$ the period of a form of the level of
$\vartheta$ equals that of its average over the level of the chosen places, because the twisted and untwisted
periods differ by translations that fix $F'$.
\end{proof}
\leanline{\lean{existsrestrictbordered}}

\begin{lemma}[the descent is adjoint to the weighted chain]\label{l:adjchain}
Let $P$ be parameters and $D$ admissible data.
\begin{enumerate}[label=\textup{(\alph*)},leftmargin=2.3em,itemsep=1pt]
\item Let $o$ be a flag, $\vartheta$ an $m$-fold twist, $n\le m$ and $y$ a form at the node $(0,0)$ of the grid
$o$; let $y_n$ be the image of $y$ under the chain of the first $n$ lifts of the triangle of the data, each
weighted by its triangle order. If $n\ne0$, then every $\dn{\vartheta}{o,n-1}F'$ and $y_n$ are fixed by the elements
of the maximal compact subgroup that reduce to $1$ modulo $\triangleModulus{o}{\vartheta}$; and for every form
$F'$ at $(0,m)$ the integral of $\dn{\vartheta}{o,n}F'\cdot y_n$ over the quotient at $(0,n)$ equals the integral
of $\dn{\vartheta}{o,0}F'\cdot y$ over the quotient at $(0,0)$.
\item $\mathcal{A}^{\pi_{\mathrm{gr}}}_0(D)$ is the image of $\mathcal{A}^{\pi_{\mathrm{gr}}}_m(D)$ under the untwisted descent
$\dn{1}{\mathrm{true},0}$.
\end{enumerate}
\end{lemma}
\begin{proof}
(a) Induction on $n$: a descent and a lift through an edge family are adjoint (Input~2), and both transform under
the elements of the level by the character of the grid, which is trivial on those that reduce to $1$. (b) The
image lies in the target by the invariance in (a). A form in the target that is orthogonal to the image has, by
(a), a weighted chain orthogonal to all forms of level $\pi_{\mathrm{gr}}$ at $(0,m)$, among them the chain itself; so the
chain vanishes, and the form vanishes because weighted chains of the triangle are injective (Input~1).
\end{proof}
\leanline{\lean{iotaPlusadjointweightedChain}\sep \lean{gridLevelFormseqmapiotaPlus}}

\begin{lemma}[steep parameters]\label{l:steep}
Let $P$ be parameters, $D$ admissible data and $o$ a flag. There is an ordering of the index set
$\GridIndex{D}{o}{0}{0}$ such that at every infinite place $\tau$ the parameter
$\gridParam{D}{o}{0}{0}{0}{\tau}$ is steep along this ordering, with the bound $w_{D,0}(\tau)$ and the signs of the
node.
\end{lemma}
\begin{proof}
Order the indices by the numbering of the base indices, compute the signs and the doubled parameters along the
ordering, bound the top weight by $2w_{D,0}(\tau)$ and compute the lowest $K$-type by induction on the rank;
steepness is then a list of inequalities.
\end{proof}
\leanline{\lean{existsisSteep}}

\begin{proposition}[a form orthogonal to all restrictions vanishes]\label{p:vanish}
Let $P$ be parameters, $D$ admissible data, $\bn$ a non-empty finite index set, $\sV$ a hermitian space over
$E(D)$ with index set $\bn$, and $\sV\oplus\langle1\rangle$ its sum with a positive line. Let $\mathrm{A}$, $\mathrm{L}$, $\mathrm{K}_0$
be natural numbers and $\mathrm{B}$, $\mathrm{M}$ real numbers; let $l_\tau$ and $l'_\tau$ be Harish-Chandra parameters
for the signs of $\sV$ and of $\sV\oplus\langle1\rangle$ at the infinite places $\tau$. Assume: each $l_\tau$ has
a large gap $4\mathrm{A}$ and each $l'_\tau$ a large gap $2\mathrm{A}$; $|l_\tau(i)|\le \mathrm{L}$ for all entries; $8(|\bn|+1)^2<\mathrm{A}$;
$\mathrm{K}_0<\mathrm{A}$; $2f(D)(|\bn|+1)+2\le \mathrm{A}$; $g(D)\,\mathrm{c}_{\mathrm{cusp}}^2\le\mathrm{B}$ for the cusp constant of the field
and $\bn$; $2000(|\bn|+\mathrm{L}+1)^6\le\mathrm{B}$; $f(D)\le\mathrm{B}$; $\mathrm{B}^{20\mathrm{A}(f(D)+1)}\le \mathrm{M}$;
$\mathrm{M}^{3|\bn|+2}\le|\pi_{\mathrm{gr}}|_\tau$ at every $\tau$; and for every $\tau$ there is an element $\mathsf{w}_\tau$ of the real
unitary group of $\sV\oplus\langle1\rangle$ with $\|\mathsf{w}_\tau\|^2\le\mathrm{B}$ such that the volume of the unit
ball of the real group of $\sV$ is at most $\mathrm{B}^{\mathrm{K}_0}$ times the absolute value of the integral over that
group of the matrix coefficient of the discrete series of $l_\tau$ against the conjugate of that of $l'_\tau$ at
$\mathsf{w}_\tau^{-1}(\,\cdot\oplus1)\mathsf{w}_\tau$. Then there is a real point $\sh$ of $\U(\sV\oplus\langle1\rangle)$
such that every form $y$ on $\sV$ of the type of $l$ and of the level of the chosen places vanishes if it is
orthogonal to $\sx\mapsto F'((\sx\oplus1)\sh)$ for every form $F'$ on $\sV\oplus\langle1\rangle$ of the type of
$l'$ and of principal level $\pi_{\mathrm{gr}}$.
\end{proposition}
\begin{proof}
Take for $\sh$ the point with the components $\mathsf{w}_\tau$. Let $y$ be orthogonal to all restrictions. The Poincaré series of the larger group are forms
$F'$ of the required kind, and the period of $y$ against the restriction of one of them is, by unfolding, an
integral of $y$ against a matrix coefficient near a point (Input~3). The lower bound in the hypothesis for this
matrix coefficient shows that $|y|$ at a point is at most a fixed fraction of its supremum on a neighbourhood,
wherever the first minimum of the lattice is large enough for the level $\pi_{\mathrm{gr}}$ to separate the terms of the
series; near the cusps $y$ is small by the bound for forms in terms of the first minimum (Input~3). The two
bounds, with the inequalities between $\mathrm{A}$, $\mathrm{B}$, $\mathrm{M}$ and $|\pi_{\mathrm{gr}}|_\tau$, show that the supremum of $|y|$ is at
most half of itself; hence $y=0$.
\end{proof}
\leanline{\lean{eqzeroofpeterssonrestricteqzero}}

\begin{corollary}[every row pairs with a restriction]\label{c:rows}
Let $P$ be parameters and $D$ admissible data. There is a real point $\sh$ of the group of the node $(0,0)$ of the
grid true such that every form $y$ of principal level on $\sY_0(D)$ vanishes if it is orthogonal to
$\sx\mapsto F'((\sx\oplus1)\sh)$ for every $F'\in\mathcal{A}^{\pi_{\mathrm{gr}}}_0(D)$.
\end{corollary}
\begin{proof}
Proposition~\ref{p:vanish} for $\sV=\sY_0(D)$, with $\mathrm{A}=r^4$ and $\mathrm{B}$ a power of $g(D)$ with exponent about
$x$; the gaps and sizes of the parameters are small facts of §\ref{s:grid}, and the lower bound for the matrix
coefficient is the Laplace method for the steep parameters of Lemma~\ref{l:steep} (Input~3).
\end{proof}
\leanline{\lean{eqzerooforthogonalgridLevelForms}}

\begin{proof}[Proof of Theorem 3]
Let $\sh$ be the point of Corollary~\ref{c:rows}, $\cL$ and $\rC$ as in Proposition~\ref{p:tidy0}\,(b), and
$\mathsf{r}_0$ the map of Lemma~\ref{l:restrict} for $\sh$. Put $\mathsf{r}=\mathsf{r}_0\circ\dn{1}{\mathrm{true},0}$
on $\mathcal{A}^{\pi_{\mathrm{gr}}}_m(D)$ and $\mathrm{a}=\rC$. \emph{Surjectivity.} By Lemma~\ref{l:adjchain}\,(b) the image of
$\mathsf{r}$ is $\mathsf{r}_0(\mathcal{A}^{\pi_{\mathrm{gr}}}_0(D))$. A form $y$ of principal level orthogonal to it satisfies
$\diagonalPeriodForm{1}{0}{\sh}(\bar y,F')=0$ for all $F'$ by Lemma~\ref{l:restrict}, that is, $y$ is orthogonal
to all restrictions, and $y=0$ by Corollary~\ref{c:rows}. \emph{The identity.} By Lemma~\ref{l:restrict}, applied
to the conjugate of $\dn{\vartheta}{\mathrm{false},0}F$, which has the level of $\vartheta$ by
Lemma~\ref{l:adjchain}\,(a), the Petersson product of that form with $\mathsf{r}(F')$ is
$\diagonalPeriodForm{\vartheta}{0}{\sh}(\dn{\vartheta}{\mathrm{false},0}F,\dn{1}{\mathrm{true},0}F')$, which
Proposition~\ref{p:tidy0}\,(b) expresses by the period of stage $0$. On the other hand, by
Lemma~\ref{l:adjchain}\,(a) for $n=m$, the same Petersson product is the integral of $F$ against the weighted
chain of $\mathsf{r}(F')$, and the weighted chain is the chain of the $m$ lifts up to the orders and indices that
appear on the two sides: the factor $\mathrm{i}_D(\vartheta)$ on the left is $\mathrm{ord}_{D,\mathrm{false}}(\vartheta,m)$,
the order of the invertible matrices of the rank of $\sY_m(D)$ modulo $\fm_D(\vartheta,1)$, while the period
carries $\mathrm{i}^{\mathrm{gr}}_D(\vartheta)$ inside.
\end{proof}

\sect{The end: stage $m$}{s:end}
\begin{definition}[the last row of the grids and the end period]\label{d:end}
Let $P$ be parameters and $D$ admissible data.
\begin{enumerate}[label=\textup{(\alph*)},leftmargin=2.3em,itemsep=1pt]
\item For a flag $o$ and a natural number $i$, $\EndIndex{o}{i}=\GridIndex{D}{o}{i}{m}$ is the index set of the
\kw{end node} $(i,m)$; the spaces at the end nodes are definite and standard (an instance). $\finWeightData{o}$
is a chosen weight for the forms at the end nodes of the grid $o$, with exponents $\finDetExp{o}$, which goes up
from each end node to the next and is the Harish-Chandra datum of minus the parameter of each end node.
\item For an $m$-fold twist $\vartheta_1$, a flag $o$ and a natural number $i$: $\finLevel{\vartheta_1}{o}{i}$ is
the level of the elements whose reduction modulo $\gridModulus{\vartheta_1}$ lies in the level of the node
$(i,m)$; $\finKernel{\vartheta_1}{o}{i}$ is the theta kernel from the end node $(i,m)$ to $(i+1,m)$ for this
weight, the characters of the grid, the jet scale of the data and that level.
\item The \kw{branching matrix} is a chosen integer matrix that pairs the weights at the nodes $(m,m)$ of the two
grids (it exists by the interlacing of their parameters), and the \kw{end pairing}
$\langle\ ,\ \rangle_{\mathrm{end}}$ is its bilinear form. For a point $\sh$ of $\U(\sV_{m,m}(D))$,
$\finBlockPoint{\vartheta_1}{\sh}$ is the point $\sh\oplus\su^{\,m}$ of the group of the node $(m,m)$ of the grid
true, read through $\gridDiagEquiv{D}{m}$.
\item $\finToAlgForm{o}$ is the isomorphism from the forms at the node $(0,m)$ of the grid $o$ to the algebraic
forms of the weight of \textup{(a)}. For a pair $y=(y_{\mathrm{false}},y_{\mathrm{true}})$ of forms at the nodes
$(0,m)$ of the two grids, $\FixedByEndLevels{\vartheta_1}{y}$ holds if for both $o$ the algebraic form
$\finToAlgForm{o}(y_o)$ is fixed by $\finLevel{\vartheta_1}{o}{0}$.
\item For a tuple $\sS$ (for a twist $\vartheta$), a flag $o$ and $i$: $\finDatum{\sS}{o}{i}$ is the vacuum scale
to the power of the number of entries, times the datum of the indicator of the matrix that $\sS$ puts on
$\rowEdge{o}{i}{m}$; and for a form $F$, $\finChain{\sS}{o}{i}{F}$ is the chain of the first $i$ lifts of
$\finToAlgForm{o}(F)$ along the end nodes, with the kernels $\finKernel{\vartheta}{o}{\cdot}$ and the data
$\finDatum{\sS}{o}{\cdot}$.
\item For a datum $\cL_0\in\GridDatum{\vartheta}$ and a pair $y$ of forms, the \kw{end period} is
\begin{multline*}
\finPeriod{\cL_0}{y}
=\cL_0\Bigl(\sS\mapsto\sum_{\sh}\bigl\langle\finChain{\sS}{\mathrm{false}}{m}{y_{\mathrm{false}}}(\sh),\\
\finChain{\sS}{\mathrm{true}}{m}{y_{\mathrm{true}}}(\finBlockPoint{\vartheta}{\sh})\bigr\rangle_{\mathrm{end}}\Bigr),
\end{multline*}
the sum being the class integral over $\U(\sV_{m,m}(D))$ at the level of the last end node.
\end{enumerate}
\end{definition}
\leanline{\lean{EndIndex}\sep \lean{endWeightData}\sep \lean{endDetExp}\sep \lean{endLevel}\sep \lean{endKernel}\sep
\lean{endBranching}\sep \lean{endPairing}\sep \lean{endBlockPoint}\sep \lean{endToAlgForm}\sep
\lean{FixedByEndLevels}\sep \lean{endDatum}\sep \lean{endChain}\sep \lean{endPeriod}\sep \lean{isStdgridSpace}}

\begin{lemma}[the end nodes]\label{l:endeq}
Let $P$ be parameters and $D$ admissible data.
\begin{enumerate}[label=\textup{(\alph*)},leftmargin=2.3em,itemsep=1pt]
\item For an $m$-fold twist $\vartheta_1$, a flag $o$ and a natural number $i$: $\finLevel{\vartheta_1}{o}{i}$ is
contained in the principal level of $\pi_{\mathrm{gr}}$.
\item For a family $\cL$ of data and a real point $\sh$ there is a complex number $\mathrm{b}$ such that for every
triple $(\vartheta,F,F')$:
$\gridPeriod{\cL}{(\vartheta,F,F')}{m}{\sh}=\mathrm{b}\cdot\finPeriod{\cL_\vartheta}{(F,F')}$.
\end{enumerate}
\end{lemma}
\begin{proof}
(a) An element of the level reduces into the level of the node, which lies in the kernel of the reduction modulo
the chosen places to the power $b_1(D)$. (b) At stage $m$ the path consists of the $m$ lifts along the end nodes,
and these lifts of a form are the chains of theta kernels of Definition~\ref{d:end}; the diagonal period of two
such chains, an integral over a compact quotient, is a class sum, and it is a constant multiple of their
branching pairing (Input~3).
\end{proof}
\leanline{\lean{endLevelleprincipalLevel}\sep \lean{gridPeriodeqendPeriod}}

\begin{proposition}[the end period as a short sum of pairings]\label{p:endmem}
Let $P$ be parameters, $D$ admissible data, $\vartheta$ an $m$-fold twist, $n$ a natural number, $\cL_0$ a datum
with $\IsTidy{n}{\stagePath{m}}{\cL_0}$ and $y$ a pair of forms with $\FixedByEndLevels{\vartheta}{y}$. Let
$\mathcal{S}$ be a property of quadruples (flag, index, point, vector), $\mathsf{F}$ a normed field,
$\sigma_1\colon\CC\to\mathsf{F}$ a ring homomorphism, $\mathrm{M}_1$ a complex number and $\mathrm{R}_1$ a real number. Assume:
$\mathcal{S}$ is preserved by the kernels, that is, for every tuple $\sx$, flag $o$, $i<m$, points $\sg$, $\sh$ and
vector $\mathsf{v}$ with $\mathcal{S}(o,i,\sg,\mathsf{v})$, the property $\mathcal{S}(o,i+1,\sh,\cdot)$ holds for
the image of $\mathsf{v}$ under the jet matrix of $\finKernel{\vartheta}{o}{i}$ at the datum
$\finDatum{\sx}{o}{i}$ moved by $(\sg,\sh)$; $\mathcal{S}$ holds at the start,
$\mathcal{S}(o,0,\sg,\finToAlgForm{o}(y_o)(\sg))$ for all $o$, $\sg$; and for vectors $\mathsf{v}'$, $\mathsf{v}$
with $\mathcal{S}$ at the last end node of the grid false at $\sg$ and of the grid true at
$\finBlockPoint{\vartheta}{\sg}$: $\|\sigma_1(\mathrm{M}_1\langle\mathsf{v}',\mathsf{v}\rangle_{\mathrm{end}})\|\le \mathrm{R}_1$. Then
there is a natural number $\mathrm{n}_1\le\tupleBound{\vartheta}^{n}$ such that $\finPeriod{\cL_0}{y}$ is a sum of $\mathrm{n}_1$
times the product of the class numbers of the end levels complex numbers $\mathrm{a}_1$, each with
$\|\sigma_1(\mathrm{M}_1\mathrm{a}_1)\|\le \mathrm{R}_1$.
\end{proposition}
\begin{proof}
Tidiness makes the summand of the end period, as a function of the point, invariant under the levels of the end
nodes. Expand $\cL_0$ as its short sum, of at most $\tupleBound{\vartheta}^{n}$ terms with coefficients that are
roots of unity; write each class integral as a sum over double cosets and each chain value as an iterated sum over
the classes of the end levels of products of jet matrices applied to the values of the form. Each term is a
pairing of two vectors that have the property $\mathcal{S}$, by its propagation along the kernels.
\end{proof}
\leanline{\lean{endPeriodmem}}

\begin{proof}[Proof of Theorem 4]
\emph{(a)} Take for $\mathrm{d}_{\mathrm{end}}$ the product, over the infinite places, of powers of $g(D)^{b_1(D)}$ and of the
denominators of the branching data, times a power of the denominator of the dual pairing; its size is bounded by
the bounds for the weights of the end nodes (small facts). Apply Proposition~\ref{p:endmem} with the completed
algebraic closure of the $\ell$-adic numbers, $\sigma_1=j$, $\mathrm{M}_1=\mathrm{d}_{\mathrm{end}}$, $\mathrm{R}_1=1$ and the property ``the gauge
of the vector is at most $1$'': it holds at the start by hypothesis, and it is preserved because the data of the
end nodes are standard data and jets of standard data are integral at the primes prime to the modulus (Input~1,
Input~4). So the end period times $\mathrm{d}_{\mathrm{end}}$ is a sum of $j$-integral numbers.

\emph{(b)} Apply Proposition~\ref{p:endmem} with $\sigma_1=\sigma$ and the property ``the norm of the
$\sigma$-conjugate of the vector at the end node $i$ of the grid $o$ is at most
$\exp(i\,x^{99}b^2)\sqrt{B_0(o)}$'', which is preserved by the bounds for the conjugates of jets of admissible
combinations (Input~4). The number of terms is bounded by the tuple bound to the power of the order and by the
class numbers of the end levels, each at most $\exp(x^{99}b^2)$; the sizes of the data give the exponent
$x^{200}b^2$.
\end{proof}

\sect{Proof of Theorem 1}{s:proof}
Let $\cL$, $\sh$, $\mathrm{a}$ and $\mathsf{r}$ be as in Theorem~3, and $\mathrm{d}_{\mathrm{end}}$ as in Theorem~4\,(a). By
Theorem~2 for $i=m$ there are a function $\rc$ with $\Cost{D}{2m^2}{\rc}$ and a family $\cL'$, tidy of order
$4m+2m^2$ on the path of stage $m$, such that proportionalities of the periods of stage $m$ of $\cL'$ descend to
those of stage $0$ of $\cL$, with the factor $1/\rc(\vartheta)$. By Lemma~\ref{l:endeq}\,(b) the periods of stage
$m$ of $\cL'$ are a constant multiple of the end periods. So there is a constant $\mathrm{b}$ with
\[\mathrm{i}_D(\vartheta)\int F\cdot\bigl(\text{chain of lifts of }\mathsf{r}(F')\bigr)
=\frac{\mathrm{b}}{\rc(\vartheta)}\;\finPeriod{\cL'_\vartheta}{(F,F')}\]
for all triples. The constant $\mathrm{b}$ may be zero, and the proof does not divide by it: there is a unit
$\mathrm{t}_0$ with $\mathrm{b}\,\mathrm{t}_0=\mathrm{b}'$, where $\mathrm{b}'$ is $1$ if $\mathrm{b}\ne0$ and $0$ if
$\mathrm{b}=0$.

\emph{From periods to coordinates.} A dual pairing of the weights of $\sY_m(D)$ (Input~3; the parameter of
$\sY_m(D)$ is a Harish-Chandra parameter by the second note) expresses the coordinates of the value of an
algebraic form at a point $\sx$ through its pairings with delta forms: for a vector $\mathsf{v}$, the delta form
at $\sx$ with value $\mathsf{v}$ is an algebraic form of the last level, and the integral of its product with
another form is a class sum with one term. Taking for $F$ the form whose algebraic form is this delta form, the
identity above shows that the coordinates at $\sx$ of
\[\frac{\rc(\vartheta)}{\rc(1)}\cdot\Tlev{\vartheta}\Bigl(\mathrm{t}_0\,\rc(1)\,\frac{\mathrm{d}_{\mathrm{end}}}{w_0}\,\mathsf{r}(F')\Bigr),\]
with $w_0$ the number of elements of the Weyl index set and $\mathrm{t}_0$ the unit above, are integer multiples of
$\mathrm{b}'\,\mathrm{d}_{\mathrm{end}}\cdot\finPeriod{\cL'_\vartheta}{(F,F')}$, up to the hook number of the weight.

\emph{Crude rows.} Call an algebraic form of level $\pi_{\mathrm{gr}}$ on the large definite space \emph{short} if its values
are integral and small in every conjugate. The pair $(F,F')$ of a delta form and a short form is fixed by the end
levels (Lemma~\ref{l:endeq}\,(a) and the level of the delta form). By Theorem~4\,(a) the numbers
$\mathrm{d}_{\mathrm{end}}\cdot\finPeriod{\cL'_\vartheta}{(F,F')}$ are $j$-integral for every prime $\ell$ prime to
$\Np Q(\vartheta)$ and every ring homomorphism $j\colon\CC\to\CC_\ell$; this is property~(1) for the row
$y=\mathrm{t}_0\,\rc(1)\,\mathrm{d}_{\mathrm{end}}\,w_0^{-1}\,\mathsf{r}(F')$. For property~(2), take $\vartheta=1$ and an automorphism
$\sigma$; a point at which the Weyl norm of the conjugate vector is attained and the dual pairing reduce the bound to
the end period of the conjugated delta form, which Theorem~4\,(b) bounds by
$\exp(x^{200}b^2)$ times the square roots of the bounds for the delta form and for the short form; with the
bounds for the hook number, for $\mathrm{d}_{\mathrm{end}}$ and the comparison of the scales of the data, the
squared Weyl norm is at most $\exp(x_1^2b^2)$.

\emph{They span.} The short forms span the forms of level $\pi_{\mathrm{gr}}$ on the large definite space (Input~3), and
$\mathsf{r}$ is surjective and linear. So the rows $y$ above span the forms of principal level on $\sY_0(D)$.
\leanline{the formal proof of \lean{cruderows} uses \lean{existsisTidyintegraleqgridPeriod}, \lean{gridPeriodclimb},
\lean{gridPeriodeqendPeriod}, \lean{existsmulendPeriodmemembValRing} and \lean{endPeriodle} as here}

\sect{Notes on the formal text}{s:notes}
\begin{enumerate}[label=\textup{\arabic*.},leftmargin=1.6em,itemsep=5pt]
\item \emph{Theorem~1 and the overview.} The overview of the handout says that in the proof of Theorem~C ``the
family is completed to a basis by crude rows''. Theorem~1 provides them: rows that span, whose twisted chains of
lifts are integral after the factor $\rc(\vartheta)/\rc(1)$, and whose untwisted lifts are bounded by
$\exp(x_1^2b^2)$ at every conjugate. The cost $2m^2$ of $\rc$ is what the fit families of twists have to pay for
these rows. In the formal statement the forms $y$ are the forms of principal level on $\sY_0(D)$, a space that
contains the rows; property~(2) is stated for the chain of lifts on all such forms. In property~(1) the formal
statement quantifies over all $\ell$-adic places, that is, over all ring homomorphisms $j\colon\CC\to\CC_\ell$,
as Theorem~C of the handout needs; an earlier edition of this note quantified over the isomorphisms of $\CC$ with
an algebraic closure of the $\ell$-adic numbers, which gives only the homomorphisms with image in that closure and
is weaker. The norm in property~(2) and in Theorem~4\,(b) is the Weyl norm of the weight representation.
\item \emph{Routes carry two numbers in one.} In the definition of a route the same number $e$ bounds the cost of
the constant and the growth of the order of tidiness. In the step from stage $i$ the constants cost $2i+1$ squares
and the order grows by $4i+2$; costs are upper bounds (the accounting is that of the second reading).
\item \emph{The constant $\mathrm{b}$ in the proof of Theorem~1.} The formal proof replaces the constant
$\mathrm{b}$ by $\mathrm{b}/\mathrm{b}$, which is $1$ or $0$, and chooses a unit accordingly; the case
$\mathrm{b}=0$ is not excluded, and cancelling $\mathrm{b}$ would lose it (second reading).
\item \emph{What Theorem~1 does not say.} It is a spanning statement: the quantifiers over twists, primes,
automorphisms and points apply to each generator $y$, and the statement names no basis and no extension of a
given independent family; the completion to a basis is made where the theorem is used. Theorem~4\,(b) is proved
for the trivial twist only, and the order $4m+2m^2$ is the order $4m$ of the data of Theorem~3 plus the bound
$2m^2$ of the climb.
\item \emph{Two census entries that are not arguments.} For the maps along a path
(Definition~\ref{d:paths}\,(b), (c)) the tool that reads the statements takes the arrow of the type ``linear
map'' for a premise. In the Lean text these definitions take the tuple, the flag and the indices and return a
linear map; here the form to which the map is applied fills the last slot of the notation.
\item \emph{Proofs as arguments.} The formal definitions of the diagonal period, of the period of a stage and of
a route take proofs of inequalities $i\le m$ as arguments; in the notation their place is taken by the index.
\item \emph{Arguments that only Lean can settle.} For the declarations that the dossier marks so, the list of
arguments was derived by Lean's rules by hand; no compiler was run. Among the statements above they are
Theorem~4\,(a), Proposition~\ref{p:move}, Proposition~\ref{p:endmem}, Proposition~\ref{p:tidy0}\,(b),
Lemma~\ref{l:order}, Lemma~\ref{l:restrict}, Lemma~\ref{l:bind}, Lemma~\ref{l:corr}\,(b),
Lemma~\ref{l:wts}\,(a), (b), (e), the definitions of tidy and level-tidy data, of the tuple bound, of the maps
along a path, of the period of a stage and of the end objects, and several small facts. The doubt concerns
implicit arguments: whether a declaration takes the data, the twist or an index set explicitly or implicitly,
and whether an instance for a prime is a separate argument.
\item \emph{Small facts.} The one-step results and computations listed in §\ref{s:grid} have no statement of
their own; the build of this note records them as absorbed, together with the definition of the harmonic of the
grid, which has more arguments than a notation can carry and is described in words. Their formal statements are
in the concordance.
\item \emph{Two indices.} The number $\mathrm{i}_D(\vartheta)$ of Theorem~3 and the number
$\mathrm{i}^{\mathrm{gr}}_D(\vartheta)$ inside the period of a stage are different: the first is the order of the
invertible matrices of the rank of $\sY_m(D)$ over the residue ring of $\fm_D(\vartheta,1)$, the second the same
order over the residue ring of $\gridModulus{\vartheta}$, whose exponent at the chosen places is $b_1(D)$. An
earlier edition wrote both with one symbol. They differ by a factor that does not depend on $\vartheta$
(Lemma~\ref{l:order}).
\item \emph{What is taken for granted.} Everything listed as Inputs: the second note (one square, the triangle,
costs), the seesaw identity for periods in the weaker form of the background library (an inference from Kudla's
identity, not a statement of his paper), the dual and branching pairings, the spanning by short forms, the bounds
of reduction theory, the Laplace method, the Poincaré series as kernels of periods (whose relation to the
literature the second reader found unresolved), and the integrality and the bounds for jets. The parts (c), (d), (e) of Lemma~\ref{l:wts} and Definition~\ref{d:wts} are used by
other units and not in this note.
\end{enumerate}

\subhead{How this note was written.} From the dossier of the unit (its third part). The program check of the
whole dossier covered \UnitSteps{} leaf steps of \UnitProofs{} proofs and found every statement at its lines. I
read all statements of this part in the dossier, the summaries of all its structured proofs, and step by step the
proof of Theorem~1. In the Lean text I read the statements of Theorems 1 to~4, the definitions of routes, of the
maps along a path, of the diagonal period and of the period of a stage. The other statements are written from
the dossier's readings and were not compared with the Lean text by me. The proofs are outlines that name every
input, shorter than the formal ones; those of Theorem~2$'$, Theorem~4, Proposition~\ref{p:move},
Proposition~\ref{p:endmem} and Proposition~\ref{p:vanish} follow the summaries of structured proofs whose long
steps the readers mark as coarse; the accounting of orders and costs in the proof of Theorem~2$'$ is that of the
second reader. This edition is a revision after a referee's report on the three notes of the unit: in
Theorem~1\,(1) the quantifier is now the formal one (all ring homomorphisms $\CC\to\CC_\ell$; I read the
statement and the definition of an $\ell$-adic place again in the Lean text), the two indices of Theorem~3 and of
Definition~\ref{d:paths}\,(e) have two symbols (read in the Lean text: the statements of Theorem~3, of the period
of a stage and of Lemma~\ref{l:order}, and the definitions of the levels of the nodes and of the triangle
orders), and the relations of the Inputs to the literature follow the second reader's verdicts.
Definition~\ref{d:grid}\,(b), (c) describe the signs and
the parameters of the grids in words; their formulas are in the dossier and in the concordance. The figure is
mine. Nothing was compiled.

\head{Appendix A. Concordance with the Lean artifact}
{\fontsize{8.6}{11.6}\selectfont\renewcommand{\leanfont}{\ttfamily\fontsize{7.0}{9.6}\selectfont\color{leancol}}
\setlength{\LTleft}{0pt}\setlength{\LTright}{\fill}\setlength{\tabcolsep}{0pt}
\begin{longtable}{@{}p{30mm}p{56mm}p{43mm}r@{}}
\textbf{In this note} & \textbf{Lean declaration} & \textbf{File} & \textbf{Lines}\\\midrule\endhead
\ConcordanceRows
\end{longtable}}

\subhead{Three statements as they are written in the artifact.}
\quoted{cruderows}
\quoted{HasRoute}
\quoted{gridPeriod}
Legend: {\leanfont D} is the data $D$; {\leanfont c} is $\rc$; {\leanfont CostsSquares} is the cost in squares;
{\leanfont iotaAtLevel} is $\Tlev{\vartheta}$ and {\leanfont D.ι .one} is $\Tall{1}$; {\leanfont ϑ.levelNorm} is
$\Np Q(\vartheta)$; {\leanfont Λ}, {\leanfont Φ} are families of data $\cL$, $\cL'$; {\leanfont IsTidy} and
{\leanfont D.path} are tidiness and the path of a stage; {\leanfont z} is the flag $o$; {\leanfont S} is a tuple
$\sS$; {\leanfont levelIndex} at the node $(0,m)$ of the grid false is $\mathrm{i}^{\mathrm{gr}}_D(\vartheta)$; the
embedding {\leanfont ι} of the first text is the ring homomorphism $j\colon\CC\to\CC_\ell$ (an $\ell$-adic place).

\head{Notation}
\begingroup\fontsize{8.6}{10.9}\selectfont\raggedright
\setlength{\columnsep}{9mm}\setlength{\parskip}{0pt}
\begin{multicols}{2}
\nt{$P$, $D$, $E(D)$, $m$, $r$, $k(D)$, $t(D)$, $f(D)$, $g(D)$, $x$, $x_1$, $b$, $b_1(D)$}{parameters, admissible data and their field, the number of lifts, the ranks of rows and columns, the step of the ranks, the number of infinite places, the second level, the three scales, an exponent of the data}
\nt{$\vartheta$, $\vartheta_1$, $Q(\vartheta)$, $\Np Q(\vartheta)$, $\theta$, $\theta_{\ell,\iota}$}{twists, their set of places and its norm; a fit family of twists, indexed by the pairs $(\ell,\iota)$}
\nt{$\tau$; $\sigma$; $\ell$, $\iota$, $j$, $\CC_\ell$; $\bar\pi_g$, $\pi_{\mathrm{gr}}$}{an infinite place; an automorphism of $\CC$; a prime, an isomorphism of $\CC$ with an algebraic closure of the $\ell$-adic numbers, a ring homomorphism from $\CC$ to the completion $\CC_\ell$ of that closure (an $\ell$-adic place); the conjugate of the integer of norm $g(D)$, the generator of the grid}
\nt{$o$, $o'$, $[o]$, $\mathrm{d}$}{flags (false or true), the number $0$ or $1$ of a flag; a direction or a function of twists}
\nt{$i$, $i'$, $i_1$, $n$, $n'$, $n_0$, $n_1$, $e$}{natural numbers: indices of nodes, stages, offsets, orders of tidiness, costs}
\nt{$\bn$, $|\bn|$; $\fp$; $\A$, $\CC$, $\cO_E$}{a finite index set and its size; a finite place; adeles, complex numbers, integers}
\nt{$\GridIndex{D}{o}{i}{n}$, $\gridNumber{D}{o}$, $\gridBlockEquiv{D}$, $\gridDiagEquiv{D}{i}$}{the index set of a node, the numbering of its base, the listing of its indices, the diagonal bijection (Definition~\ref{d:grid})}
\nt{$\gridSigns{D}{o}{n'}{i}{n}{\mathsf{x}}$, $\gridParam{D}{o}{n'}{i}{n}{\tau}$, $\signsY{D}{n}{\mathsf{x}}$, $\paramY{D}{n}{\tau}{\mathsf{x}}$, $\gridWeight{D}{o'}{i'}{\bn}{\tau}$}{signs, doubled parameters and weights of the grids (Definition~\ref{d:grid})}
\nt{$\sV_{i,n}(D)$, $\sZ_{i,n}(D)$, $\sY_n(D)$, $\sW_i(D)$, $\sV$, $\sV\oplus\langle1\rangle$; $\U(\sV)$, $d\sx$}{the spaces of the two grids, of the chain and of the auxiliary pairs; a hermitian space and its sum with a positive line; the unitary group, integration for an invariant measure}
\nt{$\gridModulus{\vartheta}$, $\triangleModulus{o}{\vartheta}$, $\reduceToTriangle{\vartheta}{o}$, $\sR_\vartheta$; $\fm_D(\vartheta,e)$}{the moduli of the grid and of the triangle, the reduction between their residue rings, the residue ring of the grid; the modulus of the chosen places to the power $e$ times $Q(\vartheta)$}
\nt{$\pow{\vartheta}{n}$, $\MFoldTwist{D}$, $\powTwist{\vartheta}{n}$, $\chainChar{D}{\vartheta}{i}$, $\vartheta^{\mathrm{t}}_{\fp}$}{powers of tame characters, the $m$-fold twists, the characters of a chain, a local tame part (Definition~\ref{d:twists})}
\nt{$\pp$, $\omega_{\pp}$, $e_{\pp}$, $\pairRatio{\pp}{\vartheta}{\pm}$; $\omega$, $\omega^{\mathrm{t}}$, $\mu$}{a pair of the data, its character and point, the ratio of $L$-values; a character of $\U(1)(\A)$, its tame part, weights}
\nt{$L^{\Sigma\cup\bar\Sigma}(s,\omega^{\mathrm{bc}})$; $\omega^{\mathrm{bc}}$; $\Sigma$, $\bar\Sigma$, $\Sigma_D$; $s$}{the Hecke $L$-function of the base change without the Euler factors at $\Sigma$ and at the conjugate places; the base change of $\omega$; a set of finite places, the set of their conjugates, the chosen places; a real number $\ge1$}
\nt{$\nStar{D}$, $\orderW{i}{\mathsf{x}}$, $\conjHolWeight{i}{\tau}$, $\kernelWeight{i}{\tau}$}{a size, an order of a basis, a lowest $K$-type, a kernel weight (Definition~\ref{d:wts})}
\nt{$\fl$, $\fa$}{an $L$-value datum, an auxiliary pair}
\nt{$\sS$, $\sS_{\mathsf{x}}$, $\EdgeTuple{\vartheta}$; $\cL$, $\cL'$, $\cL_0$, $\cL_1$, $\GridDatum{\vartheta}$}{tuples of matrices on the edges; data and families of data}
\nt{$\pp$, $\pp'$, $\sq$, $\mathsf{v}$, $\mathsf{Q}$, $\mathsf{U}$, $\mathsf{R}$, $\mathsf{r}$; $\mathsf{b}$, $\cE$, $\cE'$}{nodes; an edge, finite sets of edges}
\nt{$\IsMuEdge{\pp}{\sq}$, $\edgeChar{\vartheta}{\mathsf{b}}$, $\IsAdmissibleFamily{\sg}$, $\tupleAction{\sg}{\sS}{\mathsf{b}}$, $\edgeCharProd{\sg}{\cE}$, $\tupleBound{\vartheta}$, $\LeftInvOnMuEdges{\sS}{\cE}$}{edges that carry the twist, edge characters, admissible families, their action, products of edge characters, the tuple bound, left inverses (Definition~\ref{d:tuples})}
\nt{$\IsShortSum{n}{\cE}{\cL_0}$, $\bind{\cL_0}{\mathrm{k}}{\mathrm{t}}$, $\IsTidy{n}{\cE}{\cL_0}$, $\IsLevelTidy{n}{\cE}{\cL_0}$; $\mathcal{W}$, $\mathcal{X}$, $\rc_{\mathsf{x}}$, $\mathrm{k}$, $\mathrm{t}$}{short sums, binding, tidy and level-tidy data (Definition~\ref{d:tuples}); the sets, coefficients, weights and tuples in them}
\nt{$\verticalEdge{o}{i}{n}$, $\rowEdge{o}{i}{n}$, $\legEdges{o}{n}$, $\gridPath{o}{i}{n}$, $\stagePath{n}$}{edges and paths (Definition~\ref{d:paths})}
\nt{$\legDescents{\sS}{o}{n}{F}$, $\gridPathMap{\sS}{o}{i}{n}{F}$, $\diagonalPeriodForm{\vartheta_1}{i}{\sh}$, $\gridPeriod{\cL}{(\vartheta,F,F')}{i}{\sh}$}{the maps along a path, the diagonal period, the period of a stage (Definition~\ref{d:paths})}
\nt{$\HasRoute{D}{i}{i'}{e}$; $T$, $\mathrm{b}$, $\mathrm{b}'$, $\mathrm{t}_0$}{routes (Definition~\ref{d:route}); a function of triples, constants of proportionality, a unit}
\nt{$F$, $F'$, $y$, $y'$, $y_n$, $y_o$; $\pet{y}{y'}$; $\mathsf{l}$, $\mathsf{l}_{\sS}$; $G$, $\cG$, $\cG'$, $\cG_{\mathsf{j}}$; $B_0$, $B_1$}{forms; the Petersson product; linear functionals; functions of tuples and of matrices; a bound, a bilinear form}
\nt{$\sg$, $\sh$, $\sk$, $\sx$, $\su$, $\mathsf{x}$, $\mathsf{y}$, $\mathsf{w}_\tau$; $\mathsf{e}$, $\mathsf{e}_1$, $\mathsf{e}_2$, $\mathsf{e}_3$; $\mathsf{j}$, $\mathcal{J}$, $\mathcal{I}_0$}{points of groups, families of matrices, matrices and indices; identifications and bijections of index sets; an index, a set, a finite index set}
\nt{$\SqLaw{\chi}{\mathsf{e}}{\rc}{\cG}{\cG'}$, $\Cost{D}{n}{\rc}$; $\chi$, $\mathrm{e}_{\sR}$; $\rc$, $\rC$, $\varkappa$, $\mathrm{a}$, $\mathrm{a}_1$, $\mathrm{d}_{\mathrm{end}}$, $w_0$}{the square law and costs of the second note; a multiplicative and an additive character of a finite ring; constants}
\nt{$\Tlev{\vartheta}$, $\Tall{1}$, $\operatorname{ch}_{\vartheta}$, $w_{D,i}(\tau)$}{the composite of the $m$ lifts at principal level, the chain for the trivial twist, the chain as functions, the weights of the lifts}
\nt{$\mathrm{i}_D(\vartheta)$; $\mathrm{i}^{\mathrm{gr}}_D(\vartheta)$, $\mathrm{i}^{\mathrm{gr}}_{D,n}(\vartheta)$; $\mathrm{ord}_{D,o}(\vartheta,n)$}{the index of the last level of the chain, $\#\GL(\cO_E/\fm_D(\vartheta,1))$ in the rank of $\sY_m(D)$; the index of the level of the node $(0,m)$, of the node $(0,n)$, of the grid false, $\#\GL(\cO_E/\gridModulus{\vartheta})$ in the rank of that node; the triangle orders (second note)}
\nt{$\xi$, $\mathrm{G}$}{a coordinate vector of a weight, the Gram matrix of the weight representation (squared Weyl norm $\xi^{*}\mathrm{G}\xi$)}
\nt{$\up{}{\sx}$, $\dn{}{\sx}$, $\dn{\vartheta}{o,n}$, $\mathcal{A}^{\pi_{\mathrm{gr}}}_n(D)$, $\mathsf{r}$, $\mathsf{r}_0$}{lifts and descents along an edge with a matrix, the descent from $(0,m)$ to $(0,n)$, forms of principal level $\pi_{\mathrm{gr}}$, the restriction maps}
\nt{$\EndIndex{o}{i}$, $\finWeightData{o}$, $\finDetExp{o}$, $\finLevel{\vartheta_1}{o}{i}$, $\finKernel{\vartheta_1}{o}{i}$, $\finBlockPoint{\vartheta_1}{\sh}$, $\finToAlgForm{o}$, $\FixedByEndLevels{\vartheta_1}{y}$, $\finDatum{\sS}{o}{i}$, $\finChain{\sS}{o}{i}{F}$, $\finPeriod{\cL_0}{y}$, $\langle\ ,\ \rangle_{\mathrm{end}}$}{the objects of the end nodes and the end period (Definition~\ref{d:end})}
\nt{$\mathcal{S}$, $\mathsf{F}$, $\sigma_1$, $\mathrm{M}_1$, $\mathrm{R}_1$, $\mathrm{n}_1$; $\mathrm{A}$, $\mathrm{L}$, $\mathrm{K}_0$, $\mathrm{B}$, $\mathrm{M}$, $l_\tau$, $l'_\tau$, $\mathrm{c}_{\mathrm{cusp}}$}{the property, field, homomorphism and numbers of Proposition~\ref{p:endmem}; the numbers and parameters of Proposition~\ref{p:vanish}}

\end{multicols}
\endgroup

\end{document}
